\documentclass
[
    twoside,                 % The thesis is formatted like a book. That is, odd and even pages are handled differently.
    openany,                 % Chapters start on the next available page. Was `openright`, which forces each
                             % chapter onto an odd page and inserts a blank whenever the previous one ended
                             % on an odd page -- a binding convention that costs several empty pages in a PDF.
                             % Restore `openright` if the thesis is printed and bound double-sided.
    cleardoublepage = empty, % Any clear page that does get inserted is formatted with an empty style.
    fontsize = 12 pt,        % The size of the font.
    american,                % Support for American English.
    captions = tableheading, % Places the correct amount of space when the caption of a table is above the table.
    numbers = noenddot,      % Does not use a period at the end of numbered titles, such as sections or figures.
    footheight = 35 pt,      % Defines the height of the foot. Due to the line, it needs extra height.
%    draft,                   % Only displays boxes of figures. This option is useful if compilation takes a long time.
]
{scrbook}


% This file contains all sorts of commands that are used in order to specify certain options for the document.

\newif\ifprintVersion   % Defines a binary variable that signals whether the document is prepared for physical or digital print.
\newif\ifprofessionalPrint % Defines a binary variable that signals whether the print will be done by a professional printing service that requests extra margin for page cutting and is not bound to paper formats like A4.
\newif\iffancyTheorems  % Defines a binary variable that signals whether theorems are formatted in the classical style or in a new format that better suits the overall flavor of this thesis.
\newif\ifboldNumberSets % Defines a binary variable that signals whether the variables for number sets (like N or R) should be in bold. If not, they are in blackboard bold instead.
\newif\ifbachelorThesis % Defines a binary variable that signals whether this thesis is a bachelor thesis (true) or a master thesis (false).

% Set all variables to their default values.
\printVersionfalse
\professionalPrintfalse
\fancyTheoremstrue
\boldNumberSetstrue
\bachelorThesistrue

%%%%%%%%%%%%%%%%%%%%%%%%
% The following commands define certain strings that provide important information for the document.

% The title of the thesis.
\newcommand*{\printTitle}{}
\newcommand*{\printGermanTitle}{}
\newcommand*{\myTitle}[2]{\renewcommand*{\printTitle}{#1}\renewcommand*{\printGermanTitle}{#2}}
\newcommand*{\printTitleBold}{\textbf{\printTitle}}

% The author’s name.
\newcommand*{\printAuthor}{}
\newcommand*{\myName}[1]{\renewcommand*{\printAuthor}{#1}}

% The name of the author’s program.
\newcommand*{\printProgram}{}
\newcommand*{\myProgram}[1]{\renewcommand*{\printProgram}{#1}}

% The date when the thesis was handed in.
\newcommand*{\printDateReceived}{}
\newcommand*{\dateOfHandingIn}[1]{\renewcommand*{\printDateReceived}{#1}}

% A short description of the topic of the thesis. This string will be used for the PDF metadata.
\newcommand*{\printSubject}{}
\newcommand*{\mySubject}[1]{\renewcommand*{\printSubject}{#1}}

% A short description of the topic of the thesis. This string will be used for the PDF metadata.
\newcommand*{\printKeywords}{}
\newcommand*{\myKeywords}[1]{\renewcommand*{\printKeywords}{#1}}

% The name of the author’s supervisor.
\newcommand*{\printNameOfSupervisor}{}
\newcommand*{\nameOfMySupervisor}[1]{\renewcommand*{\printNameOfSupervisor}{#1}}

% The name of the author’s second supervisor.
\newcommand*{\printNameOfSecondSupervisor}{}
\newcommand*{\nameOfMySecondSupervisor}[1]{\renewcommand*{\printNameOfSecondSupervisor}{#1}}

% The list with the name of the additional examiners.
\newcommand*{\printAdditionalExaminers}{}
\newcommand*{\additionalExaminers}[1]{\renewcommand*{\printAdditionalExaminers}{#1}}

% Defines the extra length added to each side for the print version.
\newlength{\extraborderlength}
\newcommand*{\extraBorder}[1]{\setlength{\extraborderlength}{#1}}

% Defines the length of the binding correction. (The class ›scrbook‹ has a binding correction but it does not work due to all the other packages that are loaded.)
\newlength{\mybindingcorrection}
\newcommand*{\bindingCorrection}[1]{\setlength{\mybindingcorrection}{#1}} % Contains commands that are used for certain information that is printed.


%%%%%%%%%%%%%%%%%%%%%%%%%%%%%%%%%%%%%%
%% Please adjust your options here. %%
%%%%%%%%%%%%%%%%%%%%%%%%%%%%%%%%%%%%%%

    % This section contains commands with important data for your thesis. Please adjust them in order for the document to be printed correctly.

    % Defines the length of the amount that a printed page will be cut from EACH side (including the inner side). This option only takes effect with \printVersiontrue and \professionalPrinttrue.
    \extraBorder{3 mm}

    % Shifts the inner margin outward by the amount specified. When the book is bound, part of the page will not be seen anymore. This option compensates for this loss. It only takes effect with \printVersiontrue.
    \bindingCorrection{6 mm}

    % Use the following command if this is a master thesis.
    \bachelorThesisfalse

%    \printVersiontrue      % Use this value if you want to prepare your thesis for physical printing. In this case, links will not be colored. Without \professionalPrinttrue, the content will be moved outward by the binding correction, increasing the inner margin and decreasing the outer margin.
%    \professionalPrinttrue % Use this value if you want to have extra border for cutting and are not bound to paper formats (like A4). This option will increase the page size by the extra border on every side plus the binding correction once for the width. It only takes effect in combination with \printVersiontrue.
%    \fancyTheoremsfalse  % Use this value if you want to use the classical theorem style, where the text is italic. Further, with this style, the QED symbol is colorless.
%    \boldNumberSetsfalse % Use this value if you want variables for number sets (like N or R) to appear in blackboard bold rather than bold.

    % The title of the thesis. The first argument is for the English name, the second is for the German name.
    \myTitle{An Empirical Study of Knowledge-Graph-Driven Context Retrieval for LLM-Based Pull-Request Review Comments}{Eine empirische Studie zur wissens\-graph\-gestützten Kontext\-abfrage für LLM-basierte Pull-Request-Review-Kommentare}

    % The author’s name.
    \myName{Abdelrahman Khattab}

    % The author’s program.
    \myProgram{Computer Science}

    % The date when the thesis will be handed in.
    \dateOfHandingIn{8 October 2026}

     % The name and affiliation of the author’s supervisor.
    \nameOfMySupervisor{Prof.\,Dr. Holger Giese}

     % The name and affiliation of the author’s supervisor.
    \nameOfMySecondSupervisor{Prof.\,Dr. Walid Maalej}

    % A list with the names of the additional examiners.
    \additionalExaminers{Christian Medeiros Adriano}

    % A short summary of the thesis. These information will be used for the PDF meta data.
    \mySubject{Knowledge-graph and retrieval-augmented context for LLM-generated code review}

    % Some keywords of the thesis. These information will be used for the PDF meta data. Please use | as a separator and try to avoid commas.
    \myKeywords{code review | large language models | knowledge graphs | retrieval-augmented generation | empirical software engineering}

%%%%%%%%%%%%%%%%%%%%%%%%%%%%%%%%%%%%%%
%% End of options to adjust. %%%%%%%%%
%%%%%%%%%%%%%%%%%%%%%%%%%%%%%%%%%%%%%%


% This file includes all of the code that is used to format the thesis.
% Some packages are included if they are needed. This is done in the respective part and not at the beginning of this file.
%
% This file contains the following parts:
%   • Language an Character Set
%   • Penalties
%   • Indentation
%   • Footnotes
%   • Colors
%   • Size and Position of the Text Body
%   • Position of the Head and the Foot
%   • Margin Position and Width
%   • Header and Footer Format
%   • Caption Format
%   • Part Format
%   • Chapter Format
%   • Table of Contents


%%%%%%%%%%%%%%%%%%%%%%%%%%%%%%%%
%% Language and Character Set %%
%%%%%%%%%%%%%%%%%%%%%%%%%%%%%%%%

\usepackage[utf8]{inputenc} % Allows to input UTF8 characters.
\usepackage[T1]{fontenc}    % Allows to print special characters correctly.
\usepackage
[
    ngerman,         % German is used for the German abstract.
    main = american, % This is the main language of the thesis.
]
{babel}                     % Is responsible for sensible hyphenations.

% The following two commands make it such that the LaTeX compiler of Overleaf produces a PDF from with ligatures and mathematical symbols can be copied correctly.
% Also refer to: https://tex.stackexchange.com/questions/64188/what-are-good-ways-to-make-pdflatex-output-copy-and-pasteable
% Guarded: both primitives are pdfTeX-only, so an unguarded \input breaks any XeTeX or
% LuaTeX build. Overleaf uses pdfLaTeX and is unaffected.
\usepackage{iftex}
\ifPDFTeX
    \input glyphtounicode
    \pdfgentounicode=1
\fi


%%%%%%%%%%%%%%%
%% Penalties %%
%%%%%%%%%%%%%%%

\widowpenalties 2 10000 0


%%%%%%%%%%%%%%%%%
%% Indentation %%
%%%%%%%%%%%%%%%%%

\usepackage{calc} % Makes it easer to do math with TeX measurements.

\newlength{\myparindent}
\newlength{\myparskip}
\setlength{\myparindent}{1 em}
\setlength{\myparskip}{0 em}

\setlength{\parindent}{\myparindent}
\setlength{\parskip}{\myparskip}
\setlength{\parskip}{0 pt plus 1 pt minus 0 pt}


%%%%%%%%%%%%%%%
%% Footnotes %%
%%%%%%%%%%%%%%%

% Remove the footnote rule.
\setfootnoterule{0 cm}

% The footnote number is made bold and not in superscript.
\deffootnote[1.2 em]{1.2 em}{0 em}{\makebox[1.4 em][l]{\textbf{\thefootnotemark}}}

% The footnote number will not be reset after every chapter.
\makeatletter%
    \@removefromreset{footnote}{chapter}%
\makeatother


%%%%%%%%%%%%
%% Colors %%
%%%%%%%%%%%%

\usepackage[dvipsnames]{xcolor} % Allows it to define colors. The option says that common names can be used.

% Dark blue.
\definecolor{stroke1}{HTML}{2574A9} % This color is used as the standard color to highlight things.


% Coloring various different labels.
\colorlet{captionlabel}{black}
\colorlet{footerpagenr}{black}
\colorlet{footerchapter}{stroke1}
\colorlet{footerchaptername}{black}
\colorlet{footersection}{stroke1}
\colorlet{footersectionname}{black}
\colorlet{chapternumber}{stroke1}


%%%%%%%%%%%%%%%%%%%%%%%%%%%%%%%%%%%%%%%%
%% Size and Position of the Text Body %%
%%%%%%%%%%%%%%%%%%%%%%%%%%%%%%%%%%%%%%%%

% The new paper dimensions that are exclusively used.
\newlength{\mypaperwidth}
\setlength{\mypaperwidth}{210 mm}

\newlength{\mypaperheight}
\setlength{\mypaperheight}{297 mm}

% The text area uses aesthetically pleasing measurements in the same ratio as the page.
% These dimensions are always used, as the text area should be the same in the printed and digital version of the thesis.
\newlength{\mybodywidth}
\setlength{\mybodywidth}{140 mm}

\newlength{\mybodyheight}
\setlength{\mybodyheight}{198 mm}

\newlength{\myoutermargin}
\ifprintVersion
    \ifprofessionalPrint
        \setlength{\myoutermargin}{(\mypaperwidth - \mybodywidth) / \real{1.5} + \extraborderlength}
    \else
        \setlength{\myoutermargin}{(\mypaperwidth - \mybodywidth) / \real{1.5} - \mybindingcorrection}
    \fi
\else
    \setlength{\myoutermargin}{(\mypaperwidth - \mybodywidth) / \real{1.5}}
\fi

\newlength{\mytopmargin}
\setlength{\mytopmargin}{(\mypaperheight - \mybodyheight) / 3}
\ifprintVersion
    \ifprofessionalPrint
        \setlength{\mytopmargin}{(\mypaperheight - \mybodyheight) / 3 + \extraborderlength}
    \fi
\fi

\newlength{\myinnermargin}
\setlength{\myinnermargin}{\mypaperwidth - \mybodywidth - \myoutermargin}
\ifprintVersion
    \ifprofessionalPrint
        \setlength{\myinnermargin}{\mypaperwidth + \mybindingcorrection + 2\extraborderlength - \mybodywidth - \myoutermargin}
    \fi
\fi

\newlength{\mybottommargin}
\setlength{\mybottommargin}{\mypaperheight - \mybodyheight - \mytopmargin}
\ifprintVersion
    \ifprofessionalPrint
        \setlength{\mybottommargin}{\mypaperheight + 2\extraborderlength - \mybodyheight - \mytopmargin}
    \fi
\fi


%%%%%%%%%%%%%%%%%%%%%%%%%%%%%%%%%%%
%% Position of the Head And Foot %%
%%%%%%%%%%%%%%%%%%%%%%%%%%%%%%%%%%%

\newcommand{\goldenratio}{1.618}

\newlength{\myheadsep} % Distance from the header to the body.
\setlength{\myheadsep}{\mytopmargin / \real{\goldenratio} / \real{\goldenratio} - 1 ex}
\ifprintVersion
    \ifprofessionalPrint
        \setlength{\myheadsep}{(\mytopmargin - \extraborderlength) / \real{\goldenratio} / \real{\goldenratio} - 1 ex}
    \fi
\fi

\newlength{\myfootskip} % Distance from the body to the footer.
\setlength{\myfootskip}{\mybottommargin / \real{\goldenratio} - 1 ex}
\ifprintVersion
    \ifprofessionalPrint
        \setlength{\myfootskip}{(\mybottommargin - \extraborderlength) / \real{\goldenratio} - 1 ex}
    \fi
\fi


%%%%%%%%%%%%%%%%%%%%%%%%%%%%%%%
%% Margin Position And Width %%
%%%%%%%%%%%%%%%%%%%%%%%%%%%%%%%

\newlength{\mymargininnersep} % Distance between the margin and the body.
\setlength{\mymargininnersep}{7 mm}

\newlength{\mymarginoutersep} % Distance between the margin and the paper border.
\setlength{\mymarginoutersep}{12 mm}
\ifprintVersion
    \ifprofessionalPrint
        \setlength{\mymarginoutersep}{12 mm + \extraborderlength}
    \fi
\fi

\newlength{\mymarginwidth} % Width of the margin.
\setlength{\mymarginwidth}{\myoutermargin - \mymargininnersep - \mymarginoutersep}

\newlength{\mymarginwidthwithinnersep} % Width of the margin.
\setlength{\mymarginwidthwithinnersep}{\mymarginwidth + \mymargininnersep}

\usepackage
[
    % In the printed version, we add an extra border to each side as well as the binding correction for the width.
    \ifprintVersion
        \ifprofessionalPrint
            paperwidth = \mypaperwidth + 2\extraborderlength + \mybindingcorrection,
            paperheight = \mypaperheight + 2\extraborderlength,
        \else
            paperwidth = \mypaperwidth,
            paperheight = \mypaperheight,
        \fi
    \else
        paperwidth = \mypaperwidth,
        paperheight = \mypaperheight,
    \fi
    textwidth = \mybodywidth,
    textheight = \mybodyheight,
    outer = \myoutermargin,
    top = \mytopmargin,
    headsep = \myheadsep,
    footskip = \myfootskip,
    marginparsep = \mymargininnersep,
    marginparwidth = \mymarginwidth,
%    showframe, % Use this option for debugging purposes in order to the an outline of all of the different parts of the page layout.
]
{geometry} % Used in order to define the dimensions of the page and its layout.


%%%%%%%%%%%%%%%%%%%%%%%%%%%%%%
%% Header and Footer Format %%
%%%%%%%%%%%%%%%%%%%%%%%%%%%%%%

\usepackage
[
%    draft, % Shows a lot of rules denoting the dimensions of the head and foot. Use this option only for debugging.
]
{scrlayer-scrpage} % Allows to adjust the definitions of the head and foot of a page.

% Remove all predefined styles.
\clearpairofpagestyles

%%%%%%%%%%%%%%%%%%%%%%%%%%%%%%
% Dimensions and formats are defined.

% Define the dimensions of the head and the foot. Since we want some information to appear in the margin, we extend the head and the foot by the respective lengths.
\KOMAoptions
{%
    headwidth = \textwidth + \mymarginwidthwithinnersep,%
    footwidth = \myoutermargin : \textwidth,%
}

% Defines the formats for the chapter and section titles in the marks of the head.
\renewcommand*{\chaptermarkformat}{\normalfont\sffamily\small\color{footerchaptername}}
\renewcommand*{\sectionmarkformat}{\normalfont\sffamily\small\color{footersectionname}}

% Displays the chapter names in the head of both odd and even pages.
\automark[chapter]{chapter}
% Replaces the chapter name to the head of right pages with the section name if a section is present.
\automark*[section]{}

%%%%%%%%%%%%%%%%%%%%%%%%%%%%%%
% The head is defined.

% Head for even pages.
% Puts ›Chapter‹ followed by the current chapter number.
\lehead%
{%
    \begin{minipage}[b]{\mymarginwidth}%
        \ifnum\value{chapter}>0
            \small\raggedleft\normalfont\textsf{\textbf{\color{footerchapter}\chaptername\ \thechapter}}
        \fi
    \end{minipage}
}
% Put the title of the current chapter/section into the center of the head but push it to the border.
\cehead{\hspace*{\mymarginwidthwithinnersep}\parbox{\textwidth}{\raggedright\leftmark}}

% Head for odd pages.
\rohead%
{%
    % Check whether a section has already started or not.
    \Ifstr{\rightmark}{\leftmark}%
    {%
        \begin{minipage}[b]{\mymarginwidth}%
            \ifnum\value{chapter}>0
                \small\raggedright\normalfont\textsf{\textbf{\color{footersection}Chapter\ \thechapter}}%
            \fi
        \end{minipage}%
    }%
    {%
        \begin{minipage}[b]{\mymarginwidth}%
            \small\raggedright\normalfont\textsf{\textbf{\color{footersection}Section\ \thesection}}%
        \end{minipage}%
    }%
}
\cohead{\hspace*{-\mymarginwidthwithinnersep}\parbox{\textwidth}{\raggedleft\rightmark}}

%%%%%%%%%%%%%%%%%%%%%%%%%%%%%%
% The foot is defined.

% Displays the page number in bold in the margin, aligned toward the center. Further, a blue line is drawn above number.
% The starred variant is used, since we want the format of the foot to also apply to the pagestyle ›plain‹.
\lefoot*%
{%
    \vspace*{1 ex}%
    {\color{stroke1}\rule{\myoutermargin - \mymargininnersep}{0.5 mm}}\\
    \begin{minipage}[b]{\myoutermargin - \mymargininnersep}%
        \raggedleft\normalfont\color{footerpagenr}\textbf{\thepage}%
    \end{minipage}%
}
\rofoot*%
{%
    {\color{stroke1}\rule{\myoutermargin - \mymargininnersep}{0.5 mm}}\\
    \begin{minipage}[b]{\myoutermargin - \mymargininnersep}%
        \raggedright\normalfont\color{footerpagenr}\textbf{\thepage}%
    \end{minipage}%
}


%%%%%%%%%%%%%%%%%%%%
%% Caption Format %%
%%%%%%%%%%%%%%%%%%%%

\usepackage{caption}
\captionsetup
{
    font = small,
    labelfont = {bf, sf, color = captionlabel},
    format = plain,
    singlelinecheck = off,
}


%%%%%%%%%%%%%%%%%
%% Part Format %%
%%%%%%%%%%%%%%%%%
\usepackage{tikz} % Used in order to draw the stylistic elements.

\newlength{\mytmpa}
\setlength{\mytmpa}{1 mm}
\newlength{\mytmpb}
\newlength{\mytmpc}

%%%%%%%%%%%%%%%%%
% The following code draws the outline for a ›part‹ of the thesis.
% This command is used before the name of the part is displayed. It is void, as the part is added via \partlineswithprefixformat.
\renewcommand*{\partformat}{}
% This command calls \partformat (#2) and displays the name of the part (#3).
\renewcommand*{\partlineswithprefixformat}[3]%
{%
    #2
    \thispagestyle{empty}
    \ifprintVersion
        \setlength{\mytmpa}{0.618\mypaperwidth + \mybindingcorrection + \extraborderlength}%
        \setlength{\mytmpb}{0.382\mypaperheight + \extraborderlength}%
    \else
        \setlength{\mytmpa}{0.618\mypaperwidth}%
        \setlength{\mytmpb}{0.382\mypaperheight}%
    \fi
    \begin{tikzpicture}[overlay, remember picture]%
        \node [inner sep = 0, outer sep = 0, anchor = north] at (current page.north west)%
        {%
            \begin{tikzpicture}[overlay, remember picture]%
            \draw[color = stroke1, line width = 0.7 mm] (\mytmpa, 0) -- (\mytmpa, -\mytmpb);%
            \end{tikzpicture}%
        };%
        \node (align) [align = right, below = \mytmpb - 2 ex, inner sep = 0, outer sep = 0, anchor = north west] at (current page.north west)%
        {%
            \hspace{\mytmpa}\hspace{0.5 em}\partname\ \thepart\\[1 ex]
            \color{stroke1}#3%
        };%
    \end{tikzpicture}%
}
% This command defines various parameters for the ›part‹ format.
\RedeclareSectionCommand%
[%
    font = \normalfont\Huge\sffamily,
    prefixfont = \normalfont\Huge\sffamily,
]
{part}


%%%%%%%%%%%%%%%%%%%%%
%%% Chapter Format %%
%%%%%%%%%%%%%%%%%%%%%

\usepackage{etoolbox}
\usepackage[strict]{changepage}
\usetikzlibrary{calc}

\newbool{chapterHasANumber}
\newbool{chapterHasAStar}
\newcommand*{\thischapterlabel}{}
\renewcommand*{\chapterlinesformat}[3]%
{%
    % Check whether \chapter of \addchap has been used.
    \Ifnumbered{#1}{\setbool{chapterHasANumber}{true}}{\setbool{chapterHasANumber}{false}}%
    % Check whether \chapter* or \chapter has been used.
    \Ifstr{#2}{}{\setbool{chapterHasAStar}{true}}{\setbool{chapterHasAStar}{false}}%
    % Check whether a normal \chapter or something else is used.
    \ifboolexpr{bool{chapterHasANumber} and not bool{chapterHasAStar}}%
    {%
        \renewcommand*{\thischapterlabel}{\color{chapternumber}\thechapter}
    }%
    {%
        \renewcommand*{\thischapterlabel}{\vphantom{0}}
    }%
    \checkoddpage
    \begin{tikzpicture}[overlay, remember picture]%
    \ifoddpage%
        \coordinate (anchor) at ($(current page.north west) + (\myinnermargin, -\mytopmargin)$);
        \node [inner sep = 0, outer sep = 0, anchor = north west] (numbernode) at (anchor)%
        {%
            \sffamily\fontsize{60}{60}\selectfont\thischapterlabel
        };%
        \draw[color = stroke1, line width = 0.7 mm] ($(numbernode.south west)-(0, 1 ex)$) -- ++(\paperwidth, 0);%
        \ifboolexpr{bool{chapterHasANumber} and not bool{chapterHasAStar}}%
        {%
            \node (align) [text width = \textwidth - 2 cm, align = right, right = \mybodywidth, inner sep = 0, outer sep = 0, anchor = east] at (numbernode.west) {#3};
        }%
        {%
            \node (align) [align = left, inner sep = 0, outer sep = 0, anchor = west] at (numbernode.west) {#3};
        }%
    \else%
        \coordinate (anchor) at ($(current page.north east) - (\myinnermargin, \mytopmargin)$);
        \node [inner sep = 0, outer sep = 0, anchor = north east] (numbernode) at (anchor)%
        {%
            \sffamily\fontsize{60}{60}\selectfont\thischapterlabel
        };%
        \draw[color = stroke1, line width = 0.7 mm] ($(numbernode.south east)-(0, 1 ex)$) -- ++(-\paperwidth, 0);%
        \ifboolexpr{bool{chapterHasANumber} and not bool{chapterHasAStar}}%
        {%
            \node (align) [text width = \textwidth - 2 cm, align = left, left = \mybodywidth, inner sep = 0, outer sep = 0, anchor = west] at (numbernode.east) {#3};
        }%
        {%
            \node (align) [align = right, inner sep = 0, outer sep = 0, anchor = east] at (numbernode.east) {#3};
        }%
    \fi%
    \end{tikzpicture}%
}
\RedeclareSectionCommand%
[%
    font = \color{stroke1}\normalfont\huge\sffamily,
    afterskip = 20 pt,
]
{chapter}


%%%%%%%%%%%%%%%%%%%%%%%%
%%% Table of Contents %%
%%%%%%%%%%%%%%%%%%%%%%%%

% Format the table of contents to have a ›plain‹ page style.
\BeforeStartingTOC[toc]{\pagestyle{plain}}
\AfterStartingTOC{\thispagestyle{plain}}
                        % Contains commands that define the general format and layout of the thesis.
% This file contains all of the code that formats the bibliography. Since be package ›biblatex‹ is used, the bibliography needs to be compiled with ›biber‹.
%
% This file contains the following parts:
%   • Resources
%   • Redefined Keywords
%   • Coloring
%   • Format of the Entries
%   • Format of the Own Publications


\usepackage
[
    sortcites,              % Sort multiple references when citing them together.
    style = alphabetic,     % The style of a citation mark.
    defernumbers,           % Makes sure that references always have unique numbers. This is important if you use multiple bibliographies.
    safeinputenc,           % Allows to use UTF8 characters in the bibliography and tries to translate them into TeX automatically.
    backref = true,         % Creates back references in the bibliography.
    backrefstyle = three,   % Compresses three or more consecutive pages in the back references into a range.
    hyperref = true,        % Makes links generated by biblatex clickable. If hyperref is not used, a warning is issued.
    maxbibnames = 99,       % The maximum number of names displayed in the bibliography.
    maxcitenames = 2,       % The maximum number of names displayed when using commands like ›textcite‹. The default is 3. After that, ›et al.‹ is used.
%    useprefix,              % Prints name prefixes, such as ›von‹. The default is false. This means that prefixes are not considered to be part of the last name.
]
{biblatex} % Used in order to format the bibliography.

% The following command changes the space between the list of authors and the citation mark into a non-breaking space.
\renewcommand\namelabeldelim{\addnbspace}


%%%%%%%%%%%%%%%
%% Resources %%
%%%%%%%%%%%%%%%

\addbibresource{references/strings.bib}                     % Contains many strings for common conference names etc. These strings can then be used in the references.
\addbibresource{references/references.bib}                  % The file that contains the references that are used for the thesis.


%%%%%%%%%%%%%%%%%%%%%%%%
%% Redefined Keywords %%
%%%%%%%%%%%%%%%%%%%%%%%%

\renewbibmacro{in:}%
{%
    \ifentrytype{article}{}{\printtext{\bibstring{in}\intitlepunct}}%
}
% \renewcommand*{\volumenumberdelim}{\addcolon}

\renewbibmacro*{volume+number+eid}%
{%
    \printfield{volume}%
    \iffieldundef{number}{}{\addcolon}%
    %  \setunit*{\addnbthinspace}%
    \printfield{number}%
    \setunit*{\addcomma\space}%
    \printfield{eid}%
}

\DefineBibliographyStrings{english}%
{%
    backrefpage  = {\lowercase{s}ee page}, % For a single page number.
    backrefpages = {\lowercase{s}ee pages} % For multiple page numbers.
}


%%%%%%%%%%%%%%
%% Coloring %%
%%%%%%%%%%%%%%

\DeclareFieldFormat[article]{title}{\textbf{\color{stroke1}#1}}
\DeclareFieldFormat[inproceedings]{title}{\textbf{\color{stroke1}#1}}
\DeclareFieldFormat[thesis]{title}{\textbf{\color{stroke1}#1}}
\DeclareFieldFormat[book]{title}{\textbf{\color{stroke1}#1}}
\DeclareFieldFormat[unpublished]{title}{\textbf{\color{stroke1}#1}}
\DeclareFieldFormat[report]{title}{\textbf{\color{stroke1}#1}}
\DeclareFieldFormat[inbook]{chapter}{\textbf{\color{stroke1}#1}}
\DeclareFieldFormat[inbook]{title}{#1}
\DeclareFieldFormat{pages}{#1}


%%%%%%%%%%%%%%%%%%%%%%%%%%%
%% Format of the Entries %%
%%%%%%%%%%%%%%%%%%%%%%%%%%%

% The following toggle defines how the citation mark formats the author names. If this toggle is true, more information is used.
\newtoggle{authorend}
\togglefalse{authorend}

% Article
\DeclareBibliographyDriver{article}%
{%
  \usebibmacro{bibindex}%
  \usebibmacro{begentry}%
  \iftoggle{authorend}{}{\usebibmacro{author/translator+others}}%
  \setunit{\labelnamepunct}\newblock
  \usebibmacro{title}%
  \newunit
  \printlist{language}%
  \newunit\newblock
  \usebibmacro{byauthor}%
  \newunit\newblock
  \usebibmacro{bytranslator+others}%
  \newunit\newblock
  \printfield{version}%
  \newunit\newblock
  \usebibmacro{in:}%
  \usebibmacro{journal+issuetitle}%
  \newunit
  \usebibmacro{byeditor+others}%
  \newunit
  \usebibmacro{note+pages}%
  \newunit\newblock
  \iftoggle{bbx:isbn}
  {\printfield{issn}}
  {}%
  \newunit\newblock
  \usebibmacro{doi+eprint+url}%
  \newunit\newblock
  \usebibmacro{addendum+pubstate}%
  \setunit{\bibpagerefpunct}\newblock
  \usebibmacro{pageref}%
  \newunit\newblock
  \iftoggle{bbx:related}
  {\usebibmacro{related:init}%
    \usebibmacro{related}}
  {}%
  \usebibmacro{finentry}%
  \iftoggle{authorend}{\usebibmacro{author/translator+others}}{}%
}

% Book Chapter
\DeclareBibliographyDriver{inbook}%
{%
  \usebibmacro{bibindex}%
  \usebibmacro{begentry}%
  \iftoggle{authorend}{}{\usebibmacro{author/translator+others}}%
  \setunit{\labelnamepunct}\newblock
  % \usebibmacro{title}%
  \usebibmacro{chapter+pages}%
  % \printfield{chapter}%
  \newunit
  \printlist{language}%
  \newunit\newblock
  \usebibmacro{byauthor}%
  \newunit\newblock
  \usebibmacro{in:}%
  \usebibmacro{bybookauthor}%
  \newunit\newblock
  \usebibmacro{maintitle+booktitle}%
  \newunit\newblock
  \usebibmacro{byeditor+others}%
  \newunit\newblock
  \printfield{edition}%
  \newunit
  \iffieldundef{maintitle}
  {\printfield{volume}%
    \printfield{part}}
  {}%
  \newunit
  \printfield{volumes}%
  \newunit\newblock
  \usebibmacro{series+number}%
  \newunit\newblock
  \printfield{note}%
  \newunit\newblock
  \usebibmacro{publisher+location+date}%
  \newunit\newblock
  % \usebibmacro{chapter+pages}%
  \newunit\newblock
  \iftoggle{bbx:isbn}
  {\printfield{isbn}}
  {}%
  \newunit\newblock
  \usebibmacro{doi+eprint+url}%
  \newunit\newblock
  \usebibmacro{addendum+pubstate}%
  \setunit{\bibpagerefpunct}\newblock
  \usebibmacro{pageref}%
  \newunit\newblock
  \iftoggle{bbx:related}
  {\usebibmacro{related:init}%
    \usebibmacro{related}}
  {}%
  \usebibmacro{finentry}%
  \iftoggle{authorend}{\usebibmacro{author/translator+others}}{}%
}

% Proceedings Article
\DeclareBibliographyDriver{inproceedings}%
{%
  \usebibmacro{bibindex}%
  \usebibmacro{begentry}%
  \iftoggle{authorend}{}{\usebibmacro{author/translator+others}}%
  \setunit{\labelnamepunct}\newblock
  \usebibmacro{title}%
  \newunit
  \printlist{language}%
  \newunit\newblock
  \usebibmacro{byauthor}%
  \newunit\newblock
  \usebibmacro{in:}%
  \usebibmacro{maintitle+booktitle}%
  \newunit\newblock
  \usebibmacro{event+venue+date}%
  \newunit\newblock
  \usebibmacro{byeditor+others}%
  \newunit\newblock
  \iffieldundef{maintitle}
  {\printfield{volume}%
    \printfield{part}}
  {}%
  \newunit
  \printfield{volumes}%
  \newunit\newblock
  \usebibmacro{series+number}%
  \newunit\newblock
  \printfield{note}%
  \newunit\newblock
  \printlist{organization}%
  \newunit
  \usebibmacro{publisher+location+date}%
  \newunit\newblock
  \usebibmacro{chapter+pages}%
  \newunit\newblock
  \iftoggle{bbx:isbn}
  {\printfield{isbn}}
  {}%
  \newunit\newblock
  \usebibmacro{doi+eprint+url}%
  \newunit\newblock
  \usebibmacro{addendum+pubstate}%
  \setunit{\bibpagerefpunct}\newblock
  \usebibmacro{pageref}%
  \newunit\newblock
  \iftoggle{bbx:related}
  {\usebibmacro{related:init}%
    \usebibmacro{related}}
  {}%
  \usebibmacro{finentry}%
  \iftoggle{authorend}{\usebibmacro{author/translator+others}}{}%
}

% Thesis
\DeclareBibliographyDriver{thesis}%
{%
  \usebibmacro{bibindex}%
  \usebibmacro{begentry}%
  \iftoggle{authorend}{}{\usebibmacro{author}}%
  \setunit{\labelnamepunct}\newblock
  \usebibmacro{title}%
  \newunit
  \printlist{language}%
  \newunit\newblock
  \usebibmacro{byauthor}%
  \newunit\newblock
  \printfield{note}%
  \newunit\newblock
  \printfield{type}%
  \newunit
  \usebibmacro{institution+location+date}%
  \newunit\newblock
  \usebibmacro{chapter+pages}%
  \newunit
  \printfield{pagetotal}%
  \newunit\newblock
  \iftoggle{bbx:isbn}
  {\printfield{isbn}}
  {}%
  \newunit\newblock
  \usebibmacro{doi+eprint+url}%
  \newunit\newblock
  \usebibmacro{addendum+pubstate}%
  \setunit{\bibpagerefpunct}\newblock
  \usebibmacro{pageref}%
  \newunit\newblock
  \iftoggle{bbx:related}
  {\usebibmacro{related:init}%
    \usebibmacro{related}}
  {}%
  \usebibmacro{finentry}%
  \iftoggle{authorend}{\usebibmacro{author}}{}%
}


%%%%%%%%%%%%%%%%%%%%%%%%%%%%%%%%%%%%
%% Format of the Own Publications %%
%%%%%%%%%%%%%%%%%%%%%%%%%%%%%%%%%%%%

% The own publications are formatted using a numeric list, whereas the bibliography of the thesis uses an alphanumeric style.

% Copied from numeric.cbx in order to imitate numerical citations.
\providebool{bbx:subentry}
\newbibmacro*{citenum}%
{% Note: the original macro was called ›cite‹. I did not redefine ›cite‹ but instead defined a new macro ›citenum‹ because the author-year citations use the ›cite‹ macro too. Using ›\renewbibmacro*{cite}‹ would have caused all the author-year citations to become numeric too.
  \printtext[bibhyperref]{% If you ever want to use hyperref.
    \printfield{prefixnumber}%
    \printfield{labelnumber}%
    \ifbool{bbx:subentry}
    {\printfield{entrysetcount}}
    {}}%
}

% Copied from numeric.cbx to define a new numeric citation command for @online entries.
\DeclareCiteCommand{\conline}[\mkbibbrackets]
{\usebibmacro{prenote}}
{\usebibmacro{citeindex}%
  \usebibmacro{citenum}}% Note: this was originally "cite" but I changed it to "citenum" to avoid clashes with the author-year style.
{\multicitedelim}
{\usebibmacro{postnote}}       % Contains commands for the layout of the bibliography.
% This file contains most of the packages used for this document. If you want to add a package, do it here.
% Some packages are already included in other files in the ›core‹ folder if they were already necessary. Thus, make sure to go through these files too if you want to know whether a certain package is already included.
%
% This file contains the following parts:
%   • Typography
%   • Math
%   • Fonts
%   • Graphics
%   • Tables
%   • Enumerations
%   • Algorithms
%   • Spaces and Special Characters
%   • Miscellaneous
%   • Additional Packages
%   • Hyperlinks

%%%%%%%%%%%%%%%%
%% Typography %%
%%%%%%%%%%%%%%%%

\usepackage
[
    babel = true, % Enables language-specific tuning.
]
{microtype}           % Uses the text space more efficiently.
\usepackage{csquotes} % Uses the correct quotes according to the current language.


%%%%%%%%%%
%% Math %%
%%%%%%%%%%

% The following packages are the standard packages used in order to typeset math. They contain a lot of useful commands.
\usepackage{amsmath}
\usepackage{amssymb}
\usepackage{amsthm}
\usepackage{thmtools}
\usepackage{mathtools}
\usepackage{thm-restate}
\usepackage{dsfont}        % Yields far better blackboard-bold letters than \mathbb. Use \mathds in order to write such letters.
\usepackage{braceMnSymbol} % Adjusts overbraces and underbraces such that longer versions are put together seamlessly.


%%%%%%%%%%%
%% Fonts %%
%%%%%%%%%%%

\usepackage
[
    ttscale = 0.85, % Scales the typewriter font.
]
{libertine} % The main font used in this thesis.
\usepackage
[
    libertine,    % Changes the math font to libertine (the main font).
    slantedGreek, % Makes all greek letters italic by default. If you want to use an upright greek letter, use ›\up‹ immediately followed by the letter’s name. For example, \upGamma displays an upright uppercase gamma.
    vvarbb,       % Changes the \mathbb font to another font. However, \mathbb remains ugly and should not be used. Use \mathds instead.
    libaltvw,     % Uses different characters for v und w that look far better than the default ones.
]
{newtxmath} % The main math font of this thesis. It fits well with the main font.
\usepackage{url} % Responsible for URL formatting.
\usepackage{bm}  % Allows to use sensible bold letters in math mode. This package has to go after the font packages. Otherwise it does not work correctly!


%%%%%%%%%%%%%%
%% Graphics %%
%%%%%%%%%%%%%%

\usepackage{graphicx} % The standard package for including graphics into your document.
\usepackage
[
    subrefformat = simple, % Formats the label of the \subref command without parentheses.
    labelformat = simple,  % Formats the mark of a subfigure without parentheses.
]
{subcaption}         % Enables it to have subfigures inside of a single figure.
\usepackage{wrapfig} % Allows to put figures next to text.

% Changing the \columnsep adds some space next to a warpfigure.
\columnsep = \mymargininnersep
% The reference label of a subfigure is redefined to have a non-breaking space and parentheses. (Thus, the subfigures show parentheses although the package options removed parentheses; otherwise, two pairs of brackets would be seen.)
\renewcommand*{\thesubfigure}{~(\alph{subfigure})}


%%%%%%%%%%%%
%% Tables %%
%%%%%%%%%%%%

\usepackage{array}     % Improves the way that tables can be formatted.
\usepackage{booktabs}  % Adds lines (called ›rules‹) that can be used in tables and improves spacing.
\usepackage{longtable} % Allows to make tables that span multiple pages.
\usepackage{pdflscape} % Allows to change a page into landscape. This is handy if a table is very wide.


%%%%%%%%%%%%%%%%%%
%% Enumerations %%
%%%%%%%%%%%%%%%%%%

\usepackage{enumitem} % Adds tons of useful features to enumeration environments.


%%%%%%%%%%%%%%%%
%% Algorithms %%
%%%%%%%%%%%%%%%%

\usepackage
[
    ruled,         % Creates lines at the top and at the bottom. Further, the caption is now above the algorithm.
    vlined,        % Shows the scope of a statement spanning multiple lines via a small vertical bar. Thus, no closing tags are needed.
    linesnumbered, % Shows line numbers.
]
{algorithm2e} % Allows to write pseudocode.


%%%%%%%%%%%%%%%%%%%%%%%%%%%%%%%%%%%
%% Spaces and Special Characters %%
%%%%%%%%%%%%%%%%%%%%%%%%%%%%%%%%%%%

\usepackage{xspace}   % Adds the functionality that a space after a command will be shown as a space in the output.
\usepackage
[
    shortcuts, % Allows to use short symbols for non-breaking hyphens and dashes instead of lengthy commands.
]
{extdash}             % Adds non-breaking hyphens and dashes.
\usepackage{setspace} % Allows to easily chnage the spacing inside of the document.


%%%%%%%%%%%%%%%%%%%
%% Miscellaneous %%
%%%%%%%%%%%%%%%%%%%

\usepackage{xparse}    % Is used in order to define reasonable commands.
\usepackage{footnote}  % Allows it to extend the environments footnotes can be used in. It is said that this package is in conflict with ›hyperref‹. I did not note any troubles. However, if something is fishy, it is probably best to not use this package.
\usepackage{afterpage} % Adds the \afterpage command, which specifies that the provided argument shall be processed after the current page is finished.
\usepackage
[
    textsize = scriptsize, % Determines the text size of the TODO note.
]
{todonotes}            % Adds TODO notes to the document. These are small text areas inside of the margin of a page.


%%%%%%%%%%%%%%%%%%%%%%%%%
%% Additional Packages %%
%%%%%%%%%%%%%%%%%%%%%%%%%

% Add additional packages you would like to use here.
\usepackage{wasysym}
\usepackage{makecell}
\usepackage{multirow}
\usepackage{colortbl}
\usepackage{tabularx}

% Figures in figures/fig_*.tex are TikZ rather than bitmaps, so that they are
% vector and set in the body font. They are generated by
% scripts/generate_thesis_figures.py from the result JSONs; edit that script and
% re-run it rather than editing the generated files. TikZ itself is already
% loaded by core/preamble/format.tex, which also loads the calc library; the
% arrow tips used by the pipeline diagram need one more.
\usetikzlibrary{arrows.meta}

% One palette for every figure. Okabe-Ito: distinguishable under the common
% forms of colour blindness and distinct in greyscale, which matters because
% the thesis will be printed.
\definecolor{kgblue}{HTML}{0072B2}
\definecolor{ragorange}{HTML}{E69F00}
\definecolor{hybridgreen}{HTML}{009E73}
\definecolor{modegrey}{HTML}{6E6E6E}


%%%%%%%%%%%%%%%%
%% Hyperlinks %%
%%%%%%%%%%%%%%%%

\usepackage
[
    bookmarks = true,                 % Generates boodmarks for the PDF.
    bookmarksopen = false,            % The bookmarks are closed by default.
    bookmarksnumbered = true,         % The bookmarks use the numbers of the corresponding headline.
    pdfstartpage = 1,                 % The first page seen when opening the PDF.
    pdftitle = {{\printTitle}},       % The PDF’s title in the meta data.
    pdfauthor = {{\printAuthor}},     % The PDF’s author name in the meta data.
    pdfsubject = {{\printSubject}},   % The PDF’s subject in the meta data.
    pdfkeywords = {{\printKeywords}}, % The PDF’s keywords in the meta data.
    breaklinks = true,                % Allows it to break links.
    \ifprintVersion
        hidelinks,                    % In the printed version, links are not highlighted, as they are not clickable.
    \else
    colorlinks = true,            % The text of hyperlinks is colored instead of having a colored box around it.
    allcolors = stroke1,          % Every hyperlink uses the same color. If you want to change specific colors, use the commands below.
    %        linkcolor = stroke1,          % The color of an in-document hyperlink.
    %        citecolor = stroke1,          % The color of a citation.
    %        filecolor = stroke1,          % The color of a file link.
    %        pagecolor = stroke1,          % The color of a reference to a page.
    %        urlcolor = stroke1,           % The color of a weblink.
    \fi
]
{hyperref} % The standard package that is used for creating hyperlinks inside of a document.

\usepackage
[
    %    capitalise, % Capitalizes the words in front of the labels. This can also be done by simply using \Cref instead of \cref. In order to have a greater variety, this option is not used.
    noabbrev,   % The words in front of the labels are not abbreviated.
    nameinlink, % Extends the link of a reference to the word in front of it.
]
{cleveref} % This package must be included after ›hyperref‹. It creates clever references that know what they refer to.     % Contains the packages that this template provides.
% This file contains all sorts of macros that are globally used. Further, certain options made available through packages are set here as well.
%
% This file contains the following parts:
%   • Type of Degree
%   • Miscellaneous
%   • Footnotes
%   • Theorem Environments
%   • Meta Commands
%   • Common Commands


%%%%%%%%%%%%%%%%%%%%
%% Type of Degree %%
%%%%%%%%%%%%%%%%%%%%

% The colloquial term of the degree.
\newcommand*{\colloquialDegreeName}{Master}
\newcommand*{\colloquialDegreeNameLowercase}{master}

% The abbreviation of the degree.
\newcommand*{\degreeAbbreviation}{M.}

% Redefine the two macors above for a bachelor thesis.
\ifbachelorThesis
    \renewcommand*{\colloquialDegreeName}{Bachelor}
    \renewcommand*{\colloquialDegreeNameLowercase}{bachelor}
    \renewcommand*{\degreeAbbreviation}{B.}
\fi


%%%%%%%%%%%%%%%%%%%
%% Miscellaneous %%
%%%%%%%%%%%%%%%%%%%

% Defines the environment used at the beginning of each chapter.
\newenvironment{jointwork}
{\itshape}
{\ignorespacesafterend\bigskip}

% Defines the IfEmptyTF command. This is useful for optional arguments provided as [].
\makeatletter
    \def\IfEmptyTF#1%
    {%
        \if\relax\detokenize{#1}\relax%
            \expandafter\@firstoftwo%
        \else%
            \expandafter\@secondoftwo%
        \fi%
    }
\makeatother

% Creates an environment that automatically uses math mode if necessary and creates a space afterward if wanted. Basically, if the command \example is defined to use this environment, you can use \example without mathe mode in normal text as if it were ordinary text.
\NewDocumentCommand{\mathOrText}{m}
{%
    \ensuremath{#1}\xspace%
}

% Reduces the space around scaling bracekts.
\let\originalleft\left
\let\originalright\right
\renewcommand{\left}{\mathopen{}\mathclose\bgroup\originalleft}
\renewcommand{\right}{\aftergroup\egroup\originalright}

% Lets math text in an environment of bold text also appear bold.
\makeatletter
    \DeclareRobustCommand{\bfseries}%
    {%
        \not@math@alphabet\bfseries\mathbf%
        \fontseries\bfdefault\selectfont%
        \boldmath%
    }
\makeatother

% Adds square and curly brackets to the exceptions for xspace such that no space is used right in front of them.
\xspaceaddexceptions{]\}}

% Formats URLs by using the normal font (not the typewriter font).
\urlstyle{rm}

% Allows large display formulas to span multiple pages.
\allowdisplaybreaks

% Defines an optional argument for labels named ›ineq‹ that signals that cleveref should name the respective reference ›inequality‹ instead of its actual name.
\crefname{ineq}{inequality}{inequalities}
\creflabelformat{ineq}{#2{\upshape(#1)}#3} 

% Defines an optional argument for labels named ›term‹ that signals that cleveref should name the respective reference ›term‹ instead of its actual name.
\crefname{term}{term}{terms}
\creflabelformat{term}{#2{\upshape(#1)}#3}


%%%%%%%%%%%%%%%
%% Footnotes %%
%%%%%%%%%%%%%%%

% In the following, the command ›footnote‹ is redefined such that the footnote mark can be more easily adjusted.
\let\oldfootnote\footnote

% The following are variables used by the command.
\newlength{\spaceBeforeFootnote} % Denotes the space before the footnote mark in em.
\newlength{\spaceAfterFootnote}  % Denotes the space after the footnote mark in em.

% The new footnote command. The first three arguments are optional, the fourth mandatory. Its arguments have the following meaning:
%   1. The amount of space before the footnote mark in em. The default is 0.
%   2. The amount of space after the footnote mark in em. The default is 0.
%   3. The number of the footnote mark.
%   4. The text of the footnote.
\RenewDocumentCommand{\footnote}{o o o m}%
{%
    \IfNoValueTF{#1}%
    {%
        \oldfootnote{#4}%
    }%
    {%
        \setlength{\spaceBeforeFootnote}{\IfEmptyTF{#1}{0}{#1} em}%
        \IfNoValueTF{#2}%
        {%
            \hspace*{\spaceBeforeFootnote}\oldfootnote{#4}%
        }%
        {%
            \setlength{\spaceAfterFootnote}{\IfEmptyTF{#2}{0}{#2} em}%
            \hspace*{\spaceBeforeFootnote}\IfNoValueTF{#3}{\oldfootnote{#4}}{\oldfootnote[#3]{#4}}\hspace*{\spaceAfterFootnote}%
        }%
    }%
}

% The following commands enable it such that footnotes can be used in various other environments other than simple text.
\makesavenoteenv{figure}
\makesavenoteenv{table}
\makesavenoteenv{tabular}


%%%%%%%%%%%%%%%%%%%%%%%%%%
%% Theorem Environments %%
%%%%%%%%%%%%%%%%%%%%%%%%%%

\iffancyTheorems
    % The following theorem style uses a bold heading for the theorem and normal (upright) text. The environment begins with a triangle of color ›stroke1‹ pointing to the right and uses a QED symbol that is a triangle of the same color pointing to the left. Thus, the environment is enclosed by triangles.
    \declaretheoremstyle
    [
        spaceabove = \topsep,
        spacebelow = \topsep,
        headfont = \bfseries,
        headformat = \textcolor{stroke1}{$\blacktriangleright$} \NAME~\NUMBER \NOTE,
        notefont = \bfseries,
        notebraces = {(}{)},
        bodyfont = \normalfont,
        postheadspace = 0.5 em,
        qed = \textcolor{stroke1}{\bfseries$\blacktriangleleft$},
    ]
    {myTheoremStyle}
    
    % The QED symbol used in proofs is a squre with color ›stroke1‹ in order to look similar to the theorem environments.
    \renewcommand*{\qedsymbol}{\textcolor{stroke1}{$\blacksquare$}}
    
    \declaretheorem
    [
        style = myTheoremStyle,
        name = Conjecture,
        numberwithin = chapter,
    ]
    {conjecture}
    \declaretheorem
    [
        style = myTheoremStyle,
        name = Proposition,
        sharenumber = conjecture,
    ]
    {proposition}
    \declaretheorem
    [
        style = myTheoremStyle,
        name = Claim,
        sharenumber = conjecture,
    ]
    {claim}
    \declaretheorem
    [
        style = myTheoremStyle,
        name = Lemma,
        sharenumber = conjecture,
    ]
    {lemma}
    \declaretheorem
    [
        style = myTheoremStyle,
        name = Corollary,
        sharenumber = conjecture,
    ]
    {corollary}
    \declaretheorem
    [
        style = myTheoremStyle,
        name = Theorem,
        sharenumber = conjecture,
    ]
    {theorem}
    \declaretheorem
    [
        style = myTheoremStyle,
        name = Definition,
        sharenumber = conjecture,
    ]
    {definition}
    \declaretheorem
    [
        style = myTheoremStyle,
        name = Example,
        sharenumber = conjecture,
    ]
    {example}
    \declaretheorem
    [
        style = myTheoremStyle,
        name = Remark,
        sharenumber = conjecture,
    ]
    {remark}
\else
    % This is the default style. That is, the head is bold, the rest is italic, and there is no symbol to denote the end of the environment.
    \theoremstyle{plain}
    
    \newtheorem{conjecture}{Conjecture}[chapter]
    \newtheorem{proposition}[conjecture]{Proposition}
    \newtheorem{claim}[conjecture]{Claim}
    \newtheorem{lemma}[conjecture]{Lemma}
    \newtheorem{corollary}[conjecture]{Corollary}
    \newtheorem{theorem}[conjecture]{Theorem}
    \newtheorem{definition}[conjecture]{Definition}
    \newtheorem{example}[conjecture]{Example}
    \newtheorem{remark}[conjecture]{Remark}
\fi


%%%%%%%%%%%%%%%%%%%
%% Meta Commands %%
%%%%%%%%%%%%%%%%%%%

% A template for a function that can use an optional variable bracket size. Its arguments have the following meaning:
%   1. The name of the function.
%   2. The type of the left bracket. This should be a bracket symbol, as it will be forwarded to the command \left.
%   3. The type of the right bracket. The same restrictions as with parameter 2 hold here.
%   4. The arguments that the function takes, that is, the things that are enclosed by the brackets.
%   5. The size of the brackets. This should be a value like \big or similar, as it will be forwarded to the command \left.
\NewDocumentCommand{\functionTemplate}{m m m m o}%
{%
    \IfNoValueTF{#5}%
    {%
        \mathOrText{#1\left#2{#4}\right#3}%
    }%
    {%
        \mathOrText{#1#5#2{#4}#5#3}%
    }%
}

% The following two commands are used as variables for the following command.
\newcommand*{\leftBracketType}{(}
\newcommand*{\rightBracketType}{)}

% This is a command that creates a command that is a function as defined by the command \functionTemplate. Its arguments have the following meaning:
%   1. The name of the function command.
%   2. The name of the function itself.
%   3. The type of the left bracket. This will be forwarded to parameter 2 of \functionTemplate. The default is (. Use \lbrack for [ and \{ for }.
%   4. The type of the right bracket. This will be forwarded to parameter 3 of \functionTemplate. The default is ). The rest is similar to parameter 3.
% The command created has two optional arguments, which are as follows:
%   1. The arguments of the function. If this is empty, only the name of the function will be used.
%   2. The size of the brackets. This will be forwarded to parameter 5 of \functionTemplate.
\NewDocumentCommand{\createFunction}{m m o o}%
{%
    \renewcommand*{\leftBracketType}{\IfNoValueTF{#3}{(}{#3}}%
    \renewcommand*{\rightBracketType}{\IfNoValueTF{#4}{)}{#4}}%
    \NewDocumentCommand{#1}{o o}%
    {%
        \IfNoValueTF{##1}%
        {%
            \mathOrText{#2}%
        }%
        {%
            \functionTemplate{#2}{\leftBracketType}{\rightBracketType}{##1}[##2]%
        }%
    }%
}

% A template for a probabilistic symbol, which can make use of a condition denoted by |. Its arguments have the following meaning:
%   1. The name of the function.
%   2. The argument of the function.
%   3. The condition of the function. The default is that there is no condition.
%   4. The size of the brackets. This will be forwarded to parameter 5 of \functionTemplate.
\DeclareDocumentCommand{\probabilisticFunctionTemplate}{m m O{} o}
{%
    \functionTemplate{#1}%
    {\lbrack}%
    {\rbrack}%
    {#2\IfEmptyTF{#3}{}{\ \IfNoValueTF{#4}{\left}{#4}\vert\ \vphantom{#2}#3\IfNoValueTF{#4}{\right.}{}}}%
    [#4]%
}


%%%%%%%%%%%%%%%%%%%%%
%% Common Commands %%
%%%%%%%%%%%%%%%%%%%%%

%%%%%%%%%%%%%%%%%%%%%
% Number Sets

% Number sets appear in bold by default. The other option is to make them appear in blackboard bold.
\ifboldNumberSets
    \newcommand*{\N}{\mathOrText{\mathbf{N}}}
    \newcommand*{\Z}{\mathOrText{\mathbf{Z}}}
    \newcommand*{\Q}{\mathOrText{\mathbf{Q}}}
    \newcommand*{\R}{\mathOrText{\mathbf{R}}}
    \newcommand*{\C}{\mathOrText{\mathbf{C}}}
    \newcommand*{\indicatorFunctionSymbol}{\mathbf{1}}
\else
    \newcommand*{\N}{\mathOrText{\mathds{N}}}
    \newcommand*{\Z}{\mathOrText{\mathds{Z}}}
    \newcommand*{\Q}{\mathOrText{\mathds{Q}}}
    \newcommand*{\R}{\mathOrText{\mathds{R}}}
    \newcommand*{\C}{\mathOrText{\mathds{C}}}
    \newcommand*{\indicatorFunctionSymbol}{\mathds{1}}
\fi

%%%%%%%%%%%%%%%%%%%%%
% Probabilistic Functions
% All of these functions follow the outline of \probabilisticFunctionTemplate. That is, the syntax is, for example, \Pr{A}[B][\big], which would be shown as Pr[A | B] with \big brackets.

% Probability measure
\RenewDocumentCommand{\Pr}{m O{} o}%
{%
    \probabilisticFunctionTemplate{\mathrm{Pr}}{#1}[#2][#3]%
}

% Expected value
\NewDocumentCommand{\E}{m O{} o}%
{%
    \probabilisticFunctionTemplate{\mathrm{E}}{#1}[#2][#3]%
}

% Variance
\NewDocumentCommand{\Var}{m O{} o}%
{%
    \probabilisticFunctionTemplate{\mathrm{Var}}{#1}[#2][#3]%
}

%%%%%%%%%%%%%%%%%%%%%
% Landau Notation
% The following commands all take a mandatory argument, which is the term of the Landau notation, as well as an optional argument, which determines the size of the brackets.

% Big O
\DeclareDocumentCommand{\bigO}{m o}%
{%
    \functionTemplate{\mathrm{O}}{(}{)}{#1}[#2]%
}

% Small O
\DeclareDocumentCommand{\smallO}{m o}%
{%
    \functionTemplate{\mathrm{o}}{(}{)}{#1}[#2]%
}

% Big Theta
\DeclareDocumentCommand{\bigTheta}{m o}%
{%
    \functionTemplate{\upTheta}{(}{)}{#1}[#2]%
}

% Big Omega
\DeclareDocumentCommand{\bigOmega}{m o}%
{%
    \functionTemplate{\upOmega}{(}{)}{#1}[#2]%
}

% Small Omega
\DeclareDocumentCommand{\smallOmega}{m o}%
{%
    \functionTemplate{\upomega}{(}{)}{#1}[#2]%
}

%%%%%%%%%%%%%%%%%%%%%
% Constants

% Pi; ratio of a circle’s circumference to its diameter
\newcommand*{\circlePi}{\mathOrText{\uppi}}

% Euler’s constant. This command takes an optional parameter, which becomes the exponent of this constant.
\DeclareDocumentCommand{\eulerE}{o}%
{%
    \mathOrText{\mathrm{e}\IfNoValueTF{#1}{}{^{#1}}}%
}

% i; the imaginary unit
\newcommand*{\imaginaryUnit}{\mathOrText{\mathrm{i}}}

%%%%%%%%%%%%%%%%%%%%%
% Other

% A polynomial function. The mandatory parameter is the argument of the function, the optional one is the size of the brackets.
\DeclareDocumentCommand{\poly}{m o}%
{%
    \functionTemplate{\mathrm{poly}}{(}{)}{#1}[#2]%
}

% The identity function
\createFunction{\id}{\mathrm{id}}

% An indicator function. The first parameter is set as an index, the second is the argument of the function, and the third is the size of the brackets.
\NewDocumentCommand{\ind}{m o o}%
{%
    \IfNoValueTF{#2}%
    {%
        \mathOrText{\indicatorFunctionSymbol_{#1}}%
    }%
    {%
        \functionTemplate{\indicatorFunctionSymbol_{#1}}{(}{)}{#2}[#3]%
    }%
}

% The domain of a function. Its parameters are the same as for \poly.
\DeclareDocumentCommand{\dom}{m o}%
{%
    \functionTemplate{\mathrm{dom}}{(}{)}{#1}[#2]%
}

% The range of a function. Its parameters are the same as for \poly.
\DeclareDocumentCommand{\rng}{m o}%
{%
    \functionTemplate{\mathrm{rng}}{(}{)}{#1}[#2]%
}

% The d for an integral. The optional parameter becomes the exponent/degree of the operator.
\DeclareDocumentCommand{\d}{o}%
{%
    \mathrm{d}\IfNoValueTF{#1}{}{^{#1}}%
}

% A command that creates sets. The first parameter is the left-hand side, the second is the right-hand side, and the third (optional) parameter is the size of the brackets.
\DeclareDocumentCommand{\set}{m m o}%
{
    \mathOrText{\IfNoValueTF{#3}{\left}{#3}\{#1\ \IfNoValueTF{#3}{\left}{#3}\vert\
    \vphantom{#1}#2\IfNoValueTF{#3}{\right.}{}\IfNoValueTF{#3}{\right}{#3}\}}
}      % Contains newly defined commands useful for mathematics.
% This is where all the commands should go that you want to define yourself.
\newcommand{\ie}{i.e.,\xspace}
\newcommand{\eg}{e.g.,\xspace}

% Review-generation modes. Defined as macros so a rename stays consistent
% across every chapter; the four names are the unit of comparison throughout.
\newcommand{\baseline}{baseline}
\newcommand{\kg}{KG}
\newcommand{\rag}{RAG}
\newcommand{\hybrid}{hybrid}

\newcommand{\kappacoef}{Cohen's~$\kappa$}

% The drafting notes below cite source paths such as results/BOOTSTRAP_STATS_v2.md,
% and a bare underscore is illegal in text mode. This package makes it print literally
% in text while leaving subscripts in math mode alone.
\usepackage{underscore}

 % This is where user-defined commands should go.


% This is the thesis. The front matter as well as the references should not be changed. The main matter can be edited freely.
\begin{document}


    \frontmatter
    % This file contains the layout of the title page. It is a generous interpretation of the demo page provided by the MNF of the University of Potsdam.

% This page uses a different geometry, as the content will be centered (not including the binding correction).
\ifprintVersion
    \ifprofessionalPrint
        \newgeometry
        {
            textwidth = 134 mm,
            textheight = 220 mm,
            top = 38 mm + \extraborderlength,
            inner = 38 mm + \mybindingcorrection + \extraborderlength,
        }
    \else
        \newgeometry
        {
            textwidth = 134 mm,
            textheight = 220 mm,
            top = 38 mm,
            inner = 38 mm + \mybindingcorrection,
        }
    \fi
\else
    \newgeometry
    {
        textwidth = 134 mm,
        textheight = 220 mm,
        top = 38 mm,
        inner = 38 mm,
    }
\fi

% The format of the title page.
\begin{titlepage}
    \sffamily
    \begin{center}
        \includegraphics[height = 3.2 cm]{core/title_page/logo_UP.pdf} \hfill \includegraphics[height = 3 cm]{core/title_page/logo_HPI.pdf}\\
        \vfil
        {\LARGE
            \rule[1 ex]{\textwidth}{1.5 pt}
            \onehalfspacing\printTitleBold\\[1 ex]
            {\vspace*{-1 ex}\Large \printGermanTitle}\\
            \rule[-1 ex]{\textwidth}{1.5 pt}
        }
        \vfil
        {\Large\textbf{\printAuthor}}
        \vfil
        {\large Universitäts\colloquialDegreeNameLowercase arbeit\\[0.25 ex]
        zur Erlangung des akademischen Grades}\\[0.25 ex]
        \bigskip
        {\Large \colloquialDegreeName{} of Science}\\[0.5 ex]
        {\large\emph{(\degreeAbbreviation\,Sc.)}}\\
        \bigskip
        {\large im Studiengang\\[0.25 ex]
        \printProgram}
        \vfil
        {\large eingereicht am \printDateReceived{}\\[0.25 ex]
        Fachgebiet Systemanalyse und Modellierung der\\[0.25 ex]
        Digital-Engineering-Fakultät\\[0.25 ex]
        der Universität Potsdam}
    \end{center}
    
    \vfil
    \begin{table}[h]
        \centering
        \large
        \sffamily 
        {\def\arraystretch{1.2}
            \begin{tabular}{>{\bfseries}p{3.8 cm}p{6.4 cm}}
                Erstgutachter               & \printNameOfSupervisor\\
                Zweitgutachter               & \printNameOfSecondSupervisor\\
                Betreuer  & \printAdditionalExaminers
            \end{tabular}
        }
    \end{table}
\end{titlepage}

\restoregeometry

    \pagestyle{plain}

    \addchap{Abstract}
    Code review is limited by the evidence visible to the reviewer. When an edit
breaks code in another file, neither the diff nor the reviewer's memory
guarantees that the consequence will be found. This thesis asks how repository
context changes the coverage of LLM-generated reviews, whether structural
context enables cross-file reasoning, and whether human raters perceive the
resulting difference.

A pipeline generates reviews under four modes: a diff-only baseline, a
structural arm that supplies edges from a Joern code property graph, a
retrieval arm that supplies embedding-similar code, and a hybrid containing
both. The modes are evaluated on 35 pull requests from five repositories, on
40 controlled defect injections, and in a human study.

On the pull requests, a panel of five language-model judges scores each review
on 25 criteria. The graph arm gains $+0.69$ on the nine criteria designated
in advance as structurally relevant, and 96\,\% of its total advantage falls
within them (both $p = 0.003$). Retrieval produces the largest
full-rubric gain ($+1.03$, $p = 0.002$). After correction for multiple
comparisons, the graph retains its targeted gain and retrieval its full-rubric
gain; the hybrid retains neither.

Among 28 injected defects whose consequences lie outside the edited file, the
baseline detects none, retrieval detects 5, and the graph detects 18. Adding
inheritance edges raises detection to 26 of 28. On 12 control defects
visible in the edited file, every mode
performs at or near ceiling. The eight additional detections from inheritance
edges show that graph construction has substantial room to improve.

In the human study, 27 raters compare graph and baseline reviews.
They prefer the graph review for naming affected code (mean preference
$0.710$, where $0.5$ denotes no preference; $p = 0.0005$), show no
preference on overall usefulness, and prefer the baseline review's
commentary on code clarity.

The contribution is a bounded capability, not a general improvement in review
quality. Structural context exposes cross-file consequences that a diff does
not contain; similarity retrieval improves broader rubric coverage but does
not replace structural evidence. The benefit is large when cross-file
reasoning is required and smaller but targeted on the pull-request
benchmark.


    \selectlanguage{ngerman}
    \addchap{Zusammenfassung}
    % Deutsche Zusammenfassung; inhaltlich synchron mit preface/abstract.tex.

Code-Reviews sind auf die verfügbaren Informationen beschränkt. Verursacht
eine Änderung einen Fehler in einer anderen Datei, lässt sich diese Folge weder
aus dem Diff noch zuverlässig aus dem Gedächtnis ableiten. Diese Arbeit
untersucht, wie Repository-Kontext die inhaltliche Abdeckung LLM-generierter
Reviews verändert, ob struktureller Kontext dateiübergreifendes Schlussfolgern
ermöglicht und ob menschliche Bewertende den daraus resultierenden
Unterschied wahrnehmen.

Eine Pipeline erzeugt Reviews in vier Modi: eine reine Diff-Basislinie, ein
struktureller Arm, der Kanten aus einem mit Joern erstellten Code Property Graph
liefert, ein Retrieval-Arm, der embedding-ähnlichen Code liefert, und eine
Kombination beider Kontexte. Die Modi werden auf 35~Pull Requests aus fünf
Repositories, auf 40~kontrollierte Defektinjektionen und in einer
Nutzerstudie evaluiert.

Auf den Pull Requests bewertete ein Panel aus fünf Sprachmodellen jedes Review
anhand von 25~Kriterien. Der Graphmodus erzielt einen Zugewinn von
$+0{,}69$~Punkten bei den neun im Voraus als strukturell relevant festgelegten
Kriterien, und 96\,\% seines gesamten Zugewinns entfallen auf diese
Kriterien (jeweils $p = 0{,}003$). Retrieval erzielt den größten Zugewinn auf der
vollständigen 25-Kriterien-Skala ($+1{,}03$, $p = 0{,}002$). Nach der Korrektur
für multiple Vergleiche bleibt beim Graphmodus der gezielte Effekt und bei
Retrieval der Effekt auf der Gesamtskala signifikant; beim Hybridmodus bleibt
keiner der beiden Effekte signifikant.

Bei 28~injizierten Defekten, deren Folgen außerhalb der editierten Datei
liegen, erkennt die Basislinie keinen Defekt, Retrieval erkennt 5 und der
Graph~18. Das Hinzufügen von Vererbungskanten erhöht die Erkennung auf 26 von
28. Bei
12~Kontrolldefekten, die in der editierten Datei sichtbar sind, erreichen alle
Modi eine vollständige oder nahezu vollständige Erkennung. Die acht zusätzlich
durch Vererbungskanten erkannten Fälle zeigen erhebliches Verbesserungspotenzial
in der Graphkonstruktion.

In der Nutzerstudie verglichen 27~Teilnehmende Reviews mit Graph-Kontext
und Reviews der Basislinie. Beim Benennen betroffener Codestellen
bevorzugten sie das Review mit Graph-Kontext (mittlerer Präferenzwert
$0{,}710$; $0{,}5$ bezeichnet keine Präferenz; $p = 0{,}0005$). Bei der
wahrgenommenen Gesamtnützlichkeit zeigte sich keine Präferenz; bei den
Anmerkungen zur Codeklarheit bevorzugten sie das Review der Basislinie.

Der Beitrag besteht in einem klar abgegrenzten Fähigkeitsgewinn und nicht in
einer allgemeinen Verbesserung der Reviewqualität. Struktureller Kontext macht
dateiübergreifende Folgen sichtbar, die im Diff nicht enthalten sind.
Ähnlichkeitsbasiertes Retrieval erhöht dagegen die Abdeckung des gesamten
Kriterienrasters, ersetzt strukturelle Informationen jedoch nicht. Der Nutzen
ist groß, wenn dateiübergreifendes Schlussfolgern erforderlich ist. Im
Pull-Request-Benchmark fällt er kleiner, aber gezielt aus.

    \selectlanguage{american}

    \addchap{Acknowledgments}
    I am grateful to Christian Medeiros Adriano for his supervision throughout
this thesis. His feedback on the experimental design, the evaluation, and the
writing shaped the work at every stage.

I would like to thank my first supervisor, Prof.\ Dr.\ Holger Giese, for
making this project possible and for his guidance along the way. I also
thank Prof.\ Dr.\ Walid Maalej for agreeing to serve as second supervisor
and examiner of this thesis.

I also thank the twenty-seven participants who volunteered their time for
the human evaluation study.


    \setuptoc{toc}{totoc}
    \tableofcontents

    \pagestyle{headings}
    \mainmatter

    \chapter{Introduction}\label{chap:introduction}
    \section{Motivation}
\label{sec:intro-motivation}

Consider a change that narrows what a widely used helper function will accept:
where it previously tolerated a missing value, it now requires one. Inside the
diff the change is sound. The function's own tests are updated, and a reviewer
reading those few lines has every reason to approve. The defect, if there is
one, is somewhere else: in a caller three directories away that has always
passed nothing and relied on the old tolerance. Nothing in the diff mentions
that caller.

This scenario illustrates a recognised challenge in code review; its frequency
is not estimated in this thesis. When Bacchelli and Bird asked Microsoft
developers what they hoped to get from review and what they actually got, the
gap was defect detection, and the
reason they gave was not inattention but
\emph{understanding}~\cite{bacchelli2013}: a reviewer who does not know what
depends on a piece of code cannot judge whether changing it is safe.

A language model asked to review a pull request is usually given the diff and
nothing else. It cannot open the neighbouring file, search for usages, or run
the tests. The natural response is to put more of the repository into the
prompt, and two mechanisms are available. Retrieval embeds the repository and
returns the code most similar to the change. Structural traversal parses the
repository and returns the code \emph{connected} to the change. In the
vignette above, similarity would rank the helper's sibling utilities highly,
while the caller that is about to break may share no vocabulary with the
function it calls. Which mechanism to use, and what each is actually good for,
is the question this thesis takes up.

\section{Problem Statement}
\label{sec:intro-problem}

The problem is not that repository context is unavailable. Both retrieval and
graph construction are mature enough to deploy, and systems built on them have
demonstrated gains on completion, fault localisation and
repair~\cite{zhang2023repocoder,liu2024graphcoder,yang2025kgcompass}. The
problem is that for review-comment generation it is not known which kind of
context helps, by how much, or whether the instruments commonly used to answer
such questions can detect the answer.

Three difficulties compound. First, the tasks on which structural context has
been shown to help all have checkable outcomes (completion is right or
wrong, a test passes or fails), while a review comment is prose with no
automatic success signal, making evaluation part of the problem. Second, the
evaluation methods now used for this task are weaker than the claims made with
them: single-model rubric scoring with unreported reliability cannot
distinguish a real improvement from a preference for longer
text~\cite{zheng2023judging}. Third, the structural context must come from
a described static-analysis artefact, here a code property graph, rather
than an ad hoc text heuristic, so that the result can be attributed and
replicated.

\section{Research Questions}
\label{sec:intro-rqs}

Three research questions form a chain, each closing a gap the previous one
leaves open:

\begin{description}
    \item[RQ1: Comparative coverage.] \textit{How do \kg{}, \rag{}, and
        \hybrid{} context affect the rubric coverage of LLM-generated review
        comments compared with a diff-only \baseline{}?}
    \item[RQ2: Cross-file reasoning.] \textit{To what extent does
        structural context enable an LLM reviewer to identify cross-file
        consequences that are absent from the diff?}
    \item[RQ3: Human preference.] \textit{Do human raters prefer
        graph-augmented reviews over diff-only reviews for identifying affected
        code when the review conditions are hidden?}
\end{description}

RQ1 is evaluated with a 25-criterion rubric, including nine criteria
designated in advance as structurally relevant (Section~\ref{sec:rubric}).
The rubric measures what a review covers, not whether it is correct, so
RQ2 uses controlled defect injection, where the ground truth is known.
Both are scored by language models, so RQ3 asks people. RQ1 and RQ2 build the
knowledge graph with Joern; RQ3 uses the same graph block with edges
resolved from import statements. The hybrid arm tests whether structural and similarity context
combine their benefits or dilute attention when concatenated.

\section{Contributions}
\label{sec:intro-contributions}

\paragraph{Mechanistic evidence for the structural blind spot.}
A controlled injection experiment shows that a diff-only reviewer detects none
of 28 cross-file defects, while the graph arm detects 18. Near-ceiling
performance on local controls supports a structural interpretation of the gain.

\paragraph{A localised benefit on pull requests.}
On 35~pull requests, the graph arm concentrates its gain in
the nine criteria its mechanism predicts, while retrieval produces the largest
full-rubric gain. The hybrid mode leads on neither scale after correction.

\paragraph{Human corroboration of targeted grounding.}
In the human study, raters prefer the graph-augmented review for naming
affected code, but show no preference on overall usefulness; code clarity
favours the baseline.

\paragraph{An analysis of evaluation sensitivity.}
Judge-level and regeneration analyses show that the targeted direction is
stable, but its magnitude and significance depend on the evaluation
configuration.

\medskip\noindent
The remainder of this thesis is organised as follows.
Chapter~\ref{chap:background} situates the work in the literature on code
review, language models for code, retrieval, knowledge graphs, and evaluation.
Chapter~\ref{chap:approach} describes the pipeline and the graph
construction. Chapter~\ref{chap:edas} presents the experimental design, the
rubric, and the human study protocol. Chapter~\ref{chap:results} reports and
discusses results from all three studies together with threats to validity.
Chapter~\ref{chap:conclusion} summarises the findings and identifies future
work.



    %%%%%%%%%%%%%%%%%%%%%%%%%%%%%%%%%%%%%%%%%%%%%%%%%
    %% Please add the content of your thesis here. %%
    %%%%%%%%%%%%%%%%%%%%%%%%%%%%%%%%%%%%%%%%%%%%%%%%%
    
    \chapter{Background and Related Work}\label{chap:background}
    This chapter introduces the concepts and prior work that the rest of the thesis
builds on: code review and its automation, language models for code, the two
strategies for selecting repository context, how generated reviews are
evaluated, and the statistical tools the analysis uses. Each section closes by
identifying the gap the prior work leaves open, and the final section collects
those gaps into the position this thesis occupies.

\section{Modern Code Review and Its Automation}
\label{sec:bg-code-review}

Modern code review is the lightweight, asynchronous inspection of a proposed
change before it enters a shared branch, organised around pull requests on
hosted platforms~\cite{davila2021mcrslr}. Bacchelli and Bird found that
developers enter review expecting to find defects but most reliably obtain
knowledge transfer and awareness of alternative solutions; the gap they
identified is \emph{understanding}, meaning a reviewer who does not know what
depends on a piece of code cannot tell whether changing it is
safe~\cite{bacchelli2013}. Subsequent work confirms the pattern: useful comments
identify functional defects and consequences the author missed, while comments
about naming and formatting are judged unhelpful~\cite{bosu2015}; at scale,
review is fast and habitual, calibrated to maintain norms as much as to catch
bugs~\cite{sadowski2018}. These studies, together with the ISO/IEC~25010
quality model~\cite{iso25010}, supply the categories from which the evaluation
rubric in Chapter~\ref{chap:edas} is derived.

Automating this task has been framed as transduction from a diff to a comment.
CodeReviewer pre-trains on review data and fine-tunes for comment
generation~\cite{li2022codereviewer}; Tufano et al.\ show pre-trained models
outperform earlier baselines~\cite{tufano2022codereview}; general-purpose
instruction-following models subsequently displaced task-specific
fine-tuning~\cite{hou2023llmse,fan2023llmse}.

The same shift enabled a wave of commercial AI review tools. CodeRabbit's
code-graph analysis resolves cross-file dependencies and symbol definitions
to enrich each review, alongside the output of more than forty integrated
linters and static analysers~\cite{coderabbit2026}. Greptile builds a graph of the
repository's files, functions, and dependencies and reviews each change
against it~\cite{greptile2026}. Qodo~Merge (built on the open-source
PR-Agent) deploys specialised review agents that gather context from the
whole codebase, including dependent repositories~\cite{qodo2026}. GitHub's Copilot code-review feature was used by
over one million developers within a month of its public preview and reached
general availability in April 2025, while operating primarily on the
diff~\cite{copilotcr2025}. These systems demonstrate that
practitioners value repository-wide context, but none exposes the reasoning
to controlled evaluation: they are closed-source products, not matched
experimental arms.

Two industrial studies begin to close that gap. Cihan et al.\ deploy Qodo's
PR-Agent at Beko on 1,568 of 4,335 pull requests and find that 73.8\% of
automated comments are resolved, although average closure time
increases~\cite{cihan2025automated}. Sun et al.\ report 75\% precision and
12\,000 weekly users for ByteDance's BitsAI-CR
pipeline~\cite{sun2025bitsaicr}.

In every case the input is the change, and everything the model knows about
the rest of the repository comes from pre-training or from the tool's own
retrieval pipeline. Recent benchmarks treat that as the limiting factor.
ContextCRBench tests review with enriched
context~\cite{hu2025contextcrbench}. SWR-Bench assembles a thousand verified
pull requests~\cite{zeng2025swrbench}. CodeFuse-CR-Bench evaluates at
repository level~\cite{guo2025codefusecrbench}. MCR-Bench benchmarks
multi-round review across 2,269 tasks and finds that 23\% of missed defects
are missed because the consequence lies in a different
file~\cite{zheng2026mcrbench}. The cross-file gap that MCR-Bench quantifies
is the blind spot targeted with a knowledge graph in this thesis.

While these benchmarks measure the effect of context, a parallel line of work
combines static-analysis output directly with LLM review. Jaoua et al.\
integrate PMD and Checkstyle diagnostics with a fine-tuned CodeLlama via
three strategies and find that the hybrid improves relevance and
completeness~\cite{jaoua2025combining}. In vulnerability detection, LLMxCPG
uses Joern to extract code-property-graph slices that reduce code volume by
67--91\% while preserving vulnerability-relevant context, improving F1 by
15--40\% over prior baselines~\cite{lekssays2025llmxcpg}. These results show that
static-analysis artefacts can usefully condition an LLM, but neither applies
CPG-level structural context to pull-request review. Jaoua et al.\ supply
linter diagnostics (rule violations), not a structural graph; LLMxCPG targets
security auditing, not review-comment generation.

All of these benchmarks and combination strategies establish that context
matters but do not compare structural and similarity-based selection within
one matched review pipeline, which is the comparison this thesis runs.

\section{Large Language Models for Code}
\label{sec:bg-llms}

Every generator and judge in this thesis is a transformer~\cite{vaswani2017attention}.
Two properties of the architecture matter here: salience is learned, so a fact
placed in a long prompt competes with everything else, and attention costs grow
with sequence length, making the question of \emph{what} to put in the context
window unavoidable.

Applying transformers to code began with code-specific pre-training.
CodeBERT learns joint representations of natural language and programming
language~\cite{feng2020codebert}; GraphCodeBERT adds a data-flow objective,
and reports gains on downstream code tasks from incorporating that
structure~\cite{guo2020graphcodebert}. General-purpose
instruction-following models such as GPT-4~\cite{openai2023gpt4} subsequently
displaced fine-tuning for many software-engineering
tasks~\cite{hou2023llmse,fan2023llmse}.

Two operational properties bear on the experiments. Decoding is stochastic:
setting temperature to zero requests greedy decoding, which is nearly but not
exactly reproducible. And a model's input is bounded, so selecting a subset of a
repository is a requirement, not an optimisation. The two established ways of
making that selection, similarity-based retrieval and structural traversal,
are the subjects of the next two sections.

\section{Retrieval-Augmented Generation}
\label{sec:bg-rag}

Retrieval-augmented generation supplies a model with material fetched at
inference time rather than relying on pre-training
alone~\cite{lewis2020rag}. A retriever selects passages from a corpus and the
generator conditions on them together with the query. For software the appeal
is immediate: a repository at a particular commit is knowledge no pre-trained
model can have.

Modern retrieval is typically dense. Karpukhin et al.\ showed that a dual
encoder mapping questions and passages into a shared embedding space
outperforms sparse baselines on open-domain question
answering~\cite{karpukhin2020dpr}. Three parameters recur throughout
Chapter~\ref{chap:approach}: how the corpus is chunked, how the query is
formed, and how many chunks are kept.

The property that matters for code review is what similarity does and does not
guarantee. Similarity in embedding space is a claim about resemblance, not
about consequence. A function that resembles the changed code will rank highly
whether or not the change affects it; the caller that will actually break may
share no surface features with the callee and rank poorly. Retrieval optimises
for looking related; review needs what \emph{is} related.

Graph-structured retrieval is one response.
GraphRAG~\cite{edge2024graphrag} builds an entity graph over the corpus to
answer questions about how documents connect rather than about any single
document. RepoCoder~\cite{zhang2023repocoder} alternates retrieval and
generation for repository-level completion, on the finding that a single
similarity query over a diff retrieves poorly. Both are admissions about the
limits of similarity as a relevance criterion and both point at the comparison
this thesis makes: if similarity retrieval were adequate, neither iterative
re-querying nor graph construction would be necessary.

The retrieval arm in this thesis is a deliberately conventional instance of the
paradigm (chunked corpus, cosine similarity, fixed top-$k$), included not
as a contribution but as the comparison that makes the graph arm
interpretable.

\section{Knowledge Graphs for Software Repositories}
\label{sec:bg-kg}

A knowledge graph represents entities as nodes and relations between them as
typed edges. Applied to a repository, the entities are files, classes,
functions, tests, and the edges are relations such as imports, calls, and
tests. These are the relations a reviewer consults when deciding whether a
change is safe, and unlike similarity they are facts about the program rather
than resemblances. In this thesis, the structural context supplied to the
review model is a per-pull-request subgraph extracted from a code property
graph (CPG) for the files touched by the change; the review mode that
receives it is labelled \kg{} throughout.

\subsection{Program Representations}

The edges come from classical program representations. Parsing source into an
abstract syntax tree recovers grammatical structure~\cite{aho2006compilers}; a
control-flow graph adds execution order; a program-dependence graph adds
data-flow. Yamaguchi et al.\ combined all three into the code property graph
(CPG), a single structure over which queries can express vulnerability patterns
as graph traversals~\cite{yamaguchi2014cpg}. The CPG is the representation
behind the static-analysis platform used in Chapter~\ref{chap:approach}.
Allamanis et al.\ showed that learning over program graphs rather than token
sequences improves tasks requiring reasoning about variable
use~\cite{allamanis2018programgraphs}. This motivates an analogous hypothesis
for prompting: a relationship distant in the token stream can instead be
stated directly.

\subsection{Repository-Scale Graphs for Language Models}

Three recent systems build graphs specifically to serve language models, each
targeting a different task. GraphCoder retrieves context for code completion by
traversing a code-context graph rather than by embedding
similarity~\cite{liu2024graphcoder}. KGCompass links repository artefacts to
code entities for fault localisation, reporting that most bugs it finds are
reachable only by traversing more than one
edge~\cite{yang2025kgcompass}. Prometheus maintains a repository graph as
persistent agent memory~\cite{pan2025prometheus}. All target tasks with
checkable outcomes (completion, repair, navigation). Review-comment
generation has no such outcome, which is both why the graph literature has not
addressed it and why it is worth addressing.

All three build graphs by static analysis, so their edges are resolved rather
than guessed. This thesis adopts the same principle: the code property graph
constructed by Joern~\cite{yamaguchi2014cpg} supplies the structural context
used in Chapters~\ref{chap:approach}--\ref{chap:results}. Because the CPG
resolves calls and data-flow edges across files, the context it provides to the
language model is grounded in program semantics rather than in surface-level
co-occurrence. The human study uses the same knowledge-graph block, with dependent
edges resolved from import statements; Chapter~\ref{chap:approach} gives
the details.

\section{Evaluating Generated Reviews}
\label{sec:bg-evaluation}

How to score a generated review is contested. Text-overlap metrics against a
single human reference have been shown to misrank review comments, because a
change admits many valid reviews~\cite{hong2024deepcrceval}. Criterion-based
rubrics scored by a language model replaced them: Liu et al.'s G-Eval
established the pattern, reporting correlation with human ratings markedly
better than $n$-gram metrics~\cite{liu2023geval}. Zheng et al.\ generalised it
with MT-Bench, finding agreement between a strong judge and human preferences
comparable to agreement between two humans~\cite{zheng2023judging}. He et al.\
survey adoption of the paradigm in software
engineering~\cite{he2025llmasjudgese}.

The same literature catalogues failure modes. Zheng et al.\ document position
bias, verbosity bias, and self-enhancement bias~\cite{zheng2023judging};
Panickssery et al.\ show that models recognise and favour their own
generations~\cite{panickssery2024selfpreference}. Verbosity bias is
particularly relevant here: adding context to a prompt tends to make the output
longer, so a judge that rewards length would reward context augmentation for the
wrong reason.

The mitigations this thesis adopts follow directly. Using several judges from
different providers limits any one family's idiosyncrasies. Recomputing the
effect under individual judges makes evaluator sensitivity visible. Blinding
and pairwise presentation control position and identity
effects~\cite{chiang2024chatbotarena}. Where objective ground truth can be
constructed, adjudication is checked against it rather than trusted.

One conceptual limitation applies to every rubric of this kind. Asking whether
a review addresses a criterion measures \emph{coverage}: whether the review
says something about the subject. It does not measure whether what it says is
\emph{correct}. This thesis uses a coverage rubric for its main comparison and
supplements it with both an oracle-scored experiment and a human study,
precisely because coverage alone cannot answer whether added context was used
correctly.

The evaluation methods specific to code review share this gap. DeepCRCEval
replaces overlap metrics with a criterion-based rubric scored by a single
judge~\cite{hong2024deepcrceval}; CRScore constructs pseudo-references from
static analysers~\cite{naik2025crscore}. Neither reports agreement between
independent scorers, because neither has two. This matters more for a
context-augmentation study than for a model ranking, because verbosity bias
from the treatment cannot be separated from a real improvement without a panel.

\section{Statistical Methods}
\label{sec:bg-stats}

The analyses in Chapter~\ref{chap:results} use a small set of standard
procedures, collected here so that the results chapter can apply them without
digression. Observations are \emph{paired}: every arm is evaluated on the
same pull requests, and scores are bounded counts rather than draws from a
normal distribution.

\paragraph{Agreement.}
Cohen's $\kappa$ compares observed agreement between raters with the agreement
expected by chance, giving $1$ for perfect and $0$ for chance-level
agreement~\cite{cohen1960kappa}. Values are read against the bands of Landis
and Koch ($0.41$--$0.60$ moderate, $0.61$--$0.80$ substantial, above $0.80$
almost perfect)~\cite{landiskoch1977}.

\paragraph{Confidence intervals.}
Because the score distributions are neither normal nor of a known form,
intervals are obtained by bootstrapping: the sample is resampled with
replacement, the statistic recomputed on each resample, and the interval read
off the percentiles~\cite{efron1993bootstrap}.

\paragraph{Testing a paired difference.}
A paired permutation test flips the sign of each difference at random and
recomputes the mean; the $p$-value is the proportion of permutations producing
a mean at least as extreme as observed. Where an ordinal reading is
appropriate, the Wilcoxon signed-rank test is used
instead~\cite{wilcoxon1945}.

\paragraph{Effect size.}
For paired data, $d_z$ is the mean difference divided by the standard
deviation of the differences (small at $0.2$, medium at $0.5$, large at
$0.8$). Both $p$-values and effect sizes are reported throughout, since with
thirty-five observations a real but small effect and a null result can be hard to
tell apart from the $p$-value alone.

\paragraph{Multiple comparisons.}
Holm's sequentially rejective procedure controls the family-wise error rate
with more power than Bonferroni by ordering $p$-values and comparing each
against a relaxing threshold~\cite{holm1979}. Chapter~\ref{chap:results}
states for each family which correction was applied and over how many
hypotheses.

\section{Positioning and Research Gap}
\label{sec:bg-positioning}

Table~\ref{tab:rw-positioning} places this thesis against the six most
directly comparable systems.

\begin{table}[htbp]
    \centering
    \caption{Positioning against the closest prior work. \textbf{DC} =
             DeepCRCEval~\cite{hong2024deepcrceval}, \textbf{CR} =
             CRScore~\cite{naik2025crscore}, \textbf{GC} =
             GraphCoder~\cite{liu2024graphcoder}, \textbf{KGC} =
             KGCompass~\cite{yang2025kgcompass}, \textbf{JA} =
             Jaoua et al.~\cite{jaoua2025combining}, \textbf{LxC} =
             LLMxCPG~\cite{lekssays2025llmxcpg}. ``n/a'' marks systems scored by
             program execution.}
    \label{tab:rw-positioning}
    \small
    \begin{tabularx}{\textwidth}{@{}Xccccccc@{}}
        \toprule
          & DC & CR & GC & KGC & JA & LxC & \textbf{This thesis} \\
        \midrule
        Task is review comments
          & $\bullet$ & $\bullet$ & --- & --- & $\bullet$ & --- & $\bullet$ \\
        Adds out-of-diff context
          & --- & --- & $\bullet$ & $\bullet$ & $\bullet$ & $\bullet$ & $\bullet$ \\
        Structural (graph) context
          & --- & --- & $\bullet$ & $\bullet$ & --- & $\bullet$ & $\bullet$ \\
        Static-analysis input
          & --- & $\bullet$ & --- & --- & $\bullet$ & $\bullet$ & $\bullet$ \\
        Retrieval and structure as matched arms
          & --- & --- & $\bullet$ & --- & $\bullet$ & --- & $\bullet$ \\
        Independent LLM scorers
          & 1 & 1 & n/a & n/a & --- & --- & \textbf{5}\footnotemark \\
        Inter-scorer agreement reported
          & --- & --- & n/a & n/a & --- & --- & $\bullet$ \\
        Detection adjudication checked against oracle
          & --- & --- & n/a & n/a & --- & n/a & $\bullet$ \\
        Human corroboration of targeted criterion
          & $\bullet$ & $\bullet$ & --- & --- & --- & --- & $\bullet$ \\
        \bottomrule
    \end{tabularx}
\end{table}
\footnotetext{Five-provider panel, aggregated by a strict majority of at least
four valid votes; pairwise Cohen's $\kappa$ is reported in
Section~\ref{sec:results-exp1-agreement}. Experiment~2 additionally uses a
deterministic file oracle.}

\subsection*{Research Gap}

No existing study compares structural and similarity-based out-of-diff context
for review-comment generation within one matched generation and evaluation
pipeline, while also measuring sensitivity to the evaluator. GraphCoder runs a
matched comparison on code completion with checkable outcomes, while
DeepCRCEval and CRScore evaluate review comments but do not vary context or
report scorer sensitivity. Two recent studies narrow different parts of the
gap: Jaoua et al.\ compare three ways of injecting linter output into an LLM
reviewer but do not supply structural graph context, and LLMxCPG feeds Joern
CPG slices to an LLM but targets vulnerability detection, not review-comment
generation.

This thesis occupies the remaining gap along three dimensions the table makes
visible. It compares selection strategies rather than proposing one, using
matched arms on the same pull requests (Chapter~\ref{chap:approach}). It
treats the evaluation instrument as an object of study, reporting how much the
conclusion moves across individual judges and generators. And it corroborates
a coverage rubric with a deterministic oracle and a human study.



    \chapter{Approach}\label{chap:approach}
    \section{Overview}
\label{sec:approach-overview}

The system takes a pull request and produces a written review. Three design
properties make the comparison in Chapter~\ref{chap:edas} a comparison of
context rather than of pipelines.

The pipeline splits into a construction stage and a generation stage, with a
frozen artefact between them. Construction reads the repository once, at the
pull request's head commit, and writes an \emph{evidence pack}: a JSON file
holding the change, the structural facts, and the retrieved code snippets
(Section~\ref{sec:approach-packs}). Generation reads that file and never
touches the repository. Because all four conditions read the same pack, none
can differ by accident of timing.

The conditions are nested. \baseline{} receives the change alone. \kg{} adds a
structural block, \rag{} adds a retrieval block, and \hybrid{} adds both. No
condition is deprived of the diff. The context arms add evidence and a sentence
that names its type, so the contrast includes both supplied information and
context-specific instruction.

Structural context is produced by a code property graph (CPG) built with
Joern, a static-analysis platform that parses source into an interlinked
representation of syntax, control flow, and data flow
(Section~\ref{sec:approach-kg}). Throughout the thesis, \kg{} and ``the
graph arm'' name this same condition. The human study uses the same
knowledge-graph block, with dependent edges resolved from Python import
statements (Section~\ref{sec:approach-kg-ast}). Figure~\ref{fig:pipeline}
shows the arrangement.

\begin{figure}[htbp]
    \centering
    \resizebox{0.92\textwidth}{!}{%
    % !TeX root = ../thesis.tex
% Pipeline figure — clean academic style (thin lines, 2 accent colours).
\begin{tikzpicture}[
    >=Latex,
    every node/.style={font=\footnotesize, inner sep=0pt},
    box/.style={
      draw=black!50, line width=0.5pt, rounded corners=2pt,
      fill=white, minimum height=7mm, text centered,
      inner sep=4pt, align=center
    },
    kgbox/.style={box, draw=kgblue!70, fill=kgblue!6},
    ragbox/.style={box, draw=ragorange!70, fill=ragorange!6},
    accentbox/.style={box, line width=0.8pt},
    arr/.style={->, line width=0.5pt, black!50},
    fatarr/.style={->, line width=0.7pt, black!55}
]

% ── Row 0: Repository ─────────────────────────────────
\node[box, minimum width=6cm] (repo)
    at (5.5, 0) {Repository at PR head commit};

% ── Row 1: Three extractors ───────────────────────────
\node[box, minimum width=2.4cm] (diff) at (1.2,-1.4) {Unified diff};
\node[kgbox, minimum width=2.8cm] (cpg)  at (5.5,-1.4) {Joern CPG};
\node[ragbox, minimum width=2.8cm] (emb)  at (9.8,-1.4) {Embedding index};

\draw[arr] (repo) -- (diff);
\draw[arr] (repo) -- (cpg);
\draw[arr] (repo) -- (emb);

% ── Row 2: KG detail ─────────────────────────────────
\node[kgbox, minimum width=2.0cm, font=\scriptsize] (calls)
    at (3.7,-2.8) {Callers};
\node[kgbox, minimum width=2.0cm, font=\scriptsize] (deps)
    at (5.9,-2.8) {Dependents};
\node[box, minimum width=2.0cm, font=\scriptsize] (tests)
    at (8.1,-2.8) {Tests\\[-2pt]{\tiny filename match}};

\draw[arr] (cpg) -- (calls);
\draw[arr] (cpg) -- (deps);
\draw[arr] ([xshift=2.6cm]repo.south) -- (tests.north);

% ── Row 2: RAG detail ─────────────────────────────────
\node[ragbox, minimum width=2.4cm, font=\scriptsize] (chunks)
    at (10.5,-2.8) {Similar chunks};
\draw[arr] (emb) -- (chunks);

% ── Row 3: Evidence pack ──────────────────────────────
\node[accentbox, minimum width=10.5cm, minimum height=9mm,
      draw=black!60, fill=black!3] (pack)
    at (5.5,-4.2) {\textbf{Evidence pack}\;%
    {\scriptsize (one JSON file per PR)}};

\draw[arr] (diff.south) -- (diff.south |- pack.north);
\draw[arr] (calls.south) -- (calls.south |- pack.north);
\draw[arr] (deps.south) -- (deps.south |- pack.north);
\draw[arr] (tests.south) -- (tests.south |- pack.north);
\draw[arr] (chunks.south) -- (chunks.south |- pack.north);

% ── Annotation: "generated once, read by all modes" ───
\node[font=\scriptsize\itshape, text=black!45, anchor=west]
    at (11.0,-4.2) {generated once};

% ── Row 4: Four modes ─────────────────────────────────
\node[box, minimum width=2.1cm, fill=black!3] (bl)
    at (1.85,-5.8) {\textbf{Baseline}\\[-1pt]{\scriptsize diff}};
\node[kgbox, minimum width=2.1cm] (kgm)
    at (4.55,-5.8) {\textbf{KG}\\[-1pt]{\scriptsize diff\,+\,graph}};
\node[ragbox, minimum width=2.1cm] (ragm)
    at (7.25,-5.8) {\textbf{RAG}\\[-1pt]{\scriptsize diff\,+\,chunks}};
\node[box, minimum width=2.1cm, draw=black!50,
      fill=black!2] (hyb)
    at (9.95,-5.8) {\textbf{Hybrid}\\[-1pt]{\scriptsize diff\,+\,both}};
% Hybrid dual-colour accent
\draw[kgblue!50, line width=0.6pt, rounded corners=2.5pt]
    (8.84,-5.28) rectangle (11.06,-6.32);
\draw[ragorange!50, line width=0.6pt, rounded corners=3pt]
    (8.78,-5.22) rectangle (11.12,-6.38);

\draw[arr] (pack.south) -- ++(0,-0.35) -| (bl.north);
\draw[arr] (pack.south) -- ++(0,-0.35) -| (kgm.north);
\draw[arr] (pack.south) -- ++(0,-0.35) -| (ragm.north);
\draw[arr] (pack.south) -- ++(0,-0.35) -| (hyb.north);

% ── Row 5: Generator ─────────────────────────────────
\node[box, minimum width=9.5cm] (gen)
    at (5.5,-7.2) {GPT-4o generator \;\;$\rightarrow$\;\;
    {\scriptsize $35 \times 4 = 140$ reviews}};

\draw[arr] (bl.south)  -- (bl.south |- gen.north);
\draw[arr] (kgm.south) -- (kgm.south |- gen.north);
\draw[arr] (ragm.south) -- (ragm.south |- gen.north);
\draw[arr] (hyb.south) -- (hyb.south |- gen.north);

% ── Row 6: Judge panel ────────────────────────────────
\node[box, minimum width=10.5cm, minimum height=9mm] (judge)
    at (5.5,-8.6)
    {5-model judge panel\;\;$\times$\;\;25 criteria\;\;$=$\;\;%
     {\scriptsize 17\,500 scored cells}};

\draw[fatarr] (gen) -- (judge);

% ── Bracket + judge names on right ────────────────────
\draw[black!30, line width=0.4pt]
    (11.0,-7.9) -- (11.2,-7.9) -- (11.2,-9.3) -- (11.0,-9.3);
\node[font=\scriptsize, text=black!45, anchor=west, align=left]
    at (11.4,-8.6)
    {GPT-4o\\[-1pt]Gemini\,2.5\,Flash\\[-1pt]Claude\,Sonnet\,4.5\\[-1pt]%
     DeepSeek\,v4\,Pro\\[-1pt]Grok\,4.6};

\end{tikzpicture}
%
    }
    \caption{End-to-end pipeline. Arrows denote data flow: each stage
             reads the output of the stage above it. The repository is
             processed once into a frozen evidence pack. Each mode reads the
             same pack and differs only in which fields are rendered into the
             prompt. Related tests are matched by filename, not resolved from
             the graph. A five-model judge panel scores every review on 25
             binary criteria. The human study uses the same pipeline with the
             graph's dependent edges resolved from import statements
             (Section~\ref{sec:approach-kg-ast}).}
    \label{fig:pipeline}
\end{figure}

\section{Evidence Packs}
\label{sec:approach-packs}

An evidence pack is the unit of input to generation. It contains the pull
request's metadata, the list of changed files, the unified diff, the
structural facts from the graph builder, and the retrieved code chunks.

\section{Knowledge Graph Construction}
\label{sec:approach-kg}

The structural context supplied to the reviewer is a subgraph of the
repository's code property graph (CPG), built with Joern, a
language-agnostic static-analysis platform~\cite{yamaguchi2014cpg}. The CPG
represents syntax, control flow, and data flow in a single queryable
structure; the pipeline queries it for the cross-file relationships that a
diff-only reviewer cannot see.

\subsection{The Code Property Graph}
\label{sec:approach-kg-cpg}

At the pull request's head commit the full repository is loaded into Joern,
which parses every source file into an in-memory code property graph. The
pipeline then queries the graph for five relationship types seeded from the
changed files:
%
\begin{enumerate}[leftmargin=*,itemsep=2pt,label=(\roman*)]
  \item incoming callers of changed functions, with caller name, file, and
        line number;
  \item outgoing callees;
  \item file-level dependents (files that import a changed module);
  \item related test files, identified by filename-convention matching; and
  \item code owners read from the repository's \texttt{CODEOWNERS} file when
        present.
\end{enumerate}
%
Queries use exact name matching against the method and type nodes that Joern
extracts: a changed function's fully qualified name is looked up in the
call-graph overlay, and every node whose call target matches is returned.

Same-file and duplicate edges are removed because the diff already exposes
same-file context; the renderer presents up to twelve edges, ordered by
query type, without learned relevance ranking.
% Sources: experiments/2026-09-19_joern_replace_grep/,
% prnote/joern_kg.py, prnote/kg.py.

Experiment~2 additionally decomposes the graph arm into component
conditions: one containing only file-level dependency pointers, one
containing only the resolved call-graph edges, and one that adds
independently resolved inheritance relations. This design isolates the
contribution of call-graph edges from that of coarser pointers and measures
the gain from inheritance edges that the base Joern queries omit.
% Sources: experiments/2026-07-05_injection_exp2/harness/run_modes.py.

\subsection{The Knowledge-Graph Block in the Human Study}
\label{sec:approach-kg-ast}

The human study uses the same knowledge-graph context block as
Experiment~1: the same relations (changed files, dependents, related tests)
in the same rendered format. For its Python stimuli the dependent edges are
resolved from \texttt{import} statements rather than from the Joern CPG:
each candidate file is parsed with Python's standard library, its
\texttt{import} and \texttt{from} statements collected, and an edge
recorded only if a module name matches the changed file's module path. These edges are correct but file-level, and a
module-level import does not establish that a change to one function inside
that module reaches the importer.

\subsection{Rendering}
\label{sec:approach-kg-render}

Regardless of the builder, the pack's structural fields are rendered into a
Markdown block headed \emph{Knowledge Graph Context}, with sections for
changed files, related tests, dependents, code owners, functions, and callers.
The dependents, callers, and functions come from the CPG queries described
above or, in the human study, from import statements. Related tests are identified by filename-convention
matching (files whose path contains \texttt{test} and that share a directory
ancestor with a changed file), and code owners are read from the repository's
\texttt{CODEOWNERS} file when one exists. These two fields are therefore
text-matched, not statically resolved. Both constructions share the same
renderer and evidence-pack schema, so they differ in what they put in, not
how it is presented to the model.

Figure~\ref{fig:kg-block-example} shows a truncated graph block as the model
receives it.

\begin{figure}[htbp]
\centering
\begin{minipage}{0.92\textwidth}
{\footnotesize
\begin{verbatim}
**Knowledge Graph Context:**
**Changed Files:** .../Drawer.tsx, .../Drawer.story.tsx,
  .../SaveDashboardDrawer.tsx, .../InspectContent.tsx  [7 files]
**Related Tests:** .../RolePickerDrawer.test.tsx,
  .../SaveDashboardDrawer.test.tsx,
  .../HelpWizard.test.tsx                              [8 files]
**Files that depend on changes:**
  .../SaveDashboardDrawer.tsx,
  .../InspectContent.tsx                               [3 files]
**Functions in Changed Files:**
  SaveDashboardDrawerComponent (SaveDashboardDrawer.tsx);
  InspectContent (InspectContent.tsx);
  HelpWizard (HelpWizard.tsx)                         [15 entries]
**Call-graph edges (caller -> callee):**
  SaveDashboardDrawerComponent -> renderBody;
  InspectContent -> getErrors;
  InspectContent -> formatStats                        [4 edges]
\end{verbatim}
}
\end{minipage}
\caption{Truncated graph block from Grafana PR~\#96268 (Joern CPG). Paths
         are abbreviated; bracket counts show the full list lengths. The model
         receives this block between the diff and the closing instruction.}
\label{fig:kg-block-example}
\end{figure}

\section{Retrieval-Augmented Context}
\label{sec:approach-rag}

The retrieval arm is conventional dense retrieval, included as the comparison
that makes the graph arm interpretable. The repository is walked at the same
commit, skipping vendored directories and keeping recognised source
extensions; the first 200 candidate files are chunked and embedded with
\texttt{text-embedding-3-small}. Each changed file is embedded as a query,
and the five highest-scoring chunks from outside the changed files are
rendered into the prompt with their path, line range and a code excerpt.

\section{The Four Review-Generation Modes}
\label{sec:approach-modes}

A mode is a pair: a system prompt and a context builder.

\begin{table}[htbp]
    \centering
    \caption{The four generation modes. Every mode receives the title,
             description and diff; they differ in the context block appended.}
    \label{tab:modes}
    \small
    \begin{tabular}{@{}ll@{}}
        \toprule
        Mode & Context block \\
        \midrule
        \baseline{} & none \\
        \kg{}       & graph facts (dependents, callers, tests) \\
        \rag{}      & five retrieved snippets \\
        \hybrid{}   & both, concatenated \\
        \bottomrule
    \end{tabular}
\end{table}

The \hybrid{} mode is a concatenation, not a fusion. Its block is the graph
block followed by the retrieval block, with no joint ranking and no shared
budget. Its result therefore evaluates this simple concatenation, not all
possible ways of combining structural and similarity context.
Appendix~\ref{app:review-example-chapter} reproduces a side-by-side example
of a baseline and graph-augmented review for one pull request.

\section{Prompt Design}
\label{sec:approach-prompts}

All four modes share one system prompt, which casts the model as a senior
engineer and prescribes a five-heading output format. Each mode appends a
single paragraph naming the kind of context it has been given. The prompts are
reproduced in Appendix~\ref{app:prompts}. Two properties are load-bearing: the
fixed output format constrains what the evaluation can detect, and the per-arm
paragraph means the arms differ by one sentence of instruction as well as by
the block itself. The user message is assembled in a fixed order: title,
description, diff, context block, closing instruction.


    \chapter{Experimental Design and Setup}\label{chap:edas}
    % Chapter 5 -- Experimental Design and Setup
%
% This is the chapter to write FIRST: every input is frozen and none of it depends
% on the human study finishing.

\section{Experimental Design}

The evaluation rests on three studies that approach the same question from
complementary methodological directions.

\textbf{Experiment~1} is observational. Four review-generation modes are run
over 35 pull requests and scored against a 25-criterion rubric by a panel
of five LLM judges, with the majority verdict as the cell of record. A
language-model panel is used because the rubric requires 3\,500 binary
judgements per scorer (140 reviews $\times$ 25 criteria); a human gold
standard at that scale was not feasible, so human judgement is instead
reserved for the primary criterion in the third study. Experiment~1's
strength is direct evidence on pull-request reviews; its weakness is that it
establishes no causal mechanism and that the rubric measures coverage rather
than correctness.

\textbf{Experiment~2} is controlled. A defect is injected so that its
consequences lie in a different file from the edit; the outcome is whether the
review names a genuinely broken dependent site. Here the ground truth is known,
so a causal claim becomes available. The cost is that the defects are synthetic.

\textbf{The human study} asks people. Twenty-seven raters compare blinded
baseline and graph-augmented reviews on six pull requests, scoring
which review better identifies affected code and which is more useful overall.
Its strength is independent human judgement; its weakness is a smaller
stimulus set,
whose graph edges are resolved from import statements rather than from the
Joern CPG.

Each design covers a weakness of the others, although none removes them
all. The headline arms of Experiments~1 and~2 share their generator, judge
panel, and aggregation rule so that differences cannot be attributed to
those choices. The human study corroborates the panel's primary finding
against independent human judgement on a separate stimulus set; it does not
validate the panel itself.

\subsection{Research Questions}
\label{sec:rqs-addressed}

Experiment~1 addresses RQ1 by comparing rubric coverage across the four review
modes. Experiment~2 addresses RQ2 by testing detection of consequences outside
the diff; its component analyses explain which forms of structural evidence
support that capability. The human study addresses RQ3 through blinded
pairwise comparisons.

\subsection{Dataset}
\label{sec:dataset}

The evaluation uses 35 pull requests drawn from the pool curated by
Mariotto et~al.~\cite{mariotto2025esem}, sampling from five repositories:
\texttt{grafana/grafana}, \texttt{apache/kafka},
\texttt{scikit-learn/scikit-learn}, \texttt{godotengine/godot} and
\texttt{jenkinsci/jenkins}, spanning Go, TypeScript, Java, Scala, Python, and
C++.\footnote{Five Go-only pull requests from \texttt{grafana/grafana} are
excluded because Joern's Go frontend does not emit method-level nodes,
preventing construction of the call-graph edges required by the \kg{} arm.}
Table~\ref{tab:dataset-composition} gives the composition.

\begin{table}[htbp]
    \centering
    \caption{Composition of the evaluation dataset.}
    \label{tab:dataset-composition}
    \begin{tabular}{lrl}
        \toprule
        Repository & Total & Languages \\
        \midrule
        \texttt{grafana/grafana}             &  8 & Go, TS \\
        \texttt{apache/kafka}                &  8 & Java, Scala \\
        \texttt{scikit-learn/scikit-learn}   &  8 & Python \\
        \texttt{godotengine/godot}           &  6 & C++ \\
        \texttt{jenkinsci/jenkins}           &  5 & Java \\
        \midrule
        Total                                & 35 & \\
        \bottomrule
    \end{tabular}
\end{table}

Every candidate was screened against seven criteria, all properties of the pull
request rather than of any model output. A pull request is included only if it
was merged, is not a revert or backport, changes code (not only documentation),
touches only languages the parser supports, has a diff under 50\,000 characters, carries a
description of at least 100~characters, and does not have a purely cosmetic
title. This guarantees that all four modes, including \baseline{}, receive the
maintainer's own description.

At $n=35$, a two-sided paired permutation test at $\alpha = 0.05$ reaches
80\,\% power at
$d_z \approx 0.47$ and 70\,\% at $d_z \approx 0.42$,
bracketing the effects reported in Chapter~\ref{chap:results}.

\subsection{Rubric}
\label{sec:rubric}

Each review is scored against 25 binary criteria grouped into eight categories
(Table~\ref{tab:rubric-categories}). A criterion is satisfied if the review
addresses it; the score is the count of criteria satisfied.

\begin{table}[htbp]
    \centering
    \caption{The eight rubric categories. The KG-relevant subset comprises the
             nine criteria structural data can in principle inform.}
    \label{tab:rubric-categories}
    \begin{tabular}{llrl}
        \toprule
        Category & Criteria & $n$ & KG-relevant \\
        \midrule
        Functionality \& Integration & F1--F4 & 4 & F3, F4 \\
        Tests \& Verification        & T1--T3 & 3 & T1, T2, T3 \\
        Readability \& Structure     & R1--R3 & 3 & --- \\
        Maintainability \& Design    & M1--M3 & 3 & M1, M3 \\
        Consistency \& Style         & C1--C2 & 2 & C2 \\
        Performance                  & P1--P2 & 2 & --- \\
        Security \& Robustness       & S1--S3 & 3 & --- \\
        Review Quality (meta)        & Q1--Q5 & 5 & Q2 \\
        \midrule
        Total                        &        & 25 & 9 \\
        \bottomrule
    \end{tabular}
\end{table}

Every category is grounded in published sources: empirical studies of
review~\cite{bacchelli2013,bosu2015,sadowski2018}, evaluation frameworks for
generated comments~\cite{li2022codereviewer,hong2024deepcrceval}, and
ISO/IEC~25010~\cite{iso25010}.

\subsubsection{The KG-Relevant Subset}

Nine criteria are designated KG-relevant. The designation was fixed before
data collection and appears in the project's analysis plan committed to the
repository. A criterion qualifies if structural repository data can
substantively inform an answer to it. F4, for instance, asks whether the
review warns about impact on dependent code, which the caller fields can
answer, whereas R1 (naming and clarity) they cannot. This yields a testable
prediction: if the knowledge graph works as intended, its effect should
concentrate in these nine criteria. An independent re-annotation achieved
$\kappa = 0.615$ (84\,\% raw agreement).

\subsubsection{Rubric Discrimination}

To confirm the rubric distinguishes good from bad reviews, five deliberately
degraded review styles were scored on five pull requests. Four scored 5--9
points below the real \baseline{}. The fifth, a hallucinated review citing
fabricated files, scored 9.2 against the real \baseline{}'s 9.0. The rubric
measures whether a review \emph{addresses} a concern, not whether what it says
is \emph{true}. This is why every result is framed as coverage, not
correctness, and why Experiment~2's oracle is needed.

An operational note: seven of the 25 criteria sit at $0\,\%$ or $100\,\%$ in
both arms (Appendix, Table~\ref{tab:exp1-criteria}) and therefore cannot
register a difference in either direction, diluting the aggregate. This is
one reason the nine-criterion subscale is reported alongside the total.

\subsection{Experiment 2: Controlled Defect Injection}
\label{sec:injection-design}

Experiment~2 was designed to close the gap between coverage and correctness on
a narrow question: when a defect's consequences lie outside the edited file,
does context let the reviewer name them?

\subsubsection{Defect Taxonomy}

Injections are stratified into five structural bands (consequence in another
file) and two local control bands (consequence in the same file).

\begin{table}[htbp]
    \centering
    \caption{Injection bands. Structural bands require cross-file reasoning;
             control bands serve as negative controls, where no graph
             advantage is expected.}
    \label{tab:injection-bands}
    \small
    \begin{tabular}{llll}
        \toprule
        Band & Operator & Breaks at & Relationship \\
        \midrule
        S1 & Symbol rename          & Importers          & imports \\
        S2 & Signature change       & Callers            & calls \\
        S3 & Return-type change     & Consumers          & calls, data flow \\
        S4 & Base-class change      & Subclasses         & inheritance \\
        S5 & Import removal         & Importers          & imports \\
        \midrule
        L1 & Off-by-one             & Same file          & --- \\
        L2 & Removed null check     & Same file          & --- \\
        \bottomrule
    \end{tabular}
\end{table}

Band sizes are uneven (S1 9, S2 9, S3 1, S4 6, S5 3; six in each control
band), so S3 and S5 support no band-level claims.

The control bands carry the argument: a knowledge graph that improved detection
everywhere would be evidence of a context-volume effect rather than structural
reasoning. The prediction is asymmetric: advantage on structural bands,
parity on local. This is what makes the result falsifiable.

Forty injections were realised: 28 structural and 12 control, across three
of the Experiment~1 repositories: \texttt{scikit-learn/scikit-learn}
(Python), \texttt{apache/kafka} (Java) and \texttt{grafana/grafana}
(TypeScript). The four modes of Experiment~1 are carried over with the same
prompts, builder and renderer (Table~\ref{tab:config}). Two further arms are
included for the mechanism argument: an \emph{idealised} graph that is given
the oracle's own list of true dependents, and a graph with inheritance edges
added (\kg{} + inheritance).

Detection is binary and adjudicated in two stages. First, the review text is
searched for exact full-path or basename matches against the statically
resolved true dependents; if none is found, the verdict is deterministically
\emph{false} and no model call is made. When at least one true dependent is
mentioned, the review passages and matched evidence are sent to each judge,
which decides only whether the review links the mentioned file to the injected
change as broken or requiring an update. The majority of five judges is the
cell of record, with ties resolved to \emph{false}.

This two-stage procedure was finalised after an audit of the initial judge
prompt, which required judges to both locate filenames and assess causality in
a single step. Retained rationales exposed false negatives: reviews that
named a true dependent were sometimes judged as not detecting it. The
correction is post-hoc; the original judgments are preserved in the audit
trail and both sets of verdicts are available for recomputation.

\subsubsection{Component Analyses}

Two pre-registered follow-ups isolate the structural evidence supplied to the
reviewer. On the 28 structural injections, one ablation arm receives only the
dependent-file list (without edge annotations) and another only the resolved
call edges; both use the final five-judge adjudication. Two further
component probes, a related-test band and a context-relevance control, use
the earlier three-judge panel and are reported in
Appendix~\ref{app:supporting-components}.

\subsection{Statistical Analysis}
\label{sec:stats}

For Experiment~1's rubric scores, confidence intervals are percentile
bootstrap intervals ($B = 10\,000$).
Significance is assessed by paired sign-flip permutation tests
($B = 20\,000$), two-sided, reported alongside $d_z$. The pairing is by pull
request.

The governing RQ1 test is the pre-specified \kg{}--\baseline{} contrast on the
nine KG-relevant criteria. A rank-based multiplicity sensitivity and
per-criterion corrections are retained in
Appendix~\ref{app:supporting-multiplicity}, rather than treated as additional
headline tests.

For Experiment~2's binary detection outcomes, exact McNemar tests compare
paired arms and Wilson intervals replace bootstrap intervals (the appropriate
method for binomial proportions at small $n$).
Decision rules were fixed before the run: $+0.30$ detection
rate advantage at $p < 0.05$ for \kg{} on structural bands, parity on control
bands within $[-0.15, +0.15]$, and $+0.30$ for \kg{} + inheritance on the
structural bands.

\section{Experimental Setup}
\label{sec:setup}

\subsection{Models}

One model, gpt-4o, generates every review in Experiments~1 and~2 and in the
human study; the arms differ only in the context block and the sentence that
names it (Appendix~\ref{app:prompts}). All reviews are generated at
temperature~$0.0$, so temperature does not differ between experiments.
Absolute score levels remain not directly comparable because the samples
and outcomes differ.

One disclosure belongs here: no run log survives for the original review
generation, so the claim that Experiment~1 ran at temperature~$0.0$ rests on
the pipeline's default and recorded metadata of a later run with the same
configuration.

Scoring uses a panel of five judges (gpt-4o, gemini-2.5-flash,
claude-sonnet-4.5, deepseek-v4-pro and grok-4.6) with the majority verdict
as the cell of record; at least four valid votes are required and ties are
resolved to zero, the conservative direction for every treatment arm.
Table~\ref{tab:config} summarises the configuration.

\begin{table}[htbp]
    \centering
    \caption{Configuration of both experiments.}
    \label{tab:config}
    \small
    \begin{tabular}{lll}
        \toprule
                          & Experiment 1        & Experiment 2 \\
        \midrule
        Unit of analysis  & pull request & injected defect \\
        $n$               & 35                  & 40 (28 structural, 12 control) \\
        Repositories      & 5                   & 3 \\
        Generator         & gpt-4o              & gpt-4o \\
        Temperature       & 0.0                 & 0.0 \\
        Diff cap          & 50\,000 chars       & 12\,000 chars \\
        Graph builder     & \multicolumn{2}{l}{code property graph (Joern CPG; \S\ref{sec:approach-kg})} \\
        Chunks shown      & \multicolumn{2}{l}{5 (formatter cap)} \\
        Judges            & \multicolumn{2}{l}{five-provider panel (Table~\ref{tab:models})} \\
        Aggregation       & \multicolumn{2}{l}{strict majority, $\geq 4$ valid votes; ties to 0} \\
        Outcome           & rubric score out of 25 & defect detected \\
        Ground truth      & judge panel         & static oracle \\
        \bottomrule
    \end{tabular}
\end{table}

\subsection{Context Construction}
\label{sec:context-construction}

The four arms are nested. \baseline{} receives the title, description and
diff. \kg{} adds a graph block, \rag{} adds a
retrieval block, and \hybrid{} adds both. Because every arm contains
everything \baseline{} contains, a context arm can only lose by diluting
attention, never by withholding the change.

Both experiments build the graph block from a Joern code property
graph~(\S\ref{sec:approach-kg}). For each changed file, the CPG is
queried for method-level call and import edges that cross
file boundaries; the resulting dependent-file list and edge annotations are
rendered into the prompt.

Retrieval is conventional dense retrieval at the same commit, using the
embedding model in Table~\ref{tab:models}. Each changed file's first
4\,000~characters is embedded as the query; the five highest-scoring chunks
from outside the diff are rendered into the prompt.

\subsection{Models and Roles}
\label{sec:model-inventory}

\begin{table}[htbp]
    \centering
    \caption{Models used for review generation, judging, and embeddings.}
    \label{tab:models}
    \footnotesize
    \begin{tabularx}{\textwidth}{@{}lXc@{}}
        \toprule
        Role & Model & $T$ \\
        \midrule
        Generator (all three studies)
          & GPT-4o & 0.0 \\
        \midrule
        Substitute generators
          & Claude Haiku 4.5; Gemini 2.5 Flash; DeepSeek Chat & 0.0 \\
        \midrule
        Judge panel
          & GPT-4o; Gemini 2.5 Flash; Claude Sonnet 4.5;
            DeepSeek V4 Pro; Grok 4.6 & 0.0 \\
        \midrule
        Embeddings & \texttt{text-embedding-\allowbreak3-small} & --- \\
        \bottomrule
    \end{tabularx}
\end{table}

The APIs were invoked through provider aliases rather than immutable snapshot
identifiers, which limits exact model-level replication.

\subsection{Reproducibility}
\label{sec:repro}

Temperature zero is not bit-deterministic for this model. Regenerating the
graph arm from byte-identical inputs moved the nine-criterion score by
$0.88$~points on average, with fewer than half of the pull requests landing
on exactly the same value. The movement is unbiased (signed difference $-0.19$,
$p = 0.51$), which is why paired means over 35~pull requests remain
trustworthy, but it sets a floor beneath every small effect of roughly the
same size as the graph arm's advantage over the baseline.

What is reproducible is everything except the model calls: the stimulus set,
evidence packs, prompts, judge cache and analysis seed are all on disk. A
replication would produce different review texts and slightly different scores
from the same inputs.
\section{Human Evaluation Study}
\label{sec:human-study}

Both experiments so far rely on language models to score the output of a
language model. The third study asks people.

This study cannot validate the 35-pull-request ranking: the stimuli are a
different sample, the graph block's edges are resolved from import statements rather than
from the Joern CPG, the prompt is stricter, and the criteria are reworded for non-expert readers. What it
tests is whether human readers perceive the specific difference the graph
block produces.

\subsection{Stimuli and Design}

Six pull requests were drawn from three Python libraries, with two from each:
\texttt{psf/requests}, \texttt{pallets/flask}, and \texttt{pallets/click}.
The restriction to one language and to well-known libraries reduces the prior
repository knowledge required of raters.

Each pull request is presented with two reviews: \baseline{} (strict prompt,
no graph block) and \kg{} (a matched prompt that differs only in its
references to the available context, plus the knowledge-graph block). The on-screen diff is shared reference material for both
arms. Reviews are shown as A and B, and participants are not told which mode
produced either review. The side on which each mode appears is randomised per
rater.

The design is fully crossed: every rater sees all six pull requests. Each
task presents a plain-language summary, the two reviews side by side, and
the diff, which is expanded by default. The assisted interface also provides
arm-neutral context, review summaries, file-existence tooltips that do not
certify affectedness, and optional difference highlighting that is off by
default. The rater makes six per-criterion choices (A, B, both, or neither),
states an overall preference (A, equal, or B), optionally explains the choice
in free text, and rates the task's difficulty.
Table~\ref{tab:study-criteria} lists the criteria.

\begin{table}[htbp]
    \centering
    \caption{The six comparison criteria as worded in the instrument. H1 is
             the primary endpoint. IDs are prefixed H to distinguish them
             from the 25-criterion rubric used in Experiments~1 and~2; the
             closest rubric counterparts are F3/F4 for H1, F2 for H2, T3 for
             H3, and R1 for H5.}
    \label{tab:study-criteria}
    \small
    \begin{tabular}{llp{7.4cm}}
        \toprule
        ID & Category & Question put to the rater \\
        \midrule
        H1 & Functionality  & Which review better names the specific functions, classes, or files that this change could affect? \\
        H2 & Functionality  & Which review better describes specific situations where the changed code could fail or misbehave? \\
        H3 & Tests          & Which review better points to specific test files, or says which tests should be added or updated? \\
        H4 & Quality        & Which review better explains \emph{why} something is a problem, rather than just saying it is one? \\
        H5 & Readability    & Which review better comments on how clear and well-organised the code is? \\
        H6 & Completeness   & Which review better covers the important issues in this PR, without obvious gaps? \\
        \bottomrule
    \end{tabular}
\end{table}

\subsection{Analysis Plan}

Choices are scored directionally: a pick of the \kg{} side counts one, a pick
of the \baseline{} side counts zero, and ``both'', ``neither'', or ``equal''
count one half. ``Both'' and ``neither'' share the same directional score but
remain separate in the raw-choice descriptives. A rater's score for a criterion
is the mean over their six tasks; the unit of analysis is the rater and the
null value is~$0.5$.

The primary endpoint is H1 (naming affected components), tested against~$0.5$
with a two-sided Wilcoxon signed-rank test at $\alpha = 0.05$, reported with
a rank-biserial correlation and a bootstrap confidence interval. This
criterion was chosen because the LLM-judged result localises its effect
there, and it is where the two arms differ mechanically. Overall usefulness
is secondary and explicitly estimation-only: power simulations show it
underpowered at the intended sample size. The remaining criteria are
exploratory and Holm-corrected.

The analysis plan also pre-registered four diagnostic checks concerning
file-count association, aid use, free-text rationales, and task time. They were
not implemented in the final analyser and are not used to support any thesis
claim; this omission is reported as a limitation.

The recruitment target is 18--20 complete raters, with 15 as a lower-bound
checkpoint. Only raters who complete all six tasks enter the confirmatory
analysis. The final export exceeded that upper target.
Section~\ref{sec:results-human} reports all 27 complete raters. The
interpretation rules were fixed before data collection: a positive result
licenses the claim that raters perceive the graph arm's better naming of
affected code; a
null result is reported as a boundary finding and the plan forbids re-running
with new stimuli.

\subsection{Ethics}

Participation begins with informed consent. The instrument records a
self-chosen identifier, optional free text, role and experience descriptors,
criterion and overall choices, difficulty, task time, GitHub-link clicks, and
whether highlighting was enabled; it collects neither email nor IP address.
Processing rests on consent under Art.~6(1)(a) GDPR. Data are retained until
examination is complete and deleted no later than 31~December~2027. Only
aggregated or de-identified results appear in this thesis. No formal ethics
board review was required under the university's guidelines for this category
of research; the study collects no personally identifying information and
presents no risk beyond standard professional activity.



    \chapter{Results and Discussion}\label{chap:results}
    \section{Experiment 1: Pull-Request Evaluation}
\label{sec:results-exp1}

This section reports what adding repository context does on 35 pull
requests scored by the rubric. Five Go-only pull requests are excluded
because Joern's Go frontend does not emit the method-level nodes the
queries require. The headline is not that context helps in general (the
structural arm's aggregate effect is small and fragile) but that the two
context types help different things: the structural arm concentrates its advantage in the nine
criteria the mechanism predicts, while retrieval increases rubric coverage
broadly and evenly. That crossing pattern, not any single number, is the
principal result.

\subsection{Paired Comparisons}
\label{sec:results-exp1-paired}

Because every mode reviews every pull request, the informative comparison is
paired. Table~\ref{tab:exp1-deltas} reports each mode against \baseline{}.

\begin{table}[htbp]
    \centering
    \caption{Mean paired differences against \baseline{}, $n = 35$, with
             95\,\% bootstrap intervals.}
    \label{tab:exp1-deltas}
    \small
    \begin{tabular}{lcc}
        \toprule
        Mode & Total /25 $\Delta$ [95\,\% CI] & KG-relevant /9 $\Delta$ [95\,\% CI] \\
        \midrule
        \kg{}     & $+0.71$ [$+0.20$, $+1.26$] & $+0.69$ [$+0.29$, $+1.11$] \\
        \rag{}    & $+1.03$ [$+0.46$, $+1.60$] & $+0.43$ [$-0.03$, $+0.94$] \\
        \hybrid{} & $+0.57$ [$+0.09$, $+1.11$] & $+0.46$ [$+0.03$, $+0.91$] \\
        \bottomrule
    \end{tabular}
\end{table}

The governing pre-specified contrast is the \kg{} gain on the KG-relevant
subscale ($+0.69$, $p = 0.003$, $d_z = +0.56$). \rag{} produces the
largest total-score gain, while \hybrid{} is smaller on both scales. Full
permutation tests and effect sizes are reported in
Appendix~\ref{app:supporting-multiplicity}.
Figure~\ref{fig:exp1-effects} draws both scales together.

\begin{figure}[htbp]
    \centering
    % !TeX root = ../thesis.tex
% GENERATED FILE - do not edit by hand; your changes will be lost.
% Produced by scripts/generate_thesis_figures.py
% Regenerate: python3 scripts/generate_thesis_figures.py
% Computed from:
%   results/JOERN_UNIFIED_FIVE_JUDGE.json (Experiment 1, n = 35)
%
\begin{tikzpicture}
\node[anchor=west,font=\footnotesize] at (-4.5000,0.0000) {\textbf{\kg{}}, KG-relevant /9};
\node[anchor=west,font=\footnotesize] at (-4.5000,-0.5200) {\textbf{\kg{}}, total /25};
\draw[thin,black!55] (-4.5000,-0.7696) -- (6.0000,-0.7696);
\node[anchor=west,font=\footnotesize] at (-4.5000,-1.0400) {\textbf{\rag{}}, KG-relevant /9};
\node[anchor=west,font=\footnotesize] at (-4.5000,-1.5600) {\textbf{\rag{}}, total /25};
\draw[thin,black!55] (-4.5000,-1.8096) -- (6.0000,-1.8096);
\node[anchor=west,font=\footnotesize] at (-4.5000,-2.0800) {\textbf{\hybrid{}}, KG-relevant /9};
\node[anchor=west,font=\footnotesize] at (-4.5000,-2.6000) {\textbf{\hybrid{}}, total /25};
\draw[densely dashed,black!55] (0.3214,0.2340) -- (0.3214,-2.8340);
\draw[kgblue,line width=0.9pt] (1.3571,0.0000) -- (4.2857,0.0000);

\draw[kgblue,line width=0.9pt] (1.3571,-0.0700) -- (1.3571,0.0700);
\draw[kgblue,line width=0.9pt] (4.2857,-0.0700) -- (4.2857,0.0700);
\draw[fill=kgblue,draw=kgblue] (2.7857,0.0000) circle (2.000pt);
\draw[kgblue,line width=0.9pt] (1.0357,-0.5200) -- (4.8214,-0.5200);
\draw[kgblue,line width=0.9pt] (1.0357,-0.5900) -- (1.0357,-0.4500);
\draw[kgblue,line width=0.9pt] (4.8214,-0.5900) -- (4.8214,-0.4500);
\draw[fill=kgblue,draw=kgblue] (2.8571,-0.5200) circle (2.000pt);
\draw[ragorange,line width=0.9pt] (0.2143,-1.0400) -- (3.6785,-1.0400);
\draw[ragorange,line width=0.9pt] (0.2143,-1.1100) -- (0.2143,-0.9700);
\draw[ragorange,line width=0.9pt] (3.6785,-1.1100) -- (3.6785,-0.9700);
\draw[fill=ragorange,draw=ragorange] (1.8571,-1.0400) circle (2.000pt);
\draw[ragorange,line width=0.9pt] (1.9642,-1.5600) -- (6.0000,-1.5600);
\draw[ragorange,line width=0.9pt] (1.9642,-1.6300) -- (1.9642,-1.4900);
\draw[ragorange,line width=0.9pt] (6.0000,-1.6300) -- (6.0000,-1.4900);
\draw[fill=ragorange,draw=ragorange] (4.0000,-1.5600) circle (2.000pt);
\draw[hybridgreen,line width=0.9pt] (0.4285,-2.0800) -- (3.5714,-2.0800);
\draw[hybridgreen,line width=0.9pt] (0.4285,-2.1500) -- (0.4285,-2.0100);
\draw[hybridgreen,line width=0.9pt] (3.5714,-2.1500) -- (3.5714,-2.0100);
\draw[fill=hybridgreen,draw=hybridgreen] (1.9642,-2.0800) circle (2.000pt);
\draw[hybridgreen,line width=0.9pt] (0.6428,-2.6000) -- (4.2857,-2.6000);
\draw[hybridgreen,line width=0.9pt] (0.6428,-2.6700) -- (0.6428,-2.5300);
\draw[hybridgreen,line width=0.9pt] (4.2857,-2.6700) -- (4.2857,-2.5300);
\draw[fill=hybridgreen,draw=hybridgreen] (2.3571,-2.6000) circle (2.000pt);
\draw[thin,black!55] (0.0000,-3.0940) -- (6.5714,-3.0940);
\draw[thin,black!55] (0.3214,-3.0940) -- (0.3214,-3.1740);
\node[anchor=center,font=\scriptsize] at (0.3214,-3.3740) {0};
\draw[thin,black!55] (2.1071,-3.0940) -- (2.1071,-3.1740);
\node[anchor=center,font=\scriptsize] at (2.1071,-3.3740) {0.5};
\draw[thin,black!55] (3.8929,-3.0940) -- (3.8929,-3.1740);
\node[anchor=center,font=\scriptsize] at (3.8929,-3.3740) {1};
\draw[thin,black!55] (5.6786,-3.0940) -- (5.6786,-3.1740);
\node[anchor=center,font=\scriptsize] at (5.6786,-3.3740) {1.5};
\node[anchor=center,font=\footnotesize] at (3.0000,-3.8140) {Paired difference from the diff-only baseline (rubric points)};
\end{tikzpicture}

    \caption{Paired differences against \baseline{} on both scales. Bars are
             95\,\% bootstrap intervals. The modes differ between the total
             rubric and the KG-relevant subscale.}
    \label{fig:exp1-effects}
\end{figure}

The pre-specified \kg{}-versus-\baseline{} contrast remains the governing
analysis. A Friedman sensitivity detects an omnibus mode effect on the total
score ($p = 0.003$) but not on the KG-relevant subscale ($p =
0.084$). The planned contrast remains distinguishable under a
Bonferroni-corrected Wilcoxon test ($p = 0.011$); this is
concordance for that contrast, not evidence of an overall mode effect.
Appendix~\ref{app:supporting-multiplicity} reports the full sensitivity.

\subsection{Criterion-Level Attribution}
\label{sec:results-exp1-attribution}

Table~\ref{tab:exp1-attribution} counts criterion-level wins and losses,
split into the nine KG-relevant criteria and the sixteen others. The sixteen
act as a built-in negative control: a general coverage lift should raise both
groups together.

\begin{table}[htbp]
    \centering
    \caption{Criterion-level wins (W) and losses (L) against \baseline{},
             split into the nine KG-relevant criteria and the sixteen others,
             $n = 35$.}
    \label{tab:exp1-attribution}
    \small
    \begin{tabular}{lcccccc}
        \toprule
        & \multicolumn{3}{c}{KG-relevant (9)} & \multicolumn{3}{c}{Other (16)} \\
        \cmidrule(lr){2-4} \cmidrule(lr){5-7}
        Mode & W & L & net & W & L & net \\
        \midrule
        \kg{}  & 35 & 11 & $+24$ & 14 & 13 & $+1$ \\
        \rag{} & 36 & 21 & $+15$ & 28 & 7 & $+21$ \\
        \bottomrule
    \end{tabular}
\end{table}

This is the chapter's central observation. \kg{} concentrates its net wins
on the criteria it has a mechanism for. Among the other sixteen criteria the
gross movement is not negligible---Q3 gains $+11.4$\,pp while R1 drops $-8.6$\,pp and S1
drops $-5.7$\,pp---but these gains and losses largely cancel, netting to
$+1$. The pattern suggests a displacement effect under the fixed
five-heading format: the review allocates more attention to integration and
triage and less to clarity and input-validation commentary.
The R1 decline is consistent with the human study's later finding that
raters judged the baseline review the better commentary on code clarity
(H5, Section~\ref{sec:results-human}).
\rag{} produces a broad, even lift across both bands. The two modes differ
less in aggregate magnitude than in how selectively they act.

A permutation test confirms the localisation inferentially. If the advantage
were spread at random, the nine KG-relevant criteria would capture 36\,\% of
it. \kg{} places 96\,\% there ($p = 0.003$, $B =
200\,000$); \rag{} places 42\,\% there. This separates the two
mechanisms where aggregate scores cannot.
Figure~\ref{fig:criterion-localisation} shows the per-criterion change in
satisfaction rate for the \kg{} arm relative to the baseline.

\begin{figure}[htbp]
    \centering
    % !TeX root = ../thesis.tex
% GENERATED FILE - do not edit by hand; your changes will be lost.
% Produced by scripts/generate_thesis_figures.py
% Regenerate: python3 scripts/generate_thesis_figures.py
% Computed from:
%   experiments/2026-09-23_five_judge_panel/RESULTS.json (Experiment 1 localisation, n = 40)
%
\begin{tikzpicture}
\node[anchor=east,font=\scriptsize] at (6.5500,0.3060) {$\Delta$\,pp};
\node[anchor=west,font=\footnotesize] at (-5.0000,0.0000) {\emph{KG-relevant band (9 criteria)}};
\node[anchor=west,font=\scriptsize] at (-5.0000,-0.3600) {F3\enspace Integration with existing code};
\fill[kgblue] (1.6710,-0.4824) rectangle (5.2839,-0.2376);
\node[anchor=east,font=\scriptsize] at (6.5500,-0.3600) {+22.9};
\node[anchor=west,font=\scriptsize] at (-5.0000,-0.7200) {F4\enspace Breaking changes or impact};
\fill[kgblue] (1.6710,-0.8424) rectangle (3.9290,-0.5976);
\node[anchor=east,font=\scriptsize] at (6.5500,-0.7200) {+14.3};
\node[anchor=west,font=\scriptsize] at (-5.0000,-1.0800) {T3\enspace Specific test files referenced};
\fill[kgblue] (1.6710,-1.2024) rectangle (3.0258,-0.9576);
\node[anchor=east,font=\scriptsize] at (6.5500,-1.0800) {+8.6};
\node[anchor=west,font=\scriptsize] at (-5.0000,-1.4400) {M3\enspace Public API or interface handling};
\fill[kgblue] (1.6710,-1.5624) rectangle (3.0258,-1.3176);
\node[anchor=east,font=\scriptsize] at (6.5500,-1.4400) {+8.6};
\node[anchor=west,font=\scriptsize] at (-5.0000,-1.8000) {M1\enspace Fit with existing architecture};
\fill[kgblue] (1.6710,-1.9224) rectangle (2.5742,-1.6776);
\node[anchor=east,font=\scriptsize] at (6.5500,-1.8000) {+5.7};
\node[anchor=west,font=\scriptsize] at (-5.0000,-2.1600) {T1\enspace Need for tests discussed};
\fill[kgblue] (1.6710,-2.2824) rectangle (2.1226,-2.0376);
\node[anchor=east,font=\scriptsize] at (6.5500,-2.1600) {+2.9};
\node[anchor=west,font=\scriptsize] at (-5.0000,-2.5200) {T2\enspace Edge cases and error paths in tests};
\fill[kgblue] (1.6710,-2.6424) rectangle (2.1226,-2.3976);
\node[anchor=east,font=\scriptsize] at (6.5500,-2.5200) {+2.9};
\node[anchor=west,font=\scriptsize] at (-5.0000,-2.8800) {C2\enspace Consistency with similar solutions};
\fill[kgblue] (1.6710,-3.0024) rectangle (2.1226,-2.7576);
\node[anchor=east,font=\scriptsize] at (6.5500,-2.8800) {+2.9};
\node[anchor=west,font=\scriptsize,black!55] at (-5.0000,-3.2400) {Q2\enspace Comments anchored to code locations$^{\dagger}$};
\draw[thin,black!55] (-5.0000,-3.5280) -- (6.5500,-3.5280);
\node[anchor=west,font=\footnotesize] at (-5.0000,-3.7260) {\emph{Negative-control band (16 criteria)}};
\node[anchor=west,font=\scriptsize] at (-5.0000,-4.0860) {Q3\enspace Blocking versus non-blocking issues};
\fill[modegrey] (1.6710,-4.2084) rectangle (3.4774,-3.9636);
\node[anchor=east,font=\scriptsize] at (6.5500,-4.0860) {+11.4};
\node[anchor=west,font=\scriptsize] at (-5.0000,-4.4460) {F2\enspace Edge cases in the change identified};
\fill[modegrey] (1.6710,-4.5684) rectangle (2.1226,-4.3236);
\node[anchor=east,font=\scriptsize] at (6.5500,-4.4460) {+2.9};
\node[anchor=west,font=\scriptsize] at (-5.0000,-4.8060) {F1\enspace Change solves the stated problem};
\fill[modegrey] (1.6710,-4.9284) rectangle (2.1226,-4.6836);
\node[anchor=east,font=\scriptsize] at (6.5500,-4.8060) {+2.9};
\node[anchor=west,font=\scriptsize] at (-5.0000,-5.1660) {P1\enspace Obvious inefficiencies};
\fill[modegrey] (1.6710,-5.2884) rectangle (2.1226,-5.0436);
\node[anchor=east,font=\scriptsize] at (6.5500,-5.1660) {+2.9};
\node[anchor=west,font=\scriptsize] at (-5.0000,-5.5260) {S3\enspace Error-handling quality};
\fill[modegrey] (1.6710,-5.6484) rectangle (2.1226,-5.4036);
\node[anchor=east,font=\scriptsize] at (6.5500,-5.5260) {+2.9};
\node[anchor=west,font=\scriptsize,black!55] at (-5.0000,-5.8860) {M2\enspace Pull request too large to review$^{\dagger}$};
\node[anchor=west,font=\scriptsize,black!55] at (-5.0000,-6.2460) {P2\enspace Benchmarks for critical paths$^{\dagger}$};
\node[anchor=west,font=\scriptsize,black!55] at (-5.0000,-6.6060) {Q1\enspace Overall assessment summary$^{\dagger}$};
\node[anchor=west,font=\scriptsize,black!55] at (-5.0000,-6.9660) {Q4\enspace Clarifying questions asked$^{\dagger}$};
\node[anchor=west,font=\scriptsize,black!55] at (-5.0000,-7.3260) {Q5\enspace Rationale given for suggestions$^{\dagger}$};
\node[anchor=west,font=\scriptsize] at (-5.0000,-7.6860) {R2\enspace Unnecessary complexity flagged};
\node[anchor=west,font=\scriptsize,black!55] at (-5.0000,-8.0460) {S2\enspace Hardcoded secrets or credentials$^{\dagger}$};
\node[anchor=west,font=\scriptsize] at (-5.0000,-8.4060) {C1\enspace Project style-guide compliance};
\fill[modegrey!45] (1.2194,-8.5284) rectangle (1.6710,-8.2836);
\node[anchor=east,font=\scriptsize] at (6.5500,-8.4060) {$-$2.9};
\node[anchor=west,font=\scriptsize] at (-5.0000,-8.7660) {R3\enspace Comments explain why, not what};
\fill[modegrey!45] (1.2194,-8.8884) rectangle (1.6710,-8.6436);
\node[anchor=east,font=\scriptsize] at (6.5500,-8.7660) {$-$2.9};
\node[anchor=west,font=\scriptsize] at (-5.0000,-9.1260) {S1\enspace Input validation};
\fill[modegrey!45] (0.7677,-9.2484) rectangle (1.6710,-9.0036);
\node[anchor=east,font=\scriptsize] at (6.5500,-9.1260) {$-$5.7};
\node[anchor=west,font=\scriptsize] at (-5.0000,-9.4860) {R1\enspace Code clarity, naming, function size};
\fill[modegrey!45] (0.3161,-9.6084) rectangle (1.6710,-9.3636);
\node[anchor=east,font=\scriptsize] at (6.5500,-9.4860) {$-$8.6};
\draw[densely dashed,black!55] (1.6710,0.2160) -- (1.6710,-9.6660);
\draw[thin,black!55] (0.0000,-9.8100) -- (5.6000,-9.8100);
\draw[thin,black!55] (0.0903,-9.8100) -- (0.0903,-9.8900);
\node[anchor=center,font=\scriptsize] at (0.0903,-10.0900) {$-$10};
\draw[thin,black!55] (1.6710,-9.8100) -- (1.6710,-9.8900);
\node[anchor=center,font=\scriptsize] at (1.6710,-10.0900) {0};
\draw[thin,black!55] (3.2516,-9.8100) -- (3.2516,-9.8900);
\node[anchor=center,font=\scriptsize] at (3.2516,-10.0900) {10};
\draw[thin,black!55] (4.8323,-9.8100) -- (4.8323,-9.8900);
\node[anchor=center,font=\scriptsize] at (4.8323,-10.0900) {20};
\node[anchor=center,font=\footnotesize] at (2.8000,-10.5300) {Change in satisfaction rate, KG versus baseline (percentage points)};
\end{tikzpicture}

    \caption{Change in criterion satisfaction rate, \kg{} versus \baseline{},
             sorted by magnitude within each band. Daggers mark criteria at
             ceiling or floor in both arms.}
    \label{fig:criterion-localisation}
\end{figure}

\subsection{Inter-Judge Agreement}
\label{sec:results-exp1-agreement}

Pairwise Cohen's $\kappa$ among the five judges ranges from $0.701$ to
$0.893$, substantial to almost perfect under the interpretation in
Section~\ref{sec:bg-stats}. This supports the stability of majority
aggregation, although agreement does not establish that the shared verdict is
correct. Appendix~\ref{app:supporting-agreement} reports all ten pairs,
raw agreement, valid-cell counts, and abstentions.
% Source: experiments/2026-09-23_five_judge_panel/RESULTS.md, lines 32--47.

\subsection{Robustness}
\label{sec:results-exp1-robustness}

The effect was recomputed under several substitutions, varying one element
each time. Figure~\ref{fig:exp1-robustness} draws the estimates on one axis.
The primary row uses the final five-judge panel. The judge rows use one
judge each. The generator and held-out rows are from an earlier
configuration (text-heuristic builder, $n = 40$, three-judge panel) and are
retained as directional evidence. The final row group shows an earlier
generation of the same Joern CPG configuration under the three-judge panel,
which produced $+0.34$ ($p = 0.111$); the difference from the primary row
illustrates the regeneration variance discussed in Section~\ref{sec:repro}. Detailed statistics are in
Appendix~\ref{app:supporting-robustness}.

\begin{figure}[htbp]
    \centering
    % !TeX root = ../thesis.tex
% GENERATED FILE - do not edit by hand; your changes will be lost.
% Produced by scripts/generate_thesis_figures.py
% Regenerate: python3 scripts/generate_thesis_figures.py
% Computed from:
%   experiments/2026-09-23_five_judge_panel/RESULTS.json (headline and individual judges)
%   results/CROSS_GENERATOR_v2.json (generators)
%   results/BOOTSTRAP_STATS_joern_parity.json (builder parity)
%   results/BOOTSTRAP_STATS_confirmatory_clean.json (held-out)
%
\begin{tikzpicture}
\node[anchor=west,font=\footnotesize] at (-5.2000,0.0000) {\emph{Primary configuration}};
\node[anchor=west,font=\footnotesize] at (-5.2000,-0.4600) {Five-judge panel, $n = 35$};
\draw[thin,black!55] (-5.2000,-0.6808) -- (5.4000,-0.6808);
\node[anchor=west,font=\footnotesize] at (-5.2000,-0.9200) {\emph{Single-judge sensitivity}};
\node[anchor=west,font=\footnotesize] at (-5.2000,-1.3800) {gpt-4o only (also the generator)};
\node[anchor=west,font=\footnotesize] at (-5.2000,-1.8400) {gemini-2.5-flash only};
\node[anchor=west,font=\footnotesize] at (-5.2000,-2.3000) {claude-sonnet-4.5 only};
\node[anchor=west,font=\footnotesize] at (-5.2000,-2.7600) {deepseek-v4-pro only};
\node[anchor=west,font=\footnotesize] at (-5.2000,-3.2200) {grok-4.6 only};
\draw[thin,black!55] (-5.2000,-3.4408) -- (5.4000,-3.4408);
\node[anchor=west,font=\footnotesize] at (-5.2000,-3.6800) {\emph{Alternative generators}};
\node[anchor=west,font=\footnotesize] at (-5.2000,-4.1400) {claude-haiku-4.5};
\node[anchor=west,font=\footnotesize] at (-5.2000,-4.6000) {gemini-2.5-flash};
\node[anchor=west,font=\footnotesize] at (-5.2000,-5.0600) {deepseek-chat};
\draw[thin,black!55] (-5.2000,-5.2808) -- (5.4000,-5.2808);
\node[anchor=west,font=\footnotesize] at (-5.2000,-5.5200) {\emph{Regeneration and sample sensitivity}};
\node[anchor=west,font=\footnotesize] at (-5.2000,-5.9800) {Earlier generation, same CPG, three judges};
\node[anchor=west,font=\footnotesize] at (-5.2000,-6.4400) {Held-out pull requests, $n = 12$};
\draw[densely dashed,black!55] (1.1238,-0.2530) -- (1.1238,-6.6470);
\draw[kgblue,line width=0.9pt] (1.7728,-0.4600) -- (5.1107,-0.4600);
\draw[kgblue,line width=0.9pt] (1.7728,-0.5300) -- (1.7728,-0.3900);
\draw[kgblue,line width=0.9pt] (5.1107,-0.5300) -- (5.1107,-0.3900);
\draw[fill=kgblue,draw=kgblue] (3.4418,-0.4600) circle (2.000pt);
\draw[kgblue,line width=0.9pt] (1.9582,-1.3800) -- (5.0180,-1.3800);
\draw[kgblue,line width=0.9pt] (1.9582,-1.4500) -- (1.9582,-1.3100);
\draw[kgblue,line width=0.9pt] (5.0180,-1.4500) -- (5.0180,-1.3100);
\draw[fill=kgblue,draw=kgblue] (3.5345,-1.3800) circle (2.000pt);
\draw[kgblue,line width=0.9pt] (0.2893,-1.8400) -- (4.7398,-1.8400);
\draw[kgblue,line width=0.9pt] (0.2893,-1.9100) -- (0.2893,-1.7700);
\draw[kgblue,line width=0.9pt] (4.7398,-1.9100) -- (4.7398,-1.7700);
\draw[fill=kgblue,draw=kgblue] (2.5146,-1.8400) circle (2.000pt);
\draw[kgblue,line width=0.9pt] (1.4019,-2.3000) -- (3.8126,-2.3000);
\draw[kgblue,line width=0.9pt] (1.4019,-2.3700) -- (1.4019,-2.2300);
\draw[kgblue,line width=0.9pt] (3.8126,-2.3700) -- (3.8126,-2.2300);
\draw[fill=kgblue,draw=kgblue] (2.6073,-2.3000) circle (2.000pt);
\draw[kgblue,line width=0.9pt] (1.8655,-2.7600) -- (4.8326,-2.7600);
\draw[kgblue,line width=0.9pt] (1.8655,-2.8300) -- (1.8655,-2.6900);
\draw[kgblue,line width=0.9pt] (4.8326,-2.8300) -- (4.8326,-2.6900);
\draw[fill=kgblue,draw=kgblue] (3.3490,-2.7600) circle (2.000pt);
\draw[kgblue,line width=0.9pt] (1.4946,-3.2200) -- (4.0908,-3.2200);
\draw[kgblue,line width=0.9pt] (1.4946,-3.2900) -- (1.4946,-3.1500);
\draw[kgblue,line width=0.9pt] (4.0908,-3.2900) -- (4.0908,-3.1500);
\draw[fill=kgblue,draw=kgblue] (2.7927,-3.2200) circle (2.000pt);
\draw[kgblue,line width=0.9pt] (0.4747,-4.1400) -- (2.8854,-4.1400);
\draw[kgblue,line width=0.9pt] (0.4747,-4.2100) -- (0.4747,-4.0700);
\draw[kgblue,line width=0.9pt] (2.8854,-4.2100) -- (2.8854,-4.0700);
\draw[fill=kgblue,draw=kgblue] (1.6801,-4.1400) circle (2.000pt);
\draw[kgblue,line width=0.9pt] (0.7529,-4.6000) -- (5.0180,-4.6000);
\draw[kgblue,line width=0.9pt] (0.7529,-4.6700) -- (0.7529,-4.5300);
\draw[kgblue,line width=0.9pt] (5.0180,-4.6700) -- (5.0180,-4.5300);
\draw[fill=kgblue,draw=kgblue] (2.8854,-4.6000) circle (2.000pt);
\draw[kgblue,line width=0.9pt] (0.4747,-5.0600) -- (3.2563,-5.0600);
\draw[kgblue,line width=0.9pt] (0.4747,-5.1300) -- (0.4747,-4.9900);
\draw[kgblue,line width=0.9pt] (3.2563,-5.1300) -- (3.2563,-4.9900);
\draw[fill=kgblue,draw=kgblue] (1.8655,-5.0600) circle (2.000pt);
\draw[kgblue,line width=0.9pt] (1.0162,-5.9800) -- (3.7718,-5.9800);
\draw[kgblue,line width=0.9pt] (1.0162,-6.0500) -- (1.0162,-5.9100);
\draw[kgblue,line width=0.9pt] (3.7718,-6.0500) -- (3.7718,-5.9100);
\draw[fill=kgblue,draw=kgblue] (2.3959,-5.9800) circle (2.000pt);
\draw[kgblue,line width=0.9pt] (0.5044,-6.4400) -- (4.8326,-6.4400);
\draw[kgblue,line width=0.9pt] (0.5044,-6.5100) -- (0.5044,-6.3700);
\draw[kgblue,line width=0.9pt] (4.8326,-6.5100) -- (4.8326,-6.3700);
\draw[fill=kgblue,draw=kgblue] (2.6703,-6.4400) circle (2.000pt);
\draw[thin,black!55] (0.0000,-6.8770) -- (5.4000,-6.8770);
\draw[thin,black!55] (0.1966,-6.8770) -- (0.1966,-6.9570);
\node[anchor=center,font=\scriptsize] at (0.1966,-7.1570) {$-$0.25};
\draw[thin,black!55] (1.1238,-6.8770) -- (1.1238,-6.9570);
\node[anchor=center,font=\scriptsize] at (1.1238,-7.1570) {0};
\draw[thin,black!55] (2.0510,-6.8770) -- (2.0510,-6.9570);
\node[anchor=center,font=\scriptsize] at (2.0510,-7.1570) {0.25};
\draw[thin,black!55] (2.9782,-6.8770) -- (2.9782,-6.9570);
\node[anchor=center,font=\scriptsize] at (2.9782,-7.1570) {0.5};
\draw[thin,black!55] (3.9054,-6.8770) -- (3.9054,-6.9570);
\node[anchor=center,font=\scriptsize] at (3.9054,-7.1570) {0.75};
\draw[thin,black!55] (4.8326,-6.8770) -- (4.8326,-6.9570);
\node[anchor=center,font=\scriptsize] at (4.8326,-7.1570) {1};
\node[anchor=center,font=\footnotesize] at (2.7000,-7.5970) {Effect on the nine KG-relevant criteria (rubric points)};
\end{tikzpicture}

    \caption{The KG-relevant effect under alternative configurations. Points
             are mean paired effects and bars are 95\,\% bootstrap intervals.}
    \label{fig:exp1-robustness}
\end{figure}

Every recomputation preserves the positive direction, but the magnitude and
uncertainty depend on the judge, generator, and sample. The criterion-level
localisation is more stable than the aggregate magnitude.

\section{Experiment 2: Controlled Defect Injection}
\label{sec:results-exp2}

\subsection{Detection on the Structural Bands}
\label{sec:results-exp2-structural}

Figure~\ref{fig:exp2-detection} shows detection rates on the 28 structural
injections, where the consequence lies in a different file from the edit.%
\footnote{Numbers reflect the corrected two-stage adjudication described in
Section~\ref{sec:injection-design}, scored by the five-judge panel. Exact
rates and Wilson intervals are in
Table~\ref{tab:exp2-detection} (Appendix).}

\begin{figure}[htbp]
    \centering
    % !TeX root = ../thesis.tex
% GENERATED FILE - do not edit by hand; your changes will be lost.
% Produced by scripts/generate_thesis_figures.py
% Regenerate: python3 scripts/generate_thesis_figures.py
% Computed from:
%   experiments/2026-09-23_five_judge_panel/RESULTS.json (Experiment 2 detection)
%
\begin{tikzpicture}
\node[anchor=west,font=\footnotesize] at (-4.4000,0.0000) {\baseline{} (diff only)};
\draw[black!70,line width=0.5pt] (0.0000,0.0000) -- (0.7239,0.0000);
\draw[black!70,line width=0.5pt] (0.0000,-0.0550) -- (0.0000,0.0550);
\draw[black!70,line width=0.5pt] (0.7239,-0.0550) -- (0.7239,0.0550);
\draw[fill=white,draw=black!75,line width=0.6pt] (6.0000,0.0000) circle (2.100pt);
\node[anchor=east,font=\scriptsize] at (7.5000,0.0000) {0/28};
\node[anchor=west,font=\footnotesize] at (-4.4000,-0.5000) {\rag{}};
\fill[ragorange] (0.0000,-0.6450) rectangle (1.0714,-0.3550);
\draw[black!70,line width=0.5pt] (0.4727,-0.5000) -- (2.1355,-0.5000);
\draw[black!70,line width=0.5pt] (0.4727,-0.5550) -- (0.4727,-0.4450);
\draw[black!70,line width=0.5pt] (2.1355,-0.5550) -- (2.1355,-0.4450);
\draw[fill=white,draw=black!75,line width=0.6pt] (6.0000,-0.5000) circle (2.100pt);
\node[anchor=east,font=\scriptsize] at (7.5000,-0.5000) {5/28};
\node[anchor=west,font=\footnotesize] at (-4.4000,-1.0000) {\hybrid{}};
\fill[hybridgreen] (0.0000,-1.1450) rectangle (4.5000,-0.8550);
\draw[black!70,line width=0.5pt] (3.3987,-1.0000) -- (5.2394,-1.0000);
\draw[black!70,line width=0.5pt] (3.3987,-1.0550) -- (3.3987,-0.9450);
\draw[black!70,line width=0.5pt] (5.2394,-1.0550) -- (5.2394,-0.9450);
\draw[fill=white,draw=black!75,line width=0.6pt] (5.5000,-1.0000) circle (2.100pt);
\node[anchor=east,font=\scriptsize] at (7.5000,-1.0000) {21/28};
\node[anchor=west,font=\footnotesize] at (-4.4000,-1.5000) {\kg{} (deployed)};
\fill[kgblue] (0.0000,-1.6450) rectangle (3.8571,-1.3550);
\draw[black!70,line width=0.5pt] (2.7498,-1.5000) -- (4.7577,-1.5000);
\draw[black!70,line width=0.5pt] (2.7498,-1.5550) -- (2.7498,-1.4450);
\draw[black!70,line width=0.5pt] (4.7577,-1.5550) -- (4.7577,-1.4450);
\draw[fill=white,draw=black!75,line width=0.6pt] (6.0000,-1.5000) circle (2.100pt);
\node[anchor=east,font=\scriptsize] at (7.5000,-1.5000) {18/28};
\node[anchor=west,font=\footnotesize] at (-4.4000,-2.0000) {\kg{} + inheritance edges};
\fill[kgblue] (0.0000,-2.1450) rectangle (5.5714,-1.8550);
\draw[black!70,line width=0.5pt] (4.6413,-2.0000) -- (5.8811,-2.0000);
\draw[black!70,line width=0.5pt] (4.6413,-2.0550) -- (4.6413,-1.9450);
\draw[black!70,line width=0.5pt] (5.8811,-2.0550) -- (5.8811,-1.9450);
\draw[fill=white,draw=black!75,line width=0.6pt] (6.0000,-2.0000) circle (2.100pt);
\node[anchor=east,font=\scriptsize] at (7.5000,-2.0000) {26/28};
\node[anchor=west,font=\footnotesize] at (-4.4000,-2.5000) {\kg{} idealised (ceiling)};
\fill[kgblue] (0.0000,-2.6450) rectangle (5.5714,-2.3550);
\draw[black!70,line width=0.5pt] (4.6413,-2.5000) -- (5.8811,-2.5000);
\draw[black!70,line width=0.5pt] (4.6413,-2.5550) -- (4.6413,-2.4450);
\draw[black!70,line width=0.5pt] (5.8811,-2.5550) -- (5.8811,-2.4450);
\draw[fill=white,draw=black!75,line width=0.6pt] (6.0000,-2.5000) circle (2.100pt);
\node[anchor=east,font=\scriptsize] at (7.5000,-2.5000) {26/28};
\node[anchor=east,font=\scriptsize] at (7.5000,0.4000) {detected};
\draw[thin,black!55] (0.0000,-2.7750) -- (6.0000,-2.7750);
\draw[thin,black!55] (0.0000,-2.7750) -- (0.0000,-2.8550);
\node[anchor=center,font=\scriptsize] at (0.0000,-3.0550) {0};
\draw[thin,black!55] (1.2000,-2.7750) -- (1.2000,-2.8550);
\node[anchor=center,font=\scriptsize] at (1.2000,-3.0550) {0.2};
\draw[thin,black!55] (2.4000,-2.7750) -- (2.4000,-2.8550);
\node[anchor=center,font=\scriptsize] at (2.4000,-3.0550) {0.4};
\draw[thin,black!55] (3.6000,-2.7750) -- (3.6000,-2.8550);
\node[anchor=center,font=\scriptsize] at (3.6000,-3.0550) {0.6};
\draw[thin,black!55] (4.8000,-2.7750) -- (4.8000,-2.8550);
\node[anchor=center,font=\scriptsize] at (4.8000,-3.0550) {0.8};
\draw[thin,black!55] (6.0000,-2.7750) -- (6.0000,-2.8550);
\node[anchor=center,font=\scriptsize] at (6.0000,-3.0550) {1};
\node[anchor=center,font=\footnotesize] at (3.0000,-3.4950) {Detection rate};
\fill[black!70] (0.0000,-4.3910) rectangle (0.4200,-4.1590);
\node[anchor=west,font=\scriptsize] at (0.5500,-4.2750) {structural injections, consequence outside the diff ($n = 28$), with 95\,\% CI};
\draw[fill=white,draw=black!75,line width=0.6pt] (0.2100,-4.6750) circle (2.100pt);
\node[anchor=west,font=\scriptsize] at (0.5500,-4.6750) {local control, consequence inside the diff ($n = 12$)};
\end{tikzpicture}

    \caption{Detection rate by arm on both bands. Bars are structural
             injections with Wilson 95\,\% intervals; open markers are local
             controls. Structural-band rates increase from zero to
             near-ceiling, while local-control rates stay at or near
             ceiling.}
    \label{fig:exp2-detection}
\end{figure}

The baseline detects nothing (zero of 28) while the deployed graph arm
detects eighteen. The exact McNemar contrast is $+0.64$, $p < 0.0001$,
clearing the pre-registered threshold of $+0.30$ at $p < 0.05$.

Retrieval detects five of 28. On pull requests it had the best total coverage;
on defects whose
consequences sit outside the diff it remains substantially below the graph.
Embedding similarity retrieves code that \emph{resembles} the change, but a
file that breaks because it calls a renamed function need not resemble it.
The two mechanisms are not competing implementations of the same idea.

\hybrid{} detects twenty-one, numerically above the graph arm. The paired
difference does not provide evidence that the hybrid outperforms the graph
arm.

\subsection{The Local Control Bands}
\label{sec:results-exp2-control}

On the twelve local injections, whose consequences are visible in the diff,
all arms detect at or near ceiling (the hybrid arm misses one local case;
every other arm scores $12/12$). The graph block's advantage appears exactly
where structural reasoning is required and vanishes where it is not.

The graph block contains the dependents' filenames, so naming one is not
by itself proof of reasoning. Judges must confirm that the review links
the named file to the change, and filtering the block to true-dependent
edges raises detection further (Appendix~\ref{app:supporting-components}),
so what is named matters. A block of the same size with edges to unrelated
symbols yields 1 of 28, so the detections do not come from memory of these
repositories.

\subsection{Effect of Inheritance Edges}
\label{sec:results-exp2-inheritance}

The deployed builder carries no inheritance edges. Adding them raises
detection from 18 to 26 of 28, matching the idealised resolver. With the model
held fixed, this attributes the eight
recovered cases to relations omitted by the deployed builder. The two
remaining misses persist under the idealised resolver and cannot be assigned
to graph construction alone: one review identifies affected classes without
satisfying the exact filename gate, while the other remains generic despite
complete evidence. Perfect structural evidence therefore does not guarantee a
specific review comment. Appendix~\ref{app:supporting-exp2-misses} analyses
both cases; Appendix~\ref{app:supporting-exp2-contrasts} reports the secondary
paired contrasts.

\subsection{Supporting Mechanism Probes}

Component probes isolate dependency lists, resolved call edges, test-file
links, and matched-volume noise. Across their respective configurations they
support the same structural interpretation: relevant relationships drive
detection, while unrelated context does not. Appendix
\ref{app:supporting-components} reports the protocols and statistics.

\section{Human Evaluation Results}
\label{sec:results-human}

The study closed at 27 complete raters, above the 18--20 recruitment target,
giving 162 pairwise comparisons. Every rater who finished all six tasks is
included; one incomplete session is excluded, as specified before data
collection. Participants self-reported as students~(10), junior
developers~(7), mid-level developers~(8), and senior developers~(2), with
programming experience ranging from under one year~(3) through one to
three~(9), three to five~(12), and five to ten years~(3); 19 of 27
reported reviewing pull requests at least weekly.

\subsection{Primary Endpoint}
\label{sec:results-human-primary}

On H1 (which review better names the specific functions, classes, or files
the change could affect), the mean rater preference for \kg{} is $0.710$,
with a bootstrap interval of $[0.623, 0.787]$ excluding the no-preference
value of~$0.5$. Wilcoxon signed-rank $p = 0.0005$ against~$0.5$, with
rank-biserial $r = 0.80$. Across the 162 comparisons, raters picked the
\kg{} review 98 times and the \baseline{} review 30 times, with 29 ties and
5 ``neither'' judgements.

\subsection{Overall Usefulness}
\label{sec:results-human-overall}

Mean preference on overall usefulness is $0.512$, $[0.423, 0.605]$.
The interval includes $0.5$, so the data are consistent with no preference.
This question was set aside for estimation before the study started, so no
significance test is reported. Raters preferred the graph review when asked
which review names the affected code, and they showed no preference when asked
which review was more useful overall. The result is therefore specific to
affected-code identification.

\subsection{Exploratory Boundary}

Among the five exploratory criteria, only code clarity remains distinguishable
after correction, and it favours the \baseline{} review. The graph review is
ahead only on the primary affected-code question.
Figure~\ref{fig:human-study} summarises all criteria on one axis;
Table~\ref{tab:human-exploratory} (Appendix) reports the exploratory
estimates and Holm-corrected $p$-values.

\begin{figure}[htbp]
    \centering
    % !TeX root = ../thesis.tex
% GENERATED FILE - do not edit by hand; your changes will be lost.
% Produced by scripts/generate_thesis_figures.py
% Regenerate: python3 scripts/generate_thesis_figures.py
% Computed from:
%   experiments/2026-07-06_user_study_prs/RESULTS_HUMAN_V4.json (n = 27 raters, 6 pull requests each)
%
\begin{tikzpicture}
\node[anchor=west,font=\footnotesize] at (-5.6000,0.0000) {\emph{Primary (confirmatory)}};
\node[anchor=west,font=\footnotesize] at (-5.6000,-0.4800) {\textbf{H1} names affected code};
\draw[thin,black!55] (-5.6000,-0.7104) -- (6.8500,-0.7104);
\node[anchor=west,font=\footnotesize] at (-5.6000,-0.9600) {\emph{Secondary (estimation only)}};
\node[anchor=west,font=\footnotesize] at (-5.6000,-1.4400) {\textbf{Overall} usefulness};
\draw[thin,black!55] (-5.6000,-1.6704) -- (6.8500,-1.6704);
\node[anchor=west,font=\footnotesize] at (-5.6000,-1.9200) {\emph{Exploratory (Holm-corrected)}};
\node[anchor=west,font=\footnotesize] at (-5.6000,-2.4000) {\textbf{H2} describes failure situations};
\node[anchor=west,font=\footnotesize] at (-5.6000,-2.8800) {\textbf{H3} points to test files};
\node[anchor=west,font=\footnotesize] at (-5.6000,-3.3600) {\textbf{H4} explains why something is a problem};
\node[anchor=west,font=\footnotesize] at (-5.6000,-3.8400) {\textbf{H5} comments on code clarity};
\node[anchor=west,font=\footnotesize] at (-5.6000,-4.3200) {\textbf{H6} covers the important issues};
\draw[densely dashed,black!55] (1.9304,-0.2640) -- (1.9304,-4.5360);
\draw[kgblue,line width=0.9pt] (3.4159,-0.4800) -- (5.3000,-0.4800);
\draw[kgblue,line width=0.9pt] (3.4159,-0.5500) -- (3.4159,-0.4100);
\draw[kgblue,line width=0.9pt] (5.3000,-0.5500) -- (5.3000,-0.4100);
\draw[fill=kgblue,draw=kgblue] (4.3942,-0.4800) circle (2.000pt);
\node[anchor=east,font=\scriptsize] at (6.8500,-0.4800) {0.0005};
\draw[modegrey,line width=0.9pt] (1.0246,-1.4400) -- (3.1623,-1.4400);
\draw[modegrey,line width=0.9pt] (1.0246,-1.5100) -- (1.0246,-1.3700);
\draw[modegrey,line width=0.9pt] (3.1623,-1.5100) -- (3.1623,-1.3700);
\draw[fill=white,draw=modegrey,line width=0.6pt] (2.0754,-1.4400) circle (2.000pt);
\node[anchor=east,font=\scriptsize] at (6.8500,-1.4400) {---};
\draw[modegrey,line width=0.9pt] (0.9159,-2.4000) -- (2.2565,-2.4000);
\draw[modegrey,line width=0.9pt] (0.9159,-2.4700) -- (0.9159,-2.3300);
\draw[modegrey,line width=0.9pt] (2.2565,-2.4700) -- (2.2565,-2.3300);
\draw[fill=white,draw=modegrey,line width=0.6pt] (1.5681,-2.4000) circle (2.000pt);
\node[anchor=east,font=\scriptsize] at (6.8500,-2.4000) {0.499};
\draw[modegrey,line width=0.9pt] (0.6986,-2.8800) -- (2.0754,-2.8800);
\draw[modegrey,line width=0.9pt] (0.6986,-2.9500) -- (0.6986,-2.8100);
\draw[modegrey,line width=0.9pt] (2.0754,-2.9500) -- (2.0754,-2.8100);
\draw[fill=white,draw=modegrey,line width=0.6pt] (1.3870,-2.8800) circle (2.000pt);
\node[anchor=east,font=\scriptsize] at (6.8500,-2.8800) {0.499};
\draw[modegrey,line width=0.9pt] (1.2058,-3.3600) -- (2.4377,-3.3600);
\draw[modegrey,line width=0.9pt] (1.2058,-3.4300) -- (1.2058,-3.2900);
\draw[modegrey,line width=0.9pt] (2.4377,-3.4300) -- (2.4377,-3.2900);
\draw[fill=white,draw=modegrey,line width=0.6pt] (1.8217,-3.3600) circle (2.000pt);
\node[anchor=east,font=\scriptsize] at (6.8500,-3.3600) {1.000};
\draw[modegrey,line width=0.9pt] (0.3000,-3.8400) -- (1.6768,-3.8400);
\draw[modegrey,line width=0.9pt] (0.3000,-3.9100) -- (0.3000,-3.7700);
\draw[modegrey,line width=0.9pt] (1.6768,-3.9100) -- (1.6768,-3.7700);
\draw[fill=modegrey,draw=modegrey] (0.9159,-3.8400) circle (2.000pt);
\node[anchor=east,font=\scriptsize] at (6.8500,-3.8400) {0.028};
\draw[modegrey,line width=0.9pt] (0.9522,-4.3200) -- (2.5101,-4.3200);
\draw[modegrey,line width=0.9pt] (0.9522,-4.3900) -- (0.9522,-4.2500);
\draw[modegrey,line width=0.9pt] (2.5101,-4.3900) -- (2.5101,-4.2500);
\draw[fill=white,draw=modegrey,line width=0.6pt] (1.7130,-4.3200) circle (2.000pt);
\node[anchor=east,font=\scriptsize] at (6.8500,-4.3200) {1.000};
\draw[thin,black!55] (0.0000,-4.7760) -- (5.6000,-4.7760);
\draw[thin,black!55] (0.7565,-4.7760) -- (0.7565,-4.8560);
\node[anchor=center,font=\scriptsize] at (0.7565,-5.0560) {0.4};
\draw[thin,black!55] (1.9304,-4.7760) -- (1.9304,-4.8560);
\node[anchor=center,font=\scriptsize] at (1.9304,-5.0560) {0.5};
\draw[thin,black!55] (3.1043,-4.7760) -- (3.1043,-4.8560);
\node[anchor=center,font=\scriptsize] at (3.1043,-5.0560) {0.6};
\draw[thin,black!55] (4.2783,-4.7760) -- (4.2783,-4.8560);
\node[anchor=center,font=\scriptsize] at (4.2783,-5.0560) {0.7};
\draw[thin,black!55] (5.4522,-4.7760) -- (5.4522,-4.8560);
\node[anchor=center,font=\scriptsize] at (5.4522,-5.0560) {0.8};
\node[anchor=center,font=\footnotesize] at (2.8000,-5.4960) {Rater preference for the \kg{} review ($0.5 =$ no preference)};
\node[anchor=east,font=\scriptsize] at (6.8500,0.0000) {$p$};
\end{tikzpicture}

    \caption{Rater preference for the \kg{} review by criterion ($n = 27$).
             Bars are bootstrap intervals; the dashed line is no preference.
             Filled markers indicate significance under the governing test.}
    \label{fig:human-study}
\end{figure}

\section{Discussion}
\label{sec:discussion}

Three studies have been reported: an observational comparison, a controlled
injection experiment, and a human preference study. Together they distinguish
targeted affected-code grounding from general review quality.

\subsection{Reconciling Capability and Aggregate Gains}
\label{sec:discussion-reconciliation}

The two experiments appear to disagree about effect size by an order of
magnitude. They estimate different quantities. Experiment~2 conditions on
the mechanism applying: every injected defect was placed so that its
consequence lay outside the edited file. It measures a capability (can the
reviewer see such a consequence at all?) and answers that a diff-only
reviewer cannot and a graph-equipped one often can (18/28 in this injection
sample). Experiment~1 samples pull requests without conditioning on that
mechanism, and its average graph effect is correspondingly smaller.

The graph gives the reviewer a new ability, while its average effect is small
on the 35 pull requests in Experiment~1. The experiment does not estimate how
often that ability would matter in the broader population of pull requests.

The join is conceptual rather than based on a shared criterion label.
Experiment~1's advantage concentrates in the nine structurally relevant
criteria, with integration (F3) and breaking-change impact (F4) among the
largest movers. Experiment~2 directly tests whether a review names an affected
dependent outside the diff. The human study's primary item (H1) asks which
review better names affected code.
These results converge on affected-code grounding. The exploratory clarity
item (H5) favours the baseline, which
further limits any claim of general review improvement.

\subsection{Complementarity of Structure and Similarity}
\label{sec:discussion-complementarity}

Retrieval answers a similarity question: which parts of this repository
resemble the change. That lifts many criteria a little. A dependency graph
answers a reachability question: what is connected to what changed. That
helps a narrow set of criteria a great deal and most criteria not at all.

Experiment~2 shows the two are not interchangeable. Retrieval scored 5/28
on cross-file defects against the graph's 18/28. A file that breaks because
it calls a renamed function need not resemble the file that changed.
Retrieval is not a substitute for structural cross-file detection, although it
remains competitive on broader rubric coverage.

The hybrid arm detects twenty-one of twenty-eight, numerically above both
inputs. The paired difference against the graph is not significant
($p = 0.375$), so it provides evidence of neither superiority nor dilution.
On the pull requests the picture is less favourable: the hybrid trails the
leading single arm on both scales and both arms on the total
(Table~\ref{tab:exp1-deltas}), so under naive concatenation the rubric
shows no additive gain.
Together with the crossing pattern in Experiment~1, the result is consistent
with complementary mechanisms, but does not show that naive concatenation
combines their benefits.

\subsection{Implications for Evaluation Design}
\label{sec:discussion-measurement}

The rubric estimate has two measured weaknesses. It is panel-dependent: the
same reviews yield effects ranging from significant to not, depending on
which judges are polled. And it is generator-dependent: three substitute
generators yield no significant effect
(Appendix~\ref{app:supporting-robustness}).

The practical recommendation generalises past this thesis: a
generate-then-LLM-judge design should not be the sole evidence for a context
mechanism. Where a mechanism can be stated precisely enough to build an
oracle, the oracle should be built.

\subsection{Practical Implications}
\label{sec:discussion-practical}

Commercial review tools such as CodeRabbit and Greptile already index
repositories and hand the model pre-digested dependency
findings~\cite{coderabbit2026,greptile2026}; this thesis shows that a model
given raw graph context can itself name affected code outside the diff, with
the gain localised to cross-file criteria. Three positions follow for
pipeline builders.

\emph{Supply the graph selectively.} The graph earns its place on changes
with cross-file consequences and contributes little elsewhere. A pull
request's structural fan-out can be computed before any review is generated,
so the decision whether to include graph context can itself be automated.

\emph{Resolve structure at function granularity where the analyser supports
it.} The human study's graph block carries file-level import edges, and a
module-level import does not imply that a change reaches the importer. This is a precision argument, not evidence that
function-level edges are always necessary: in Experiment~2 the file-level
dependency list alone recovers sixteen of the graph arm's eighteen
detections.

\emph{Invest in the analyser, not only the model.} Within Experiment~2 the
detection ladder moves substantially when the analyser changes and the model
is held fixed. Whether analyser improvements should take priority for
broader pull-request review remains untested.

\section{Threats to Validity}
\label{sec:threats}

\paragraph{Construct validity.}
The rubric asks whether a review addresses a topic, not whether what it says
is true. A hallucinated review scored $9.2$ against a real review's $9.0$.
All rubric numbers are therefore coverage measurements; correctness claims
rest on Experiment~2's oracle. The human-study interface verifies only that a
referenced file exists, not that the change genuinely affects it.
Binary scoring and seven immovable criteria (at $0\,\%$ or $100\,\%$ in both
arms) further dilute the aggregate, which is one reason the pre-specified
subscale is reported alongside the total. The graph arm also writes the
longest reviews ($+12\,\%$ in characters; Table~\ref{tab:review-length}).
Retrieval adds far more input, writes reviews of baseline length, and still
gains most on the total, so length alone does not explain the pattern; a
length contribution to the graph-arm estimate cannot be ruled out.

\paragraph{Internal validity.}
Language models judge language-model output. The panel and the generator share
training data, so a judge may reward its own family's output. Mitigations:
the five-judge panel draws from five providers, per-judge verdicts are
published for re-aggregation, and judging runs at $T = 0.0$. Tested alone,
four of five judges reproduce the targeted effect at $p < 0.05$ and all five
preserve its direction.

The arms differ not only in the context block but also in the instruction
paragraph that names it (Appendix~\ref{app:prompts}), so the contrast
confounds content with instruction. Experiment~2's oracle mitigates this:
detection requires naming a genuinely broken site, which the instruction
alone cannot supply. The retrieval index covers only the first 200
candidate source files per repository (Section~\ref{sec:approach-rag}), so
in large repositories the retrieval arm is bounded by its index as well as
by its method.

Generation is not bit-reproducible: regenerating the graph arm from identical
inputs moves the nine-criterion score by $0.88$~points on average, unbiased
but of the same magnitude as the effect. The reviews and per-judge verdicts
are retained, so every number recomputes from stored outputs. Both
experiments generate at $T=0.0$, but they use different samples, outcomes,
and diff-length caps. Their effect magnitudes are therefore not compared
directly.

Experiment~2's injections are drawn from mutation-testing operators; selection
was direction-blind, but a seeded rename is not a defect a human introduced.
Its two-stage adjudication was introduced after an audit exposed false
negatives in the original judge prompt. The correction is post-hoc; the
frozen reviews and original verdicts remain available, and the thesis reports
the corrected measurement as such.

The human study uses an assisted interface with summaries, file tooltips, and
a diff shown by default. Its primary result therefore measures preference
under that interface, not unaided review reading. Four pre-registered
diagnostic checks of aid use, rationale content, file-count association, and
task time were not implemented and support no thesis claim.

\paragraph{External validity.}
Three substitute generators preserve the direction but none reaches
significance; both the generator and the held-out rows are from the earlier
text-heuristic configuration at $n = 40$ under a three-judge panel, so they
are directional evidence rather than tests of the current configuration. The
rubric-measured benefit is conditional on the generator.

Two graph constructions appear (the Joern code property graph in
Experiments~1 and~2, and import-statement edges in the human study), and
they are not interchangeable. The human study's edges are file-level; its
effect is an effect of that construction. Claims about graph
constructions are claims about these constructions.

All 35 pull requests were merged (survivorship). The selection acts
identically on all arms, which protects the within-sample comparison, but it
can still limit how the treatment effect generalises beyond the selected
sample. The human study covers six pull requests from three Python
libraries, so its preference result does not establish cross-language
generality.

\paragraph{Conclusion validity.}
With $n = 35$ the study is powered for $d_z \geq 0.47$ at 80\,\%. The graph arm
reaches significance on the full rubric ($p = 0.017$,
$d_z = +0.44$) but does not survive Bonferroni correction
($p_{\text{corr}} = 0.066$). The conclusion is therefore framed in
terms of the subscale, where the effect is larger
($d_z = +0.56$, $p = 0.003$,
$p_{\text{corr}} = 0.011$). The human study's 27
raters are sufficient for the primary endpoint (rank-biserial $r = 0.80$)
but explicitly underpowered for the overall-usefulness endpoint, which is
designated estimation-only. On 12 further pull requests not used in the
main analysis, evaluated under the earlier text-heuristic configuration, the
KG-relevant estimate remains positive ($+0.42$) but does not confirm the
effect ($p = 0.312$); the small replication was underpowered and is
interpreted as directional consistency rather than confirmation.



    %%%%%%%%%%%%%%%%%%%%%%%%%%%%%%%%%
    %% End of adding your content. %%
    %%%%%%%%%%%%%%%%%%%%%%%%%%%%%%%%%


    % Add the following chapters not to the current ›part‹ but one level above instead.
    \makeatletter
        \def\toclevel@chapter{-1}
        \def\toclevel@section{0}
    \makeatother

    \chapter{Conclusions \& Future Work}\label{chap:conclusion}
    \section{Summary of Findings}
\label{sec:conclusion-summary}

This thesis asked how repository context changes review coverage, whether
structural context enables cross-file reasoning, and whether human raters
perceive that capability.

\subsection*{RQ1: Comparative coverage}

There is no single best review mode. Retrieval produces the largest gain on
the full rubric ($+1.03$, $p = 0.002$), while the structural
arm produces the largest gain on the nine structurally relevant criteria
($+0.69$, $p = 0.003$).
After correction for multiple comparisons, each retains the gain on the scale
on which it leads; the hybrid retains neither. The graph-arm estimate is
sensitive to regeneration and to the judge panel, so its magnitude should be
read with the uncertainty reported in Section~\ref{sec:threats}.

\subsection*{RQ2: Cross-file reasoning}

Structural context enables a capability that the diff-only reviewer lacks. On
28 defects whose consequences lie outside the edited file, the baseline
detects none and the graph arm detects 18 ($p < 0.0001$); retrieval
detects 5. The hybrid detects 21 but is not significantly better than the graph
arm alone. Adding inheritance edges raises detection to 26 of 28. These
results establish the capability on controlled
synthetic defects; they do not estimate its prevalence in the broader
population of pull requests.

\subsection*{RQ3: Human preference}

Twenty-seven raters preferred the graph-augmented review when asked which review names the affected code
(0.710, $p = 0.0005$). They showed no preference on
overall usefulness (0.512). After correction they preferred the baseline
review on code clarity. The human-study result therefore corroborates the
targeted affected-code benefit, not a general improvement in review quality.

\subsection*{Synthesis}

In the controlled experiment, structural context gives the language model a
capability unavailable from the diff alone: identifying a consequence that is
not in the diff. On the
35-pull-request benchmark its average gain is small but sharply localised:
96\% ($p = 0.003$) lands in the criteria the mechanism
predicts.

\section{Limitations}
\label{sec:conclusion-limitations}

This work has three boundaries. First, it models static relationships within
one repository, not dynamic or cross-repository dependencies. Second, the
evaluation uses selected pull requests, synthetic defects, and a small
Python-only human study. Third, it measures review coverage, controlled
defect detection, and preference rather than production defects or developer
productivity. Section~\ref{sec:threats} discusses the methodological threats
in detail.

\section{Future Work}
\label{sec:conclusion-future}

\paragraph{Issue-guided cross-repository graphs.}
The graphs in this thesis stop at the repository boundary and are seeded from
the changed files. Future work should connect related repositories through
APIs, schemas, package manifests, events, and test dependencies. Issue and
pull-request descriptions can identify relevant entities and intended
behaviour, then guide a demand-driven traversal of this larger graph. Textual
descriptions should rank and restrict the retrieved subgraph rather than create
structural edges by themselves. A controlled evaluation could compare a
single-repository graph, a complete cross-repository graph, and an issue-guided
cross-repository subgraph under the same context budget. Relevant outcomes
include detection of downstream breakage, false-positive rate, context size,
latency, and token cost.

\paragraph{Review-history memory.}
The current system treats each pull request independently. A temporal graph
could connect code entities to earlier issues, pull requests, review comments,
and the outcomes of those comments, including whether a finding was accepted,
fixed, or rejected. A reviewer could then retrieve relevant precedents when a
similar component or contract changes again. Evaluation should use a strict
temporal split, allowing access only to history that predates each test pull
request. Relevant outcomes include precision on later reviews, false-positive
reduction, developer acceptance, and adaptation when project conventions
change. Such a system would also need mechanisms for forgetting stale
preferences and avoiding the reinforcement of incorrect historical decisions.

\paragraph{Adaptive context selection.}
The pipeline supplies the same context block to every pull request, although
the graph arm's benefit is concentrated in changes with cross-file
consequences (Section~\ref{sec:discussion-practical}). A lightweight gate
could instead decide per pull request whether graph context is warranted.
The simplest gate is deterministic: compute the change's structural fan-out
from the code property graph before generation and supply the graph block
only when the changed functions have dependents outside the edited files. A
learned alternative would follow adaptive retrieval, where a small classifier
or the generator itself predicts whether retrieval will help and skips it
otherwise~\cite{jeong2024adaptiverag,asai2024selfrag,jiang2023flare};
Repoformer applies this to repository-level code completion and reports that
retrieval helps on only a minority of instances~\cite{wu2024repoformer}.
Pull-request change-impact analysis already derives per-file risk from
call-graph dependencies and history~\cite{gocmen2025changeimpact}, so the
gate signal is available at review time. Evaluation should follow the
routing literature: compare the gated pipeline with always-graph and
never-graph baselines on rubric score and detection, report tokens and
latency, and measure the gate's agreement with an oracle label marking the
pull requests on which the graph arm outperformed the
baseline~\cite{ong2025routellm,chen2024frugalgpt}. The risk is that a gate's
false negatives reinstate the blind spot the graph was introduced to close,
so recall on cross-file changes should be weighted above precision.

\paragraph{Correctness-oriented and cross-language evaluation.} A rubric that
asks whether a review's claims are \emph{grounded} (whether the component it
names is genuinely affected) would measure what matters. Cross-repository
contract tests and runtime traces could verify predicted breakage, while
developer acceptance or rejection of review findings would measure practical
utility. A follow-up human study should include pull requests from non-Python
repositories and language-appropriate graph builders to test whether the
affected-code preference and clarity trade-off generalise across languages.
Together, these evaluations would expose cases in which syntactically valid
but coarse-grained edges support irrelevant claims.

\paragraph{Human calibration of the judge panel.}
This thesis validates its LLM judges only against each other and never
against human judgement. The human study and Experiment~1 also cannot be
compared directly, since they differ in stimuli, graph construction and
criteria. A direct calibration check would close this gap at low cost. The
five-judge panel would score the twelve human-study reviews in the same
pairwise format the raters used, with the H1--H6 wording unchanged. For each
pull request and criterion, the panel's majority verdict would then be
compared with the raters' majority preference, and their agreement reported.
Each human criterion would also be mapped to its closest rubric criterion: H1
to F3 and F4, H2 to F2, H3 to T3, and H5 to R1. This tests whether the
effects point in the same direction in both instruments. Such a check would
show whether the panel reproduces human preferences, which is the
precondition for reading the 35-pull-request rubric results as evidence of
review quality. With only six pull requests, the agreement estimates would be
imprecise. The check should therefore be repeated on a larger stimulus set
built with the Joern pipeline.


    % Following are the files and commands for the bibliography and the author’s publications.
    \pagestyle{plain}

    \renewcommand*{\bibfont}{\small}
    \printbibheading
    \addcontentsline{toc}{chapter}{Bibliography}
    \printbibliography[heading = none]

    \appendix
    \chapter{Supporting Analyses}\label{app:supporting-analyses}
    \section{Experiment 1 Supporting Statistics}
\label{app:supporting-multiplicity}

Table~\ref{tab:exp1-means} gives the absolute mode scores. The paired
differences in the main chapter remain the informative comparison.

\begin{table}[htbp]
    \centering
    \caption{Mean rubric score per mode, $n = 35$, with 95\,\% bootstrap
             confidence intervals ($B = 10\,000$).}
    \label{tab:exp1-means}
    \small
    \begin{tabular}{lcc}
        \toprule
        Mode & Total /25 [95\,\% CI] & KG-relevant /9 [95\,\% CI] \\
        \midrule
        \baseline{} & 8.03 [7.51, 8.54] & 4.66 [4.29, 5.03] \\
        \kg{}       & 8.74 [8.11, 9.37] & 5.34 [4.94, 5.74] \\
        \rag{}      & 9.06 [8.57, 9.54] & 5.09 [4.69, 5.49] \\
        \hybrid{}   & 8.60 [8.20, 9.00] & 5.11 [4.77, 5.46] \\
        \bottomrule
    \end{tabular}
\end{table}

Table~\ref{tab:exp1-deltas-full} gives the inferential details behind the
paired estimates in Table~\ref{tab:exp1-deltas}.

\begin{table}[htbp]
    \centering
    \caption{Paired permutation tests and effect sizes against \baseline{},
             $n = 35$ ($B = 20\,000$, two-sided, uncorrected).}
    \label{tab:exp1-deltas-full}
    \small
    \begin{tabular}{lcccccc}
        \toprule
        & \multicolumn{3}{c}{Total /25} & \multicolumn{3}{c}{KG-relevant /9} \\
        \cmidrule(lr){2-4} \cmidrule(lr){5-7}
        Mode & $\Delta$ [95\,\% CI] & $p$ & $d_z$ & $\Delta$ [95\,\% CI] & $p$ & $d_z$ \\
        \midrule
        \kg{}     & $+0.71$ [$+0.20$, $+1.26$] & 0.017 & $+0.44$ & $+0.69$ [$+0.29$, $+1.11$] & \textbf{0.003} & $+0.56$ \\
        \rag{}    & $+1.03$ [$+0.46$, $+1.60$] & 0.002 & $+0.59$ & $+0.43$ [$-0.03$, $+0.94$] & 0.127 & $+0.29$ \\
        \hybrid{} & $+0.57$ [$+0.09$, $+1.11$] & 0.056 & $+0.35$ & $+0.46$ [$+0.03$, $+0.91$] & 0.068 & $+0.34$ \\
        \bottomrule
    \end{tabular}
\end{table}

The governing Experiment~1 analysis is the pre-specified paired contrast
between each context arm and \baseline{}, tested by sign-flip permutation.
As a sensitivity to test choice and multiplicity, Table~\ref{tab:exp1-corrected}
also reports a Friedman omnibus and Bonferroni-corrected Wilcoxon signed-rank
tests. The Friedman test does not detect an omnibus mode difference on the
KG-relevant subscale. The corrected \kg{} comparison is therefore read only as
an alternative analysis of the planned contrast, not as post-hoc evidence
following a significant omnibus result.

\begin{table}[htbp]
    \centering
    \caption{Bonferroni-corrected Wilcoxon sensitivity. Friedman omnibus:
             $\chi^2 = 14.08$, $p = 0.003$ on total;
             $\chi^2 = 6.63$, $p = 0.084$ on KG-relevant.}
    \label{tab:exp1-corrected}
    \small
    \begin{tabular}{lcccc}
        \toprule
        & \multicolumn{2}{c}{Total /25} & \multicolumn{2}{c}{KG-relevant /9} \\
        \cmidrule(lr){2-3} \cmidrule(lr){4-5}
        Mode vs \baseline{} & $p$ raw & $p$ corrected & $p$ raw & $p$ corrected \\
        \midrule
        \kg{}     & 0.022 & 0.066 & 0.004 & \textbf{0.011} \\
        \rag{}    & 0.001 & \textbf{0.004} & 0.118 & 0.354 \\
        \hybrid{} & 0.048 & 0.143 & 0.064 & 0.192 \\
        \bottomrule
    \end{tabular}
\end{table}

\section{Experiment 1 Inter-Judge Agreement}
\label{app:supporting-agreement}

Table~\ref{tab:exp1-agreement} reports agreement before majority aggregation.
Pairs are evaluated on their shared valid rubric cells; Gemini abstained on
104 of 3\,500 cells, while the other judges returned a verdict for every cell.

\begin{table}[htbp]
    \centering
    \caption{Pairwise agreement in the final five-judge Experiment~1 panel.
             G4 = gpt-4o, GM = gemini-2.5-flash, CL = claude-sonnet-4.5,
             DS = deepseek-v4-pro, and GR = grok-4.6.}
    \label{tab:exp1-agreement}
    \small
    \begin{tabular}{lrrr}
        \toprule
        Pair & Valid cells & Raw agreement & Cohen's $\kappa$ \\
        \midrule
        G4--GM & 3\,396 & 85.7\,\% & 0.701 \\
        G4--CL & 3\,500 & 89.8\,\% & 0.788 \\
        G4--DS & 3\,500 & 90.3\,\% & 0.786 \\
        G4--GR & 3\,500 & 92.0\,\% & 0.822 \\
        GM--CL & 3\,396 & 85.8\,\% & 0.708 \\
        GM--DS & 3\,396 & 88.4\,\% & 0.753 \\
        GM--GR & 3\,396 & 87.3\,\% & 0.729 \\
        CL--DS & 3\,500 & 87.1\,\% & 0.727 \\
        CL--GR & 3\,500 & 87.5\,\% & 0.735 \\
        DS--GR & 3\,500 & 95.3\,\% & 0.893 \\
        \bottomrule
    \end{tabular}
\end{table}
% Source: results/JOERN_UNIFIED_FIVE_JUDGE.json, agreement section.

\section{Experiment 1 Robustness Statistics}
\label{app:supporting-robustness}

Table~\ref{tab:exp1-robustness} gives the statistics behind
Figure~\ref{fig:exp1-robustness}. The primary and single-judge rows use the
Joern CPG at $n = 35$. The generator and held-out rows are from an earlier
configuration that used a text-heuristic file matcher at $n = 40$ under a
three-judge panel; they are retained as directional evidence that alternative
generators preserve the sign of the effect but are not directly comparable
to the primary row. Values are therefore interpreted within rows.

\begin{table}[htbp]
    \centering
    \caption{The \kg{} effect under alternative generators, judges,
             and samples.}
    \label{tab:exp1-robustness}
    \small
    \setlength{\tabcolsep}{5pt}
    \begin{tabular}{llcccc}
        \toprule
        & & \multicolumn{2}{c}{Total /25} & \multicolumn{2}{c}{KG-relevant /9} \\
        \cmidrule(lr){3-4} \cmidrule(lr){5-6}
        Varied & Configuration & $\Delta$ & $p$ & $\Delta$ & $p$ \\
        \midrule
        ---       & primary configuration, $n = 35$ & $+0.71$ & \textbf{0.017} & $+0.69$ & \textbf{0.003} \\
        \midrule
        Judge     & gpt-4o alone             & --- & --- & $+0.63$ & \textbf{0.007} \\
                  & gemini-2.5-flash alone   & --- & --- & $+0.74$ & \textbf{0.005} \\
                  & claude-sonnet-4.5 alone  & --- & --- & $+0.26$ & 0.192 \\
                  & deepseek-v4-pro alone    & --- & --- & $+0.60$ & \textbf{0.019} \\
                  & grok-4.6 alone           & --- & --- & $+0.60$ & \textbf{0.009} \\
        \midrule
        Generator$^{*}$ & claude-haiku-4.5    & $+0.30$ & 0.407 & $+0.15$ & 0.457 \\
                  & gemini-2.5-flash         & $+0.33$ & 0.465 & $+0.47$ & 0.136 \\
                  & deepseek-chat            & $+0.47$ & 0.117 & $+0.20$ & 0.383 \\
        \midrule
        Sample$^{*}$ & held-out PRs, $n = 12$   & $+0.42$ & 0.426 & $+0.42$ & 0.312 \\
        \bottomrule
        \multicolumn{6}{l}{\footnotesize $^{*}$\,Text-heuristic builder, $n = 40$, three-judge panel; retained as directional evidence.} \\
    \end{tabular}
\end{table}

\section{Experiment 2 Detection Rates and Secondary Contrasts}
\label{app:supporting-exp2-contrasts}

Table~\ref{tab:exp2-detection} gives the exact detection rates behind
Figure~\ref{fig:exp2-detection}.

\begin{table}[htbp]
    \centering
    \caption{Detection rate on structural bands, $n = 28$, with Wilson
             95\,\% confidence intervals.}
    \label{tab:exp2-detection}
    \small
    \begin{tabular}{lcc}
        \toprule
        Arm & Detected & Rate [95\,\% CI] \\
        \midrule
        \baseline{}                 & 0/28  & 0.00 [0.00, 0.12] \\
        \rag{}                      & 5/28  & 0.18 [0.08, 0.36] \\
        \kg{} (deployed)            & 18/28 & 0.64 [0.46, 0.79] \\
        \hybrid{}                   & 21/28 & 0.75 [0.57, 0.87] \\
        \midrule
        \kg{} + inheritance edges   & 26/28 & 0.93 [0.77, 0.98] \\
        \kg{} idealised (ceiling)   & 26/28 & 0.93 [0.77, 0.98] \\
        \bottomrule
    \end{tabular}
\end{table}

Table~\ref{tab:exp2-perband} breaks detection down by structural band and
local control. The five structural bands (S1--S5) represent different
dependency types; the two local bands (L1--L2) serve as negative controls
whose consequences are visible in the diff.

\begin{table}[htbp]
    \centering
    \caption{Per-band detection counts. $n$ is the number of injections in
             each band. Structural bands (S1--S5) require cross-file
             reasoning; local bands (L1--L2) do not.}
    \label{tab:exp2-perband}
    \small
    \begin{tabular}{lrcccccc}
        \toprule
        Band & $n$ & \baseline{} & \kg{} & \rag{} & \hybrid{} & Inherit & Idealised \\
        \midrule
        S1 & 9  & 0/9  & 4/9  & 3/9  & 6/9  & 9/9  & 8/9  \\
        S2 & 9  & 0/9  & 6/9  & 1/9  & 7/9  & 7/9  & 8/9  \\
        S3 & 1  & 0/1  & 1/1  & 0/1  & 1/1  & 1/1  & 1/1  \\
        S4 & 6  & 0/6  & 4/6  & 1/6  & 4/6  & 6/6  & 6/6  \\
        S5 & 3  & 0/3  & 3/3  & 0/3  & 3/3  & 3/3  & 3/3  \\
        \midrule
        L1 & 6  & 6/6  & 6/6  & 6/6  & 6/6  & 6/6  & 6/6  \\
        L2 & 6  & 6/6  & 6/6  & 6/6  & 5/6  & 6/6  & 6/6  \\
        \bottomrule
    \end{tabular}
\end{table}

The structural bands confirm the aggregate pattern: baseline scores zero in
every band, and the eight cases recovered by inheritance edges fall in S1
($+5$), S2 ($+1$) and S4 ($+2$). The inheritance and idealised arms match
in aggregate, not band by band. Local controls are at or near ceiling for
all arms (one hybrid miss on L2).

The primary Experiment~2 contrast is \kg{} versus \baseline{}. Table
\ref{tab:exp2-secondary-contrasts} retains the secondary comparisons used to
interpret retrieval, hybrid concatenation, and inheritance edges.

\begin{table}[htbp]
    \centering
    \caption{Secondary exact McNemar contrasts on the 28 structural
             injections.}
    \label{tab:exp2-secondary-contrasts}
    \small
    \begin{tabular}{lccc}
        \toprule
        Contrast & $\Delta$ detection & Discordant wins/losses & $p$ \\
        \midrule
        \rag{} vs \baseline{} & $+0.179$ & 5/0 & 0.0625 \\
        \hybrid{} vs \kg{} & $+0.107$ & 4/1 & 0.3750 \\
        \kg{} + inheritance vs \kg{} & $+0.286$ & 8/0 & 0.0078 \\
        \bottomrule
    \end{tabular}
\end{table}

\subsection{Residual Idealised-Arm Misses}
\label{app:supporting-exp2-misses}

The two idealised-arm misses arise at different stages. In
\texttt{kafka\_S2\_01}, the review names genuine dependent classes, including
\texttt{CachingKeyValueStore} and \texttt{CachingSessionStore}, but does not
emit the exact file token required by the deterministic filename gate. This is
a boundary of the operational definition rather than an absence of affected
code in the review. In \texttt{kafka\_S1\_02}, whose resolver returns 82 true
dependents, the review warns generically about broken dependencies but names
none. That case is a generation or evidence-salience failure despite complete
structural evidence.
% Sources: experiments/2026-07-05_injection_exp2/out/manifest.json;
% out/reviews/{kafka_S2_01,kafka_S1_02}/kg_idealised.md;
% experiments/2026-09-23_five_judge_panel/RESULTS.json.

\section{Component Contributions}
\label{sec:results-ablations}
\label{app:supporting-components}

These component analyses support RQ2 by testing which parts of the graph block
contribute to cross-file detection. The rubric's regeneration variance
(Section~\ref{sec:repro}) makes it too coarse to separate sub-components, so
the analyses use the detection oracle instead.

The deployed graph arm supplies two kinds of dependency evidence
simultaneously: a file-level dependency list and resolved call-graph edges from
the code property graph. Table~\ref{tab:exp2-features} isolates each.

\begin{table}[htbp]
    \centering
    \caption{Detection on structural bands with one kind of dependency
        evidence at a time. Reviews use gpt-4o at $T=0.0$; rates are majority
        votes from five judges at $T=0.0$. The $p$~values are exact McNemar,
        Holm-corrected.}
    \label{tab:exp2-features}
    \begin{tabular}{llcc}
        \toprule
        Arm & Evidence supplied & Detected & Rate \\
        \midrule
        \baseline{}        & none                        & 0/28  & 0\,\% \\
        \texttt{deps only} & file-level dependency list   & 16/28 & 57\,\% \\
        \texttt{edges only}& resolved call edges         & 17/28 & 61\,\% \\
        \kg{}              & both                        & 18/28 & 64\,\% \\
        \midrule
        \multicolumn{2}{l}{Contrast} & $\Delta$ & $p$ (Holm) \\
        \midrule
        \multicolumn{2}{l}{\texttt{deps only} $-$ \baseline{}}  & $+57$pp & $< 0.001$ \\
        \multicolumn{2}{l}{\texttt{edges only} $-$ \baseline{}} & $+61$pp & $< 0.001$ \\
        \bottomrule
    \end{tabular}
\end{table}

Either component alone lifts detection from zero to over half, and both
survive correction. The dependency list recovers sixteen of the eighteen
combined-arm detections, while resolved call edges recover seventeen. Neither
component is significantly below the combined arm at $n=28$. The result
therefore supports partial substitutability, not the necessity or superiority
of resolved edges: the causal claim is that useful structural pointers enable
cross-file detection, while this sample cannot identify one representation as
uniquely responsible.

A pre-registered, Grafana/TypeScript-only test band, evaluated by the earlier
three-judge panel, addressed the test component. On 28 injections where a
renamed symbol provably breaks test files, the filename-convention test
matcher of Section~\ref{sec:approach-kg-render} named a broken test in 24 of 28 cases (86\,\%), compared with 27/28 (96\,\%)
when given the ground truth; the difference is not significant ($p = 0.38$).
Without the test
section, neither \baseline{} nor \kg{} named a broken test in any case. This
shows that the test block contributes detections the dependency block alone
does not.

\subsection{Context Relevance}
\label{sec:results-context-relevance}

The component ablation shows that dependency edges help, but does not rule
out a prompt-length confound: the reviewer might try harder whenever it sees
\emph{any} dependency block, regardless of content. A post-hoc control tests
this possibility.
Table~\ref{tab:exp2-signal-noise} tests this by holding the number of
context items constant and varying only their relevance.
% Source: experiments/2026-09-19_joern_replace_grep/context_ablation/RESULTS_SIGNAL_NOISE_V3.md
% Design follows Yang et al. (EMNLP 2025) GSM-DC matched-distractor methodology.

\begin{table}[htbp]
    \centering
    \caption{Detection under matched-quantity contexts that differ only in
        relevance. \texttt{noise} draws the same number of edges from
        unrelated symbols in the same repository; \texttt{relevant} filters
        the deployed edges to those pointing at true dependents.
        All arms use gpt-4o at $T=0.0$ and the earlier three-judge panel.}
    \label{tab:exp2-signal-noise}
    \begin{tabular}{llcc}
        \toprule
        Arm & Context relevance & Detected & Rate \\
        \midrule
        \baseline{}             & none             & 0/28  & 0\,\% \\
        \texttt{noise}          & wrong edges, matched count & 1/28  & 4\,\% \\
        \kg{}                   & deployed mix     & 18/28 & 64\,\% \\
        \texttt{relevant only}  & true-dependent edges only & 25/28 & 89\,\% \\
        \texttt{idealised}      & ground-truth resolver & 26/28 & 93\,\% \\
        \bottomrule
    \end{tabular}
\end{table}

The noise arm differs little from \baseline{} ($1/28$ versus $0/28$):
matched-volume edges drawn from unrelated symbols provide no comparable
detection gain. Filtering the deployed edges to only those pointing at true
dependents raises detection from 18 to 25 of 28. This post-hoc staircase is
consistent with the context-rot finding that
irrelevant tokens degrade output quality even when the answer is present in
the input~\cite{hong2025contextrot, gupta2025contextlength}.

\section{Human-Study Exploratory Endpoints}
\label{app:supporting-human}

The five non-primary criteria are exploratory and Holm-corrected. Only H5
(code clarity) remains distinguishable after correction, favouring the
\baseline{} review.

\begin{table}[htbp]
    \centering
    \caption{Exploratory human-study criteria. Values above $0.5$ favour
             \kg{}. IDs are prefixed H to distinguish them from the
             25-criterion rubric.}
    \label{tab:human-exploratory}
    \small
    \begin{tabular}{llcccc}
        \toprule
        ID & Criterion & Mean [95\,\% CI] & Non-tie $n$ & $p$ & $p_{\text{Holm}}$ \\
        \midrule
        H2 & Describes failure cases     & 0.469 [0.414, 0.528] & 25 & 0.146 & 0.499 \\
        H3 & Points to tests             & 0.454 [0.395, 0.512] & 19 & 0.125 & 0.499 \\
        H4 & Explains why                & 0.491 [0.438, 0.543] & 20 & 0.735 & 1.000 \\
        H5 & Comments on code clarity    & 0.414 [0.370, 0.460] & 24 & 0.006 & 0.028 \\
        H6 & Covers important issues     & 0.481 [0.417, 0.548] & 22 & 0.636 & 1.000 \\
        \bottomrule
    \end{tabular}
\end{table}

\section{Prompt and Review Length}
\label{app:supporting-length}

Table~\ref{tab:prompt-length} reports the total input size (system prompt
plus user prompt) and the augmentation overhead for each arm, computed
from all 40 stored evidence packs (including the five Go pull requests
excluded from the $n = 35$ analysis sample) and the generation template.
Table~\ref{tab:review-length} reports the output review length.

\begin{table}[htbp]
    \centering
    \caption{Total input size per arm ($n = 40$ stored evidence packs).
             Token estimates use the chars\,/\,4 heuristic.}
    \label{tab:prompt-length}
    \small
    \begin{tabular}{lrrrr}
        \toprule
        Arm & Mean chars & Median & $\approx$\,Mean tokens & Overhead vs \baseline{} \\
        \midrule
        \baseline{} & 16\,926 & 12\,338 & 4\,232 & --- \\
        \kg{}       & 18\,384 & 14\,240 & 4\,596 & $+$10.9\,\% \\
        \rag{}      & 21\,292 & 15\,733 & 5\,323 & $+$43.4\,\% \\
        \hybrid{}   & 22\,538 & 17\,016 & 5\,635 & $+$52.3\,\% \\
        \bottomrule
    \end{tabular}
\end{table}

\begin{table}[htbp]
    \centering
    \caption{Review output length per arm ($n = 40$ stored evidence packs).}
    \label{tab:review-length}
    \small
    \begin{tabular}{lrrr}
        \toprule
        Arm & Mean chars & Mean words & $\approx$\,Mean tokens \\
        \midrule
        \baseline{} & 1\,774 & 243 & 444 \\
        \kg{}       & 1\,991 & 266 & 498 \\
        \rag{}      & 1\,782 & 243 & 445 \\
        \hybrid{}   & 1\,857 & 252 & 464 \\
        \bottomrule
    \end{tabular}
\end{table}

The graph arm adds the least input overhead ($+$10.9\,\%) but produces
the longest reviews ($+$12\,\% in characters and $+$9\,\% in words relative
to baseline). RAG
adds more input ($+$43\,\%) yet produces reviews of the same length as
baseline. This asymmetry suggests the model converts structural context
into additional commentary more readily than it converts similarity
context.



    \chapter{Review Example}\label{app:review-example-chapter}
    \label{app:review-example}

This appendix reproduces excerpts from the baseline and graph-augmented
reviews of \texttt{click} pull request \#3493, which narrows the
\texttt{echo} function's \texttt{message} parameter from
\texttt{typing.Any} to \texttt{object} and replaces the type-dispatch
chain with a \texttt{match} statement.  The two reviews were generated by
GPT-4o under matched conditions; the prompts differ only in their
references to the available context, and the graph-augmented arm
additionally receives the structural block.
Sections are condensed for space; the full reviews are archived in the
study repository.

\bigskip

\noindent
\begin{minipage}[t]{0.48\textwidth}
\noindent\textbf{Baseline review} (diff only)\par
\smallskip\hrule\smallskip
{\small
\textbf{Problems found}
\begin{enumerate}[leftmargin=*,itemsep=2pt]
\item The change from \texttt{t.Any} to \texttt{object} may break callers
      that rely on more permissive typing.
\item The \texttt{match} statement may not be compatible with older Python
      versions.
\item Handling of \texttt{None} as a message is not explicitly tested.
\item The interaction of \texttt{nl} with \texttt{None} and empty strings
      is not thoroughly tested.
\end{enumerate}

\textbf{Files referenced}
\begin{itemize}[leftmargin=*,itemsep=1pt]
\item \texttt{src/click/utils.py:236}
\item \texttt{src/click/utils.py:287--295}
\item \texttt{tests/test\_utils.py:98--102}
\end{itemize}
}
\end{minipage}%
\hfill
\begin{minipage}[t]{0.48\textwidth}
\noindent\textbf{Graph-augmented review} (diff + KG context)\par
\smallskip\hrule\smallskip
{\small
\textbf{Problems found}
\begin{enumerate}[leftmargin=*,itemsep=2pt]
\item The type change from \texttt{t.Any} to \texttt{object} might break
      callers that pass types not explicitly handled.
\item The change could violate existing API contracts if callers expect
      \texttt{echo} to accept additional types.
\item \texttt{echo} is widely used across the codebase, including in
      \texttt{src/click/core.py} and \texttt{src/click/termui.py}, which
      could lead to unexpected behaviour.\label{itm:kg-callers}
\item No test covers \texttt{echo} with complex objects that are not
      \texttt{str}, \texttt{bytes}, or \texttt{bytearray}.
\end{enumerate}

\textbf{Files referenced}
\begin{itemize}[leftmargin=*,itemsep=1pt]
\item \texttt{src/click/utils.py:232--311}
\item \texttt{tests/test\_utils.py:98--104}
\item \texttt{src/click/core.py}\;\textsuperscript{$\star$}
\item \texttt{src/click/termui.py}\;\textsuperscript{$\star$}
\item \texttt{src/click/testing.py}\;\textsuperscript{$\star$}
\end{itemize}

\smallskip
{\footnotesize $\star$\;File absent from the diff; surfaced by the
knowledge graph as dependent on the changed module.}
}
\end{minipage}

\bigskip

The baseline review identifies risks that are visible in the diff---the
type narrowing and the new \texttt{match} syntax---but cannot name
\emph{where} those risks materialise beyond the changed file itself.  The
graph-augmented review names three additional files
(\texttt{core.py}, \texttt{termui.py}, \texttt{testing.py}) that depend on
the changed module, two of which call \texttt{echo} at runtime.  The
study's answer key confirms that
\texttt{core.py} and \texttt{termui.py} contain direct call sites of the
changed function; \texttt{testing.py} imports utilities from the same
module.  This is the pattern the primary endpoint (H1) was designed to
detect: the graph supplies evidence that lets the reviewer ground a claim
about affected code in a concrete file reference rather than a generic
warning.


    \chapter{Prompts}\label{chap:appendix-prompts}
    \label{app:prompts}

This appendix reproduces the prompts, in discharge of guideline~3 of
\textcite{baltes2026llmguidelines}. They are
transcribed from \texttt{prnote/note.py} for generation and from
\texttt{scripts/evaluate\_reviews.py} for scoring, the two files that produced
the reviews and verdicts of Experiments~1 and~2 analysed in
Chapter~\ref{chap:results}. The transcription is character-exact except that
one em dash in the required output heading is set here as a hyphen and the
code fence that wraps the output format in the source is omitted. The human
study used stricter variants of the generator prompts
(Section~\ref{sec:human-study}); these are archived with the study materials
in the replication package rather than reproduced here.

\section{Generator Prompts}
\label{app:prompts-generator}

Every arm shares the following system prompt.

{\footnotesize
\begin{verbatim}
You are a senior software engineer conducting a thorough code review of a
pull request.

Your task is to generate a structured review note that helps the PR author
understand:
1. What potential issues exist
2. Evidence from the code supporting your concerns
3. The impact of these issues
4. Concrete recommendations

Output Format (follow exactly):

# Review Note - Evidence-Anchored

**Scope:** [One sentence describing what this PR does]

## Problem
[Numbered list of 2-3 specific concerns or issues you identified]

## Evidence
[Bullet points with specific file:line references supporting your concerns]

## Impact
[Explain the technical impact - what could go wrong, what risks exist]

## Recommendation (Fix / Tests / Risks)
[Numbered list of concrete, actionable recommendations]

## Traceability
[List relevant code owners or teams if known, otherwise state
"Not specified"]

Guidelines:
- Be specific and cite actual file paths and line numbers from the diff
- Focus on meaningful issues, not style nitpicks
- Consider integration impacts and test coverage
- Be constructive and professional
\end{verbatim}
}

The imposed output structure bears on the measurement. Because all four arms are
told to produce the same five headings, the rubric cannot reward an arm for
organising its output; it can only reward what fills those headings. This narrows
the space in which a difference between arms can appear, and it is one reason the
effects reported in Chapter~\ref{chap:results} are small in absolute terms.

Each arm then appends one paragraph naming the context it has been given. This is
the instruction asymmetry disclosed in Section~\ref{sec:setup}.

\paragraph{\baseline{}.}
{\footnotesize
\begin{verbatim}
Context: You only have access to the PR diff. Make reasonable inferences
about integration and testing needs based on the changes you can see.
\end{verbatim}
}

\paragraph{\kg{}.}
{\footnotesize
\begin{verbatim}
Context: You have access to Knowledge Graph context including:
- Which files are changed and their dependencies
- Which tests cover the changed files
- Who owns the affected code areas
- Which other files call or import the changed code

Use this structural context to provide deeper insights about integration
risks, missing test coverage, and code ownership concerns.
\end{verbatim}
}

\paragraph{\rag{}.}
{\footnotesize
\begin{verbatim}
Context: You have access to semantically similar code chunks from the
repository. These show patterns used elsewhere that may be relevant to
this change.

Use this context to identify consistency issues, suggest similar patterns
to follow, and flag potential conflicts with existing code.
\end{verbatim}
}

\paragraph{\hybrid{}.}
{\footnotesize
\begin{verbatim}
Context: You have access to both:
1. Knowledge Graph context (tests, dependencies, ownership)
2. Semantically similar code from the repository

Combine structural insights (what depends on what, test coverage) with
semantic insights (similar patterns elsewhere) to provide comprehensive
review feedback.
\end{verbatim}
}

The user message is assembled from the frozen evidence pack in a fixed order: the
pull request title; its description, capped at 8\,000 characters and omitted
entirely when empty; the unified diff, capped as given in
Table~\ref{tab:config}; the context block if the arm has one; and the closing
instruction. The diff therefore precedes the context, so in the two arms that
carry a block the added material is the last thing the model reads before the
instruction.
% dataset_v2/scripts/regenerate_reviews_v2.py:110-117 (Experiment 1) and
% prnote/note.py:661-669 (Experiment 2) both order title, body, diff, context.

\section{Per-Criterion Satisfaction Rates}

Table~\ref{tab:exp1-criteria} gives the satisfaction rate behind the
criterion-level analysis. It is here rather than in
Chapter~\ref{chap:results} because no
individual criterion reaches significance once the 17 that move are corrected
for multiplicity, so the chapter's claim is about the band and this table is
reference material for checking it.

\begin{table}[htbp]
    \centering
    \caption{Satisfaction rate per rubric criterion, \baseline{} versus \kg{},
             $n = 35$ (majority of five judges), for all 25 criteria. $\Delta$ is
             in percentage points. The seven criteria marked $\dagger$ sit at $0\,\%$ or
             $100\,\%$ in both arms and cannot register a difference in either
             direction.}
    \label{tab:exp1-criteria}
    \small
    \begin{tabular}{llccc}
        \toprule
        ID & Criterion (abbreviated) & \baseline{} & \kg{} & $\Delta$ \\
        \midrule
        \multicolumn{5}{l}{\emph{KG-relevant band (9 criteria)}} \\
        F3 & Integration with existing code       & 60.0\,\%  & 82.9\,\%  & $+22.9$ \\
        F4 & Breaking changes or impact           & 60.0\,\%  & 74.3\,\%  & $+14.3$ \\
        T1 & Need for tests discussed             & 97.1\,\%  & 100.0\,\% & $+2.9$ \\
        T2 & Edge cases and error paths in tests  & 65.7\,\%  & 68.6\,\%  & $+2.9$ \\
        T3 & Specific test files referenced       & 62.9\,\%  & 71.4\,\%  & $+8.6$ \\
        M1 & Fit with existing architecture       & 8.6\,\%   & 14.3\,\%  & $+5.7$ \\
        M3 & Public API or interface handling     & 8.6\,\%   & 17.1\,\%  & $+8.6$ \\
        C2 & Consistency with similar solutions   & 2.9\,\%   & 5.7\,\%   & $+2.9$ \\
        Q2$^{\dagger}$ & Comments anchored to code locations & 100.0\,\% & 100.0\,\% & $0.0$ \\
        \midrule
        \multicolumn{5}{l}{\emph{Negative-control band (16 criteria)}} \\
        F1 & Change solves the stated problem     & 5.7\,\%   & 8.6\,\%   & $+2.9$ \\
        F2 & Edge cases in the change identified  & 65.7\,\%  & 68.6\,\%  & $+2.9$ \\
        R1 & Code clarity, naming, function size  & 8.6\,\%   & 0.0\,\%   & $-8.6$ \\
        R2 & Unnecessary complexity flagged       & 2.9\,\%   & 2.9\,\%   & $0.0$ \\
        R3 & Comments explain why, not what       & 5.7\,\%   & 2.9\,\%   & $-2.9$ \\
        M2$^{\dagger}$ & Pull request too large to review & 0.0\,\% & 0.0\,\% & $0.0$ \\
        C1 & Project style-guide compliance       & 2.9\,\%   & 0.0\,\%   & $-2.9$ \\
        P1 & Obvious inefficiencies               & 8.6\,\%   & 11.4\,\%  & $+2.9$ \\
        P2$^{\dagger}$ & Benchmarks for critical paths  & 0.0\,\%  & 0.0\,\%  & $0.0$ \\
        S1 & Input validation                     & 14.3\,\%  & 8.6\,\%   & $-5.7$ \\
        S2$^{\dagger}$ & Hardcoded secrets or credentials & 0.0\,\% & 0.0\,\% & $0.0$ \\
        S3 & Error-handling quality               & 20.0\,\%  & 22.9\,\%  & $+2.9$ \\
        Q1$^{\dagger}$ & Overall assessment summary       & 100.0\,\% & 100.0\,\% & $0.0$ \\
        Q3 & Blocking versus non-blocking issues  & 2.9\,\%   & 14.3\,\%  & $+11.4$ \\
        Q4$^{\dagger}$ & Clarifying questions asked       & 0.0\,\% & 0.0\,\% & $0.0$ \\
        Q5$^{\dagger}$ & Rationale given for suggestions  & 100.0\,\% & 100.0\,\% & $0.0$ \\
        \bottomrule
    \end{tabular}
\end{table}

\section{Judge Prompt}
\label{app:prompts-judge}

The panel scores one review against all 25 criteria in a single call. The system
prompt is:

{\footnotesize
\begin{verbatim}
You are an expert code review evaluator. Your task is to assess whether a
PR review meets specific quality criteria.

For each criterion, you must:
1. Determine if the review addresses that criterion (score: 1) or not
   (score: 0)
2. Provide a brief quote or explanation as evidence

Be strict but fair:
- Score 1 only if the review CLEARLY addresses the criterion
- Score 0 if the criterion is not addressed or only vaguely mentioned
- Base your judgment on meaning and intent, not just keywords
\end{verbatim}
}

The user message interpolates the pull request under review, the review text, and
the numbered criteria of Section~\ref{sec:rubric}:

{\footnotesize
\begin{verbatim}
## PR Being Reviewed
Title: {title}
Description: {body, first 400 characters}

## Review to Evaluate
{review, first 4000 characters}

## Criteria to Check
1. [F1] {description}
...
25. [Q5] {description}

## Instructions
Evaluate the review against each criterion. Respond with a JSON object in
this exact format:

{
  "scores": [
    {"id": "F1", "score": 0,
     "evidence": "Review does not verify the change addresses the problem"},
    {"id": "F2", "score": 1,
     "evidence": "Review mentions: 'what about edge cases where X is null'"}
  ]
}

Important:
- Include ALL 25 criteria (F1-F4, T1-T3, R1-R3, M1-M3, C1-C2, P1-P2,
  S1-S3, Q1-Q5)
- Each score must be 0 or 1
- Evidence should be a short quote or explanation
\end{verbatim}
}

Three properties of this prompt bear on how its output should be read.

The two character caps could in principle bias the comparison, since an arm whose
reviews ran longer would have more of its text withheld from the judge. Neither
cap binds in practice. The longest of the 140 canonical reviews is 2\,620
characters and the longest of the 12 human study stimuli is 2\,053, both well
inside the 4\,000-character limit, so no review was truncated before scoring. The
400-character description cap does bind, but it applies identically to all four
arms of a given pull request.

The instruction requires evidence for every verdict, which makes the verdicts
auditable rather than opaque. The stored evidence strings are what
Section~\ref{sec:results-exp1-attribution} inspects when characterising what a
criterion was awarded for.

The judge is asked whether a criterion is \emph{addressed}, not whether the
review is correct about it. This is the coverage-not-correctness limitation that
governs the interpretation of every rubric number in this thesis, stated where
the rubric is defined (Section~\ref{sec:rubric}) and revisited in
Section~\ref{sec:threats}. Reading the prompt shows why the limitation is
structural rather than incidental: a judge given only the review and 400
characters of description has no access to the repository against which
correctness would have to be established.


    \addchap{Statement of Independent Work}
    \noindent
\textbf{Name:} Abdelrahman Khattab \hfill \textbf{Student ID:} 825360

\bigskip
\noindent
I hereby declare that I have prepared the present work without the assistance
of third parties and without the use of any sources or aids other than those
indicated. All passages of the work that I have taken from these sources and
aids, either verbatim or in terms of content, have been identified as such and
are listed in the bibliography. I further declare that neither I nor anyone
else has previously submitted this work, in its present or any modified form,
as academic credit.

\bigskip
\textbf{Information on the Use of (Generative) AI}

\medskip
\noindent
I declare in lieu of an oath that I have independently authored the present
work and have only used the sources and tools indicated. I declare that I have
verified that the content generated by AI systems and incorporated by me is
factually accurate. I am aware that, as the author of this work, I bear
responsibility for the statements and claims made therein.

\medskip
\noindent
AI tools were used for the following purposes:
\begin{itemize}[nosep]
    \item[\checkmark] Coding: support in programming experiment and evaluation scripts
    \item[\checkmark] Data Collection: generation of review comments and judge
          scores as part of the experimental methodology
          (described in Chapters~\ref{chap:approach} and~\ref{chap:edas})
\end{itemize}

\noindent
The tools used and their roles are documented in the AI Usage table on the
following page.

\vspace{6ex}
\begin{flushleft}
    Potsdam, 8 October 2026
    \hspace*{2em}
    \raisebox{-0.9\baselineskip}{%
        \begin{tabular}{p{5cm}}
            \hline
            \centering\footnotesize\printAuthor
        \end{tabular}%
    }
\end{flushleft}

\clearpage

\section*{Documentation Table for AI Usage}

\noindent
\small
\begin{longtable}{@{}c >{\raggedright}p{3.4cm} p{2.0cm} >{\raggedright}p{4.2cm} >{\raggedright\arraybackslash}p{3.8cm}@{}}
\toprule
\textbf{No.} & \textbf{AI Tool (Date)} & \textbf{Type of Use} &
\textbf{Key Inputs / Functions} &
\textbf{Outcome} \\
\midrule
\endhead
1 & OpenAI GPT-4o (2026) &
    Data collection &
    Review-comment generation for 35~PRs $\times$ 4~conditions
    (Ch.~\ref{chap:results}) &
    140 reviews used as experimental data \\[6pt]
2 & GPT-4o, Gemini-2.5-Flash, Claude-Sonnet-4.5, DeepSeek-v4-pro, Grok-4.6 (2026) &
    Data collection &
    Rubric scoring and defect-detection adjudication; 5-judge panel
    (Ch.~\ref{chap:results}) &
    Judge scores and detection verdicts \\[6pt]
3 & OpenAI text-embedding-3-small (2026) &
    Data collection &
    Embedding repository files for RAG retrieval
    (Ch.~\ref{chap:approach}) &
    Retrieval index for RAG condition \\[6pt]
4 & Claude-Haiku-4.5, Gemini-2.5-Flash, DeepSeek-Chat (2026) &
    Data collection &
    Substitute generators for sensitivity analysis
    (Ch.~\ref{chap:results}) &
    Robustness check \\[6pt]
\bottomrule
\end{longtable}



\end{document}